% Define Class: %%%%%%%%%%%%%%%%%%%%%%%%%%%%%%%%%%%%%%%%%%%%%%%%%%%%%%%%
\documentclass[template=tech-report]{nasa-latex-docs}

% Preamble Section: %%%%%%%%%%%%%%%%%%%%%%%%%%%%%%%%%%%%%%%%%%%%%%%%%%%%

% Acronyms
% List of acronyms (printed by the class). One \newacronym per entry; keep alphabetical.
\newacronym{ahi}{AHI}{Astronaut Health Intelligence}
\newacronym{ai}{AI}{Artificial Intelligence}
\newacronym{api}{API}{Application Programming Interface}
\newacronym{aqm}{AQM}{Air Quality Monitor}
\newacronym{ar}{AR}{Anomaly Recognition}
\newacronym{ared}{ARED}{Advanced Resistive Exercise Device}
\newacronym{ble}{BLE}{Bluetooth Low Energy}
\newacronym{bp}{BP}{Blood Pressure}
\newacronym{ccmc}{CCMC}{Community Coordinated Modeling Center}
\newacronym{cdf}{CDF}{Cumulative Distribution Function}
\newacronym{cevis}{CEVIS}{Cycle Ergometer with Vibration Isolation and Stabilization}
\newacronym{cme}{CME}{Coronal Mass Ejection}
\newacronym{csa}{CSA}{Canadian Space Agency}
\newacronym{csv}{CSV}{Comma-Separated Values}
\newacronym{dba}{dBA}{A-weighted decibel}
\newacronym{donki}{DONKI}{Space Weather Database Of Notifications, Knowledge, Information}
\newacronym{ecg}{ECG}{Electrocardiogram}
\newacronym{eda}{EDA}{Environmental Data Application}
\newacronym{eeg}{EEG}{Electroencephalogram}
\newacronym{esa}{ESA}{European Space Agency}
\newacronym{ewma}{EWMA}{Exponentially Weighted Moving Average}
\newacronym{fhir}{FHIR}{Fast Healthcare Interoperability Resources}
\newacronym{hera}{HERA}{Human Exploration Research Analog}
\newacronym{hr}{HR}{Heart Rate}
\newacronym{hrp}{HRP}{Human Research Program}
\newacronym{hrv}{HRV}{Heart Rate Variability}
\newacronym{http}{HTTP}{Hypertext Transfer Protocol}
\newacronym{iftn}{i18n}{Internationalization}
\newacronym{icd}{ICD}{International Classification of Diseases}
\newacronym{iot}{IoT}{Internet of Things}
\newacronym{iss}{ISS}{International Space Station}
\newacronym{jaxa}{JAXA}{Japan Aerospace Exploration Agency}
\newacronym{json}{JSON}{JavaScript Object Notation}
\newacronym{llm}{LLM}{Large Language Model}
\newacronym{lsda}{LSDA}{Life Sciences Data Archive}
\newacronym{lstm}{LSTM}{Long Short-Term Memory}
\newacronym{mad}{MAD}{Median Absolute Deviation}
\newacronym{mfa}{MFA}{Multi-Factor Authentication}
\newacronym{ml}{ML}{Machine Learning}
\newacronym{msv}{mSv}{Millisievert}
\newacronym{mvp}{MVP}{Minimum Viable Product}
\newacronym{nasa}{NASA}{National Aeronautics and Space Administration}
\newacronym{nlp}{NLP}{Natural Language Processing}
\newacronym{nlsp}{NLSP}{NASA Life Sciences Portal}
\newacronym{oct}{OCT}{Optical Coherence Tomography}
\newacronym{od}{OD}{Operational Data}
\newacronym{ooha}{OOHA}{On-Orbit Hearing Assessment}
\newacronym{osdr}{OSDR}{Open Science Data Repository}
\newacronym{ppg}{PPG}{Photoplethysmography}
\newacronym{rbac}{RBAC}{Role-Based Access Control}
\newacronym{rest}{REST}{Representational State Transfer}
\newacronym{roc}{ROC}{Receiver Operating Characteristic}
\newacronym{rss}{RSS}{Risk Scoring System}
\newacronym{sans}{SANS}{Spaceflight Associated Neuro-ocular Syndrome}
\newacronym{sep}{SEP}{Solar Energetic Particle}
\newacronym{slm}{SLM}{Sound Level Meter}
\newacronym{spo2}{SPO2}{Peripheral Oxygen Saturation}
\newacronym{sql}{SQL}{Structured Query Language}
\newacronym{stt}{STT}{Speech to Text}
\newacronym{t2}{T2}{Treadmill 2}
\newacronym{tts}{TTS}{Text to Speech}
\newacronym{ui}{UI}{User Interface}
\newacronym{utc}{UTC}{Coordinated Universal Time}
\newacronym{ux}{UX}{User Experience}
\newacronym{voc}{VOC}{Volatile Organic Compound}
\newacronym{wcag}{WCAG}{Web Content Accessibility Guidelines}
\newacronym{who}{WHO}{World Health Organization}

% Every entry is printed in the Acronyms list (\glsaddall in \printLOA, nasa-latex-docs.cls)


% Add bibliography file
\addbibresource{main.bib}

% Figure source line: printed under the caption; the link lives here, not in the references.
\newcommand{\figsource}[1]{\par\vspace{0.6em}\parbox{\linewidth}{\footnotesize\raggedright\textbf{Source:} #1}\par\vspace{0.4em}}
% Japanese text (status word translations): pdflatex with CJKutf8.
\usepackage{CJKutf8}
\newcommand{\ja}[1]{\begin{CJK}{UTF8}{min}#1\end{CJK}}
% Hardware figure: width, file, caption, label, source.
% Placed in the text flow (not a float), because floats followed by text overlap in this template.
\newcommand{\hwfigure}[5]{%
\par\bigskip\noindent
\begin{minipage}{\linewidth}
\centering
\includegraphics[width=#1\textwidth]{#2}
\captionof{figure}{#3}
\label{#4}
\figsource{#5}
\end{minipage}
\par\bigskip}

% Document parameters (generic placeholder values - replace with your own) %%%
\docTypeTitle[NASA Space Apps Challenge - 2026 Project] % Project or Team Name
\docTitle[Astronaut Health Intelligence (AHI)] % Technical Report Title
\docNumber[TR-0001]
\docVersion[1.1]
\docDate[November 15, 2026]
\docOrganization[AstroHealth Tech Team \\ Santo André, São Paulo - Brazil] % Team name \\ City, Country
\docHeader[Document Header]
\docFooter[] % Document Footer
\docAbstract[
    Human health monitoring is a critical component of crew safety and mission success in space exploration. As human spaceflight missions become longer and extend farther from Earth, astronauts are increasingly exposed to a combination of physiological, environmental, and operational stressors that may affect their health, well-being, and performance. Factors such as microgravity, space radiation, isolation, confinement, disrupted sleep cycles, and demanding workloads create complex health risks that require continuous monitoring and timely interpretation.
    
    Although space missions generate valuable physiological, behavioral, and environmental data, an important challenge remains: transforming heterogeneous information into a coherent understanding of an astronaut's health status. Identifying meaningful changes, recognizing abnormal patterns, and anticipating potential risks require more than data collection. They require integrated analysis, individual health baselines, and the ability to interpret trends over time.

    The \textbf{Astronaut Health Intelligence (AHI)} project proposes a predictive human health monitoring platform designed to address this challenge. Its primary objective is to transform diverse health and mission-related data into actionable insights through data integration, statistical analysis, anomaly detection, risk assessment, and artificial intelligence-assisted interpretation.

    The proposed platform assumes as its data sources the equipment in use on the International Space Station: wearables such as the Bio-Monitor, the Actiwatch Spectrum and the Dry-EEG Headband; medical equipment such as the blood pressure and ECG monitor, the Holter monitor, the pulse oximeter, the Ultrasound 2, the urine analyzer and sample collection kits; the optical coherence tomography and the On-Orbit Hearing Assessment for periodic eye and ear examinations; the CEVIS, ARED and T2 exercise devices; and the environmental monitors for air quality, carbon dioxide, combustion products, noise, radiation, water and microbes. Space weather notifications from the NASA DONKI service, accessed through the NASA Open APIs, complement the on board radiation monitors. Publicly available scientific datasets from NASA complement these sources. By organizing these sources into a standardized data model, AHI aims to provide a unified and traceable representation of astronaut health throughout a space mission.

    At its analytical core, AHI works on three levels: real-time monitoring, predictive analytics and prescriptive intelligence. It combines individual baseline modeling, time-series forecasting, statistical anomaly detection, and machine learning models such as Isolation Forest and XGBoost to identify deviations from expected physiological and behavioral patterns across ten health dimensions, including vision and hearing as two separate dimensions. Instead of relying exclusively on generalized thresholds, the platform emphasizes changes relative to each astronaut's historical baseline. This approach aims to account for individual variability and identify potentially relevant trends that could otherwise remain unnoticed.

    The platform also includes a \textbf{Large Language Model (LLM)} as an AI-assisted interpretation layer, under strict guardrails: it writes the statistical predictions in direct, plain language and powers a health assistant that answers questions by text and voice. Health monitoring cannot depend on internet access, so the target architecture runs a local model specialized in AHI on board the ISS or the Artemis spacecraft. The demonstrator uses the Claude API as a stand-in for that local model, to show how the functionality would work. Using structured analytical outputs and curated scientific context from NASA and other reliable sources, this component can explain detected patterns, summarize potential concerns, and present relevant scientific information in a human-readable format. The LLM acts as an explanatory and contextual assistance mechanism rather than an independent diagnostic system or autonomous medical decision-maker.

    The AHI platform is organized around five complementary capabilities:

    \begin{enumerate}
    \item \textbf{Data Integration:} Consolidation of physiological, behavioral, environmental, equipment, and scientific research data into a unified and standardized representation.

    \item \textbf{Health Monitoring:} Visualization of health indicators, historical trends, individual baselines, and data quality information.

    \item \textbf{Anomaly Detection:} Identification of unusual variations and deviations in physiological or behavioral patterns through statistical and analytical methods.

    \item \textbf{Predictive Risk Assessment:} Analysis of historical trends and potential risk scenarios to support early risk awareness and operational planning.

    \item \textbf{AI-Assisted Interpretation:} Generation of contextual explanations and analytical summaries supported by evidence and curated scientific references.

    \end{enumerate}

    AHI is also designed to be global and inclusive. The interface is available in English, Portuguese, Spanish, Japanese and German; access is role based, so that crew members, crew medical officers and flight surgeons each see what their role requires; and audio description lets blind and low vision users, people with dyslexia and astronauts busy with a task follow and operate the platform without looking at a screen.

    The project is grounded in publicly available NASA resources, including the \textit{NASA Open Science Data Repository (OSDR)}, Human Research Program publications, and relevant mission health measurement references. These resources provide scientific foundations and opportunities to investigate historical research data. Since access to real astronaut health records and operational telemetry is restricted, AHI explicitly distinguishes between device data, laboratory results, self report, publicly available research data, derived analytical outputs and simulated observations used for development and demonstration.

    The expected outcome is a modular and extensible software platform demonstrating how data-driven health monitoring and AI-assisted interpretation can improve the identification and understanding of potential health changes during human spaceflight missions. Its architecture is designed to support future integration with wearable sensors, environmental monitoring systems, mission data infrastructure, and additional validated analytical models.

    \textbf{AHI is designed as a health-monitoring and decision-support tool, not as a medical diagnostic system.} It does not replace qualified medical professionals, clinical evaluation, or established mission medical protocols. Instead, its purpose is to organize complex information, highlight potentially relevant changes, and provide interpretable evidence that can assist medical and mission personnel in their assessments.

    By connecting health data, scientific knowledge, and predictive analytics, \textbf{Astronaut Health Intelligence} proposes a data-driven approach to one of the central challenges of human space exploration: identifying relevant health signals early enough to support crew well-being, informed decision-making, and mission success.
] % Executive summary of the report.

% Remove the logos (the class falls back to the NASA/Artemis logos when unset)
% \docLogo[Artemis_Logo]
% \docOrgLogo[NASA_Logo]
\docLogo[SpaceApps]
\docOrgLogo[NASA_Logo]

% Remove the revision-log instruction page
\docRevisionLog[]

% Authors (the template ships with 2 default placeholder slots; empty slots hide them)
\docAuthor[1]{Igor Fernando Casita}{Team Leader \& AI Lead}{University of São Paulo (USP)}{Brazil}
\docAuthor[2]{Lucas Palhares Barbosa}{Senior Software Engineer \& AI Engineer}{FATEC Rio Preto}{Brazil}
\docAuthor[3]{Gabriel da Silva Simões}{Senior Data Scientist}{University of São Paulo (USP)}{Brazil}
\docAuthor[4]{Daniel Hiro Batista}{Data Researcher}{University of São Paulo (USP)}{Brazil}
\docAuthor[5]{Ricardo Camargo}{Senior Data Analyst}{USP ESALQ}{Brazil}
% \docAuthor[6]{Bruna Debianchi}{}{}{Brazil}

% Electronic signature on the author/approver pages ("/s/ Name"). Set only when you intend to sign.
\docSigned[true] % /s/ - eletronic signature
\docSignDate[November 15, 2026]

% Contributors and approvers: empty entries hide those pages. Fill them to show the pages.
\docContributor[1]{}{}{}{}
\docContributor[2]{}{}{}{}
\docApprover[1]{}{}{}{}
\docApprover[2]{}{}{}{}

% Main Document Content: %%%%%%%%%%%%%%%%%%%%%%%%%%%%%%%%%%%%%%%%%%%%%%%
\begin{document}

\section*{Abstract}
\phantomsection % TOC link target; \section* creates none
\addcontentsline{toc}{section}{Abstract}

Human spaceflight missions expose astronauts to complex physiological, environmental, and operational stressors that may compromise their health, well-being, and performance. Continuous health monitoring is therefore essential for identifying potentially adverse changes and supporting informed decision-making throughout space missions. However, the integration and interpretation of heterogeneous health-related data remain significant challenges.

This work presents \textbf{Astronaut Health Intelligence (AHI)}, a predictive health monitoring platform that moves from health monitoring to predictive health protection. AHI assumes as its data sources the equipment in use on the International Space Station: wearables such as the Bio-Monitor and the Actiwatch Spectrum, medical equipment such as the blood pressure and ECG monitor and the Ultrasound 2, the optical coherence tomography and the hearing assessment used for periodic eye and ear examinations, the CEVIS, ARED and T2 exercise devices, and the environmental monitors for air quality, carbon dioxide, noise and radiation. It organizes health into ten dimensions, and uses the five hazards of human spaceflight defined by NASA as the axis of its risk analysis. It follows the career radiation dose of each astronaut against the limit published in NASA-STD-3001, and it recommends preparing the radiation shelter when space weather notifications or the on board monitors indicate a solar particle event.

The analysis works on three levels. Real-time monitoring compares each astronaut with their own baseline and detects acute changes, slow drifts and missing data. Predictive analytics projects trends and estimates risks, starting with transparent statistical rules and adding machine learning models such as Isolation Forest and XGBoost as real data allow. Prescriptive intelligence recommends exercise, nutrition, sleep, well-being and medical follow-up actions, each with its evidence and an indication of whether the crew can act on board. A Large Language Model (LLM) writes the statistical predictions in plain language under strict guardrails and never computes them, and a health assistant answers questions about the data by text and voice. Internet access is useful but not essential: the demonstrator is a web application that uses the Claude API to simulate the AI language features, and, because a real system on board cannot depend on the internet, the architecture includes an on board backend and a contingency plan that keep monitoring, alerts and recommendations working without a connection, with a local model specialized in AHI for the language features. The platform is designed for a global and diverse crew, with the interface in English, Portuguese, Spanish, Japanese and German, role-based access for crew members and flight surgeons, and audio description for blind and low vision users, for people with dyslexia and for astronauts whose eyes and hands are busy with a task. Public NASA data from the Open Science Data Repository and related sources ground and test the analysis.

AHI is designed as a decision-support tool. It does not diagnose, does not prescribe, and does not replace clinical evaluation or qualified medical personnel.

\noindent\textbf{Keywords:} Astronaut Health Monitoring; Predictive Analytics; Anomaly Detection; Space Medicine; Artificial Intelligence; Human Spaceflight; International Space Station.


\section{Introduction}
\label{sec:introduction}

\subsection{Problem Statement}
\label{subsec:problem-statement}

Human space exploration introduces significant challenges to astronaut health and performance, particularly during long-duration missions beyond low Earth orbit. Unlike terrestrial environments, space missions expose crew members to a combination of environmental and operational stressors, including space radiation, isolation and confinement, altered gravity, distance from Earth, and hostile or closed environments. These conditions can contribute to physiological and behavioral changes, increasing the complexity of maintaining crew health throughout a mission \cite{nasa_human_spaceflight_risks}.

The effects of these hazards may involve multiple physiological systems. Exposure to microgravity, for example, is associated with cardiovascular adaptations and musculoskeletal changes, including bone loss. Spaceflight-related stressors may also affect immune function, sensorimotor performance, sleep, and behavioral health. The interaction between these factors makes astronaut health a multidimensional concern that requires continuous observation and the interpretation of multiple health indicators over time \cite{nasa_human_research_program,nasa_human_health_countermeasures}.

Although NASA and its international partners have established medical monitoring procedures, research programs, and technologies for collecting astronaut health information, the effective interpretation of health-related data remains an important operational consideration. Existing approaches include wearable physiological monitoring systems, medical assessments, exercise monitoring, and longitudinal health surveillance. For example, the Bio-Monitor system combines wearable technologies to collect physiological measurements aboard the International Space Station, while NASA's Lifetime Surveillance of Astronaut Health program analyzes medical and exposure data to investigate acute and long-term health outcomes \cite{canadian_space_agency_biomonitor,nasa_lsah}.

The challenge becomes particularly relevant during exploration missions in which communication delays, limited medical resources, and operational constraints may increase crew autonomy. Astronauts may need to assess their health status and recognize potentially relevant changes while performing mission activities. Consequently, health monitoring software must do more than display isolated measurements: it should support the organization, visualization, and interpretation of health indicators in a way that facilitates awareness of individual health conditions and potential changes over time.

This problem is also the central motivation of the 2026 NASA Space Apps Challenge, under the challenge \textit{Create Health Monitoring Software for Astronauts on Space Missions}. The challenge calls for the development of software capable of gathering health indicators and enabling astronauts to evaluate and act upon their health status during space missions \cite{nasa_health_monitoring_challenge_2026}.

In this context, the engineering problem addressed by this project is the development of an integrated software approach capable of consolidating heterogeneous health-related information, establishing individual reference patterns, identifying relevant deviations, and presenting analytical findings in an understandable and traceable manner. The challenge involves not only data integration, but also the distinction between normal individual variability and potentially significant changes, while preserving data quality, scientific context, and the limitations of automated analysis.

Therefore, there is an opportunity to investigate how data integration, statistical analysis, anomaly detection, and artificial intelligence-assisted interpretation can be combined into a unified astronaut health monitoring platform. Addressing this problem is the primary motivation behind the development of the \textbf{Astronaut Health Intelligence (AHI)}.

\subsection{Objectives}
\label{subsec:objectives}
The primary objective of the \textbf{Astronaut Health Intelligence (AHI)} project is to design a software platform for predictive astronaut health monitoring during space missions, integrating heterogeneous health-related data and applying analytical methods to identify relevant patterns, detect potential anomalies, and support the interpretation of individual health conditions over time.

To achieve this objective, the project establishes the following specific objectives:

\begin{enumerate}
\item \textbf{Identify and integrate relevant data sources:} define the International Space Station equipment, wearable, medical, exercise and environmental, that serves as the data source of the platform, and organize publicly available NASA research datasets that ground and test the analysis.

\item \textbf{Establish a standardized health data model:} define a structured representation for health indicators, observations, data sources, and individual historical records, ensuring consistency, traceability, and interoperability across different data categories.

\item \textbf{Implement health monitoring and visualization capabilities:} provide an interface for presenting health indicators, historical measurements, individual baselines, and relevant variations over time.

\item \textbf{Develop anomaly detection mechanisms:} investigate and implement statistical and computational methods capable of identifying deviations from individual reference patterns and highlighting potentially relevant changes in health-related indicators.

\item \textbf{Develop predictive risk assessment:} analyze temporal patterns and relationships between health indicators across ten health dimensions, to forecast trends and estimate risks for the following 24 to 72 hours.

\item \textbf{Recommend interventions:} turn identified risks into non-prescriptive recommendations on exercise, nutrition, sleep, well-being and medical follow-up, each with its evidence, autonomy and follow-up.

\item \textbf{Integrate AI-assisted contextual interpretation:} include a Large Language Model (LLM) to explain analytical findings using curated scientific knowledge and reliable reference sources, while maintaining transparency regarding the origin and limitations of generated information, and offer a health assistant that answers questions by text and voice.

\item \textbf{Treat the internet as useful, not essential:} use internet services when a connection exists, as the demonstrator does with the Claude API to simulate the AI language features, and keep monitoring, alerts, recommendations, explanations and the assistant working on board when there is no connection, through an on board backend and a contingency plan with a local language model (\ref{subsec:offline-llm}).

\item \textbf{Protect the crew from radiation:} follow the career radiation dose of each astronaut against the limit published in NASA-STD-3001, and recommend preparing the radiation shelter when a solar particle event is reported or detected.

\item \textbf{Make the platform global and inclusive:} offer the interface in English, Portuguese, Spanish, Japanese and German, control access by role, and support blind and low vision users, people with dyslexia and astronauts busy with a task through audio description and voice.

\item \textbf{Evaluate the proposed platform:} define a staged evaluation, from simulated scenarios to the replay of real public data and an operational shadow mode, with every data origin identified.

\end{enumerate}

Through these objectives, AHI aims to demonstrate the feasibility of combining heterogeneous health data, analytical techniques, and AI-assisted interpretation into a unified software approach for astronaut health monitoring. The resulting platform is intended to support health awareness and informed assessment rather than provide medical diagnoses or replace qualified medical personnel.


\section{Scientific Foundation}
\label{sec:scientific-foundation}

This section summarizes the public scientific material that shaped the design of AHI. It explains which health risks the platform is built around, which indicators it follows and why, and which published studies support those choices. It is scientific context for the design. It is not medical guidance, and it is not evidence that the platform detects any real condition (\ref{sec:limitations-ethics}).

\subsection{Spaceflight Challenges}
\label{subsec:spaceflight-challenges}

The NASA Human Research Program (HRP) groups the health threats of human spaceflight into five hazards: space radiation, isolation and confinement, distance from Earth, gravity fields, and hostile or closed environments \cite{nasa_hrp_five_hazards}. NASA notes that these hazards do not act alone, since they can feed off one another and exacerbate effects on the human body \cite{nasa_hrp_five_hazards}. AHI adopts this framework to decide what to monitor, as summarized in \ref{tab:hazards}.

\begin{table}[htbp]
\centering
\vspace{0.8em}
\caption{The five HRP hazards and the data AHI uses for each one.}
\label{tab:hazards}
\small
\begin{tabular}{p{3.2cm}p{5.6cm}p{5.4cm}}
\hline
\textbf{Hazard} & \textbf{What the public sources describe} & \textbf{What AHI represents} \\
\hline
Space radiation & Galactic cosmic rays, solar particle events and trapped particles \cite{nasa_srag_space_radiation}; one of the most hazardous aspects of spaceflight \cite{nasa_hrp_five_hazards} & Area radiation monitors and personal dosimeters: dose rate, mission dose and career dose; space weather notifications \\
Isolation and confinement & Can reduce cognitive and behavioral functioning \cite{nasa_behavioral_health_risk} & Self report of mood, stress and workload; sleep; cognition tests \\
Distance from Earth & Communication delay up to 20 minutes one way to Mars, and a need for crew self-sufficiency \cite{nasa_hrp_distance_from_earth}; about 4 to 24 minutes depending on planetary positions \cite{nasa_space_communications_seven_things} & Communication delay of the mission profile; autonomy of each recommendation; on board first steps \\
Gravity fields & Bone mineral density loss of 1\% to 1.5\% per month, headward fluid shift, effects on the eyes \cite{nasa_hrp_gravity_fields} & CEVIS, ARED and T2 performance; heart rate and blood pressure; ultrasound and periodic examinations \\
Hostile or closed environments & Temperature, pressure, lighting, noise and microbes inside a closed ecosystem \cite{nasa_hrp_five_hazards} & CO\textsubscript{2}, volatile organic compounds, combustion products, noise, temperature, humidity, water and microbial monitoring \\
\hline
\end{tabular}
\vspace{1.2em}
\end{table}

\medskip\noindent\textbf{Radiation.} Outside Earth's atmosphere and magnetic field, crews are exposed to galactic cosmic rays, solar particle events and trapped radiation. In low Earth orbit, flights cross the South Atlantic Anomaly six or seven times a day \cite{nasa_srag_space_radiation}. NASA links space radiation to increased lifetime cancer risk, central nervous system effects and degenerative disease \cite{nasa_why_space_radiation_matters}. NASA-STD-3001 sets a career effective dose limit below 600~mSv and a solar particle event limit below 250~mSv \cite{nasa_ochmo_tb020_radiation}. These limits are expressed as effective dose. AHI compares a dosimeter value with them only when the value is reported as effective dose, and it does not estimate radiation health risk.

\medskip\noindent\textbf{Closed environment.} In a spacecraft carbon dioxide depends on scrubbing equipment. The NASA-STD-3001 technical brief reports headaches among ISS crew and links CO\textsubscript{2} level to them, and it limits the average one hour CO\textsubscript{2} partial pressure in the habitable volume to 3~mmHg \cite{nasa_ochmo_tb004_co2}. AHI uses this standard as the cited limit for its cabin CO\textsubscript{2} rule.

\medskip\noindent\textbf{Noise.} A closed spacecraft is noisy: fans, pumps and life support run day and night. NASA-STD-3001 sets acoustic limits to preclude noise-related hearing loss and interference with communication, sleep and alarms, and it requires that noise exposure be monitored and quantified for each crew member \cite{nasa_ochmo_tb035_acoustics}. The 24 hour noise exposure at the ear of a crew member is limited to a noise dose of at most 100, equivalent to an 8 hour time-weighted average of 85~dBA \cite{nasa_ochmo_tb035_acoustics}. On the ISS, crew-worn acoustic dosimeters showed that NASA noise exposure limits were exceeded in 41\% of the dosimetry measurements made since 2008, without accounting for the effect of hearing protection \cite{limardo2021noise}. AHI uses the NASA-STD-3001 noise dose as the cited limit for its personal noise exposure rule.

\medskip\noindent\textbf{Gravity.} Microgravity unloads bones and muscles and shifts body fluids toward the head \cite{nasa_hrp_gravity_fields}. The cardiovascular system is strongly gravity dependent, and its altered function causes problems after flight \cite{hughson2018heart}. These effects motivate monitoring exercise and cardiovascular signals against each person's own normal.

\medskip\noindent\textbf{Distance and autonomy.} The longer the distance, the longer the wait for help. The Mars delay cited above means a question to Earth and its answer can take close to 50 minutes at the greatest distance \cite{nasa_space_communications_seven_things}. This is why AHI places on board steps first when the Earth link is long. \ref{subsec:mission-scenarios} compares how these hazards change from launch to the ISS, the Moon and Mars.

\subsection{Mission Scenarios}
\label{subsec:mission-scenarios}

The five hazards are present in every mission, but their weight changes with where the crew goes and for how long. NASA contrasts a three-day trip to the Moon with a journey of roughly three years for Mars, during which the crew would live in three-eighths of Earth's gravity for up to two years \cite{nasa_hrp_five_hazards}. Radiation follows the same direction: as crews travel farther from the protective shields of the Earth, they are exposed to the full radiation spectrum, and beyond low Earth orbit space radiation may place them at significant risk for radiation sickness \cite{nasa_why_space_radiation_matters}. The longer and the farther the mission, the larger the accumulated exposure and the longer the crew must manage alone.

AHI considers four scenarios: \textbf{launch}, which is a phase of every mission rather than a destination, and three destinations, the \textbf{ISS} in low Earth orbit, the \textbf{Moon} and a journey to \textbf{Mars}. \ref{tab:mission-scenarios} summarizes what the public sources describe for each one and what AHI does in it.

\begin{itemize}
\item \textbf{Launch and the other dynamic phases.} During launch the dominant hazards are mechanical. A NASA-STD-3001 technical brief lists, for sustained acceleration, limited movement and breathing and possible unconsciousness; for transient acceleration, mostly traumatic injuries, vertebral injuries being the most common; and for whole-body vibration, fatigue, discomfort, degraded vision and reduced fine motor control \cite{nasa_ochmo_tb024_occupant}. The same brief notes that G loading is expected to peak at 3.8~Gx nominally on ascent and even higher during re-entry, that back and leg pain and incontinence have been experienced on the launch pad, and it cites the requirement that the time on the back with feet elevated in launch configuration not exceed 3 hours and 15 minutes [V2 6113]. NASA-STD-3001 also requires that vehicle and crew acceleration parameters be collected during ascent, abort, entry, descent, landing, postlanding and extravehicular activity, to correlate them with any injuries [V2 6111] \cite{nasa_ochmo_tb024_occupant}.
\item \textbf{ISS.} The station flies inside the magnetic field of the Earth, but crosses the South Atlantic Anomaly six or seven times a day \cite{nasa_srag_space_radiation}. A dose equivalent of 731~\textmu Sv per day measured inside the ISS is used as a reference in the lunar study below \cite{zhang2020lunar}. For spacecraft close to Earth, the communication delay is almost negligible \cite{nasa_space_communications_seven_things}. The ISS is the reference scenario of this report and the basis of its hardware premise (\ref{sec:hardware-sensors}).
\item \textbf{Moon.} The Moon has neither a substantial atmosphere nor a global magnetic field, so the crew is exposed to galactic cosmic rays, to sporadic solar particle events and to neutrons and gamma rays produced in the lunar soil. The Chang'E~4 lander measured an average dose equivalent of 1369~\textmu Sv per day on the lunar surface, and its authors estimate the daily dose from galactic cosmic rays to be about 2.6 times the dose inside the ISS \cite{zhang2020lunar}. The link to Earth can be interrupted: during Artemis~II, the Orion spacecraft lost contact for about 40 minutes behind the Moon \cite{ctv_artemis2_blackout}.
\item \textbf{Mars.} During the 253-day cruise of the Curiosity rover to Mars, an on board detector measured the radiation inside the spacecraft; from these data, the dose equivalent of even the shortest round trip with current propulsion and comparable shielding was estimated at 0.66~$\pm$~0.12~Sv \cite{zeitlin2013mars}. This is a dose equivalent and not the effective dose of the NASA career limit (\ref{subsec:spaceflight-challenges}), so it is not a comparison with the limit, but it shows the scale of the exposure. The communication delay is about 4 to 24 minutes one way \cite{nasa_space_communications_seven_things}, and the crew must be largely self-sufficient \cite{nasa_hrp_distance_from_earth}.
\end{itemize}

The radiation values above come from different instruments, shielding and solar conditions, so they show the direction of the risk and not a precise ranking. AHI uses them only as context and never as limits; the only radiation limits it applies are those of NASA-STD-3001 (\ref{subsec:spaceflight-challenges}).

\begin{table}[htbp]
\centering
\vspace{0.8em}
\caption{Mission scenarios considered by AHI.}
\label{tab:mission-scenarios}
\footnotesize
\begin{tabular}{p{1.7cm}p{4.3cm}p{3.6cm}p{4.6cm}}
\hline
\textbf{Scenario} & \textbf{Main risks in the public sources} & \textbf{Radiation and link to Earth} & \textbf{Focus of AHI} \\
\hline
Launch (and entry, landing) & Sustained and transient acceleration, vibration, time in launch orientation \cite{nasa_ochmo_tb024_occupant} & Short phase; mechanical hazards dominate & Records the phase on the timeline, the time in launch orientation and a check-in after launch and after landing; imports the acceleration record after flight; never acts on the vehicle \\
ISS (low Earth orbit) & The five hazards over long stays; inside the magnetic field of the Earth & 731~\textmu Sv per day measured inside the ISS \cite{zhang2020lunar}; almost negligible delay & Reference scenario: all devices of \ref{sec:hardware-sensors}, ground support available \\
Moon & Radiation without a magnetic field, reduced gravity, surface activity & 1369~\textmu Sv per day on the surface \cite{zhang2020lunar}; blackouts behind the Moon \cite{ctv_artemis2_blackout} & Shelter rule including surface activity; smaller device set; operation during blackouts \\
Mars & Roughly three years, reduced gravity, isolation, self-sufficiency \cite{nasa_hrp_five_hazards} & 0.66~Sv dose equivalent for a round trip \cite{zeitlin2013mars}; 4 to 24 minutes one way & Maximum on board autonomy, career dose projection, on board steps for every high item \\
\hline
\end{tabular}
\vspace{1.2em}
\end{table}

\subsection{Health Indicators}
\label{subsec:health-indicators}

AHI follows indicators grouped by body system, plus the cabin environment. \ref{tab:indicators} lists the systems, how each is measured on real missions according to public sources, and what AHI uses.

\begin{table}[htbp]
\centering
\vspace{0.8em}
\caption{Monitored systems, measurement on real missions and AHI use.}
\label{tab:indicators}
\small
\begin{tabular}{p{2.6cm}p{6.4cm}p{5.2cm}}
\hline
\textbf{System} & \textbf{Measurement on real missions (public sources)} & \textbf{In AHI} \\
\hline
Cardiovascular & Bio-Monitor smart shirt records heart electrical activity, blood pressure and physical activity \cite{canadian_space_agency_biomonitor}; ISS monitoring includes blood pressure, ECG and heart rate \cite{nasa_jsc_exercise_countermeasures} & Heart rate, heart rate variability, blood pressure, ECG \\
Respiratory & Bio-Monitor records breathing rate and blood oxygen saturation \cite{canadian_space_agency_biomonitor}; a crew wearable vest records heart rate and breathing during sleep \cite{nasa_wearable_tech_iss} & Blood oxygen, breathing rate and volume \\
Thermal & Bio-Monitor records skin temperature; core temperature needs dedicated devices such as the T-Mini headband \cite{nasa_wearable_tech_iss} & Skin temperature \\
Sleep and activity & Wrist actigraphy (Actiwatch Spectrum) tracks motion and light to estimate sleep \cite{nasa_wearable_tech_iss} & Sleep duration, efficiency and stages, activity, light \\
Musculoskeletal & Exercise on CEVIS, treadmills and ARED \cite{nasa_jsc_exercise_countermeasures}; bone density by densitometry before and after flight \cite{lang2004bone} & Exercise performance against the personal baseline; examination results \\
Behavior and cognition & Standard Measures include cognition tests \cite{nasa_standard_measures}; the Cognition battery has 10 tests \cite{basner2015cognition} & Self report and cognition test results \\
Hydration & Periodic fluid samples \cite{nasa_standard_measures}, not a continuous sensor & Urinalysis, blood chemistry and reported intake \\
Vision & Vision screening (visual acuity, Amsler grid, survey), retinal photography, optical coherence tomography and ocular ultrasound before, during and after flight \cite{brunstetter2023brief} & Structured eye examination results as point measurements; visual symptoms reported by the crew \\
Hearing & On-Orbit Hearing Assessment for crew members on orbit for more than 30 days \cite{coste2024hearing}; crew-worn acoustic dosimetry \cite{limardo2021noise} & Hearing thresholds by frequency and ear; personal noise exposure \\
Bone density & Densitometry before and after flight \cite{lang2004bone} & Examination results as point measurements \\
\hline
\end{tabular}
\vspace{1.2em}
\end{table}

Three points about this table matter for how the platform works.

\begin{itemize}
\item \textbf{Real telemetry is not public.} NASA publishes how crews are monitored, but no continuous physiological telemetry of astronauts is available to the public. Data from the devices are a design premise of AHI (\ref{sec:hardware-sensors}).
\item \textbf{Some risks are measured only by periodic exams.} Bone density, muscle volume and the eye findings of spaceflight associated neuro-ocular syndrome are assessed by examinations and not continuously \cite{lang2004bone,nasa_hrr_sans}. AHI receives their results as point measurements.
\item \textbf{Some indicators are weaker than others.} Self report is not a validated clinical instrument, and estimated values such as sleep stages are marked as estimates.
\end{itemize}

\subsection{Supporting Scientific Studies}
\label{subsec:supporting-studies}

The studies below support the choice of indicators and the decision to compare each person with their own normal. Each is cited only for the statement it supports. None of them is evidence about the performance of AHI.

\medskip\noindent\textbf{Cardiovascular.} Hughson et al. review how the extraterrestrial environment affects the cardiovascular system, including postflight problems, a possible late cardiovascular risk from radiation, and exercise as a countermeasure \cite{hughson2018heart}. In the NASA Twins Study, a longitudinal before, during and after design with a matched ground control, the flight twin showed carotid artery changes among other findings \cite{garrettbakelman2019twins}.

\medskip\noindent\textbf{Sleep.} In a study of 64 Space Shuttle and 21 ISS astronauts, sleep measured by wrist actigraphy was short in flight, at about 6.09 hours per night on ISS missions, and also during preflight training \cite{barger2014sleep}. The authors compare flight with each person's own ground intervals, which is the same logic as the personal baseline in AHI.

\medskip\noindent\textbf{Bone, muscle and exercise.} Lang et al. measured areal bone mineral density loss of 0.9\% per month at the spine and 1.4\% to 1.5\% per month at the hip in 14 ISS crew members \cite{lang2004bone}. Trappe et al. found a 13\% average decrease in calf muscle volume in 9 crew members after about six months, despite an exercise program that varied between individuals \cite{trappe2009exercise}. Resistive exercise with adequate nutrition maintained bone in most regions in one study \cite{smith2012bone}, while a later study found only partial protection and no prevention of trabecular loss at the hip \cite{sibonga2019resistive}. Vico and Hargens link skeletal changes with fluid shifts, the vascular system and metabolism \cite{vico2018skeletal}. Because exercise regimens differ between people, AHI judges exercise performance against each person's own usual sessions.

\medskip\noindent\textbf{Bed rest as a ground analog.} LeBlanc et al. review the skeletal data from the U.S. and Russian space programs together with ground-based disuse studies, in which volunteers remain in bed rest as an analog of spaceflight \cite{leblanc2007skeletal}. The review covers bone mineral density measured by X-ray absorptiometry and CT, bone related hormones and biochemical markers of bone turnover, and calcium excretion and balance. Its central point for this project is that bed rest studies give information complementary to the spaceflight data and allow countermeasures to bone loss to be tested on the ground \cite{leblanc2007skeletal}. The same review notes that many questions remain about bone loss, bone recovery and the factors that affect them \cite{leblanc2007skeletal}. Bed rest is an analog and not a replacement for flight data: it does not reproduce radiation, the closed environment or the other spaceflight hazards. AHI follows bone through exercise data and periodic examination results.

\medskip\noindent\textbf{Cognition and behavior.} The Cognition test battery was developed and validated for astronauts, and its validation shows that practice effects differ between tests, which makes normative data important \cite{basner2015cognition}. NASA states that isolation and confinement can reduce cognitive and behavioral functioning \cite{nasa_behavioral_health_risk}. AHI uses the results of cognition tests and self report as indicators and does not interpret them clinically.

\medskip\noindent\textbf{Eyes.} Spaceflight associated neuro-ocular syndrome (SANS) comprises findings such as optic disc edema and globe flattening \cite{lee2020sans}. NASA's Human Research Roadmap reports that 70\% of crew members experience one or more hallmark signs \cite{nasa_hrr_sans}, and a 2023 clinical update estimates about 66\% of long duration astronauts show the earliest indication \cite{brunstetter2023sans}. The leading hypothesis links SANS to the headward fluid shift \cite{nasa_fluid_shifts_study}. A NASA briefing for astronaut candidates adds that 18\% of ISS crew members develop clinically concerning SANS, that cases are typically asymptomatic unless blurred vision develops, that there is currently no evidence of permanent vision loss in any astronaut, and that severity is related to mission duration \cite{brunstetter2023brief}. Because most cases have no symptoms, SANS is followed with scheduled examinations: the same briefing lists vision screening (visual acuity, Amsler grid and a survey), retinal photography, optical coherence tomography (OCT) and ocular ultrasound before flight, around flight days 30 and 90, 30 days before return and after flight, with MRI and a comprehensive eye examination only on the ground, and it describes OCT as the most critical tool to detect and monitor SANS, on board the ISS since 2013 \cite{brunstetter2023brief}. NASA follows these changes with optical coherence tomography on board: in one study, the same Spectralis system used before flight was used on the ISS around flight days 10, 30, 60, 90, 120 and 150 \cite{macias2020ocular}. AHI places structured eye examination results on the timeline and does not interpret images.

\medskip\noindent\textbf{Hearing.} Hearing loss is part of the acoustics risk followed by NASA \cite{nasa_ochmo_tb035_acoustics}. Crew members on orbit for more than 30 days perform the On-Orbit Hearing Assessment, which flight surgeons on the ground review for significant changes. A NASA poster reports a low frequency hearing threshold shift in 40\% of these tests and in 60\% of crew members, and notes that such low frequency changes are not typically caused by noise exposure \cite{coste2024hearing}. AHI therefore follows exposure and hearing thresholds side by side, compares each result with the person's own preflight audiogram, and reports any link between them as co-occurrence, never as cause.

\medskip\noindent\textbf{Mission monitoring.} NASA Standard Measures and Artemis~II Standard Measures list the standardized measures collected before, during and after missions \cite{nasa_standard_measures,nasa_artemis_ii_standard_measures}. Astroskin, the technology behind the Bio-Monitor, was developed under Canadian Space Agency contracts \cite{csa_wireless_body_sensors}. In testing with 16 crew members in the HERA analog, 24 hour monitoring showed lessons on data loss and Bluetooth connectivity, with no comparison against laboratory standard devices \cite{toscano2017astroskin}. This is the reason AHI treats missing data as \texttt{NO\_DATA} and never as normal.

\medskip\noindent\textbf{Scope of the evidence.} The sources above justify what AHI monitors and why a personal baseline is appropriate. They do not by themselves justify numeric thresholds. Every fixed threshold in AHI must cite a standard, such as NASA-STD-3001, or be reviewed by a flight surgeon, and none is presented as an official NASA limit without a source (\ref{sec:future-work}).


\section{Data Sources}
\label{sec:data-sources}

This report describes the final scope of AHI. In that scope the platform draws on three families of data. The first and main family is \textbf{operational data from equipment in use on the International Space Station (ISS)}: wearables, medical equipment, exercise equipment and environmental monitors. The second is \textbf{space weather notifications} from NASA, which warn of solar events before or while their radiation reaches the spacecraft. The third is \textbf{public NASA research data}, used to ground the analysis in real measurements and to support the development and evaluation of the models. \ref{tab:data-sources} summarizes both families and the role of each.

\medskip\noindent\textbf{Design premise.} AHI assumes that every operational value comes from the ISS equipment described in \ref{sec:hardware-sensors}, for both the crew and the cabin. Only equipment that public NASA, Canadian Space Agency (CSA) or Japan Aerospace Exploration Agency (JAXA) sources describe as in use on the ISS is part of this premise. The premise defines what the platform is designed to receive. It does not mean that the project has access to these devices or to their data. Access to real crew data would require authorization, documentation of each interface and a data handling agreement, and it is treated in \ref{sec:limitations-ethics}.

\begin{table}[htbp]
\centering
\vspace{0.8em}
\caption{Data families in the final scope of AHI and their role.}
\label{tab:data-sources}
\small
\begin{tabular}{p{3.4cm}p{5.2cm}p{5.4cm}}
\hline
\textbf{Family} & \textbf{Examples} & \textbf{Role in AHI} \\
\hline
Wearables & Bio-Monitor, Actiwatch Spectrum, Dry-EEG Headband & Continuous or session based physiology, activity and sleep \\
Medical equipment & BP/ECG monitor, Holter Monitor 2, pulse oximeter, Ultrasound 2, optical coherence tomography, fundoscope, vision screening, On-Orbit Hearing Assessment, urine analyzer, blood and saliva collection & Scheduled and periodic clinical measurements, eye and hearing examinations \\
Exercise equipment & CEVIS, ARED, T2 with COLBERT & Physical performance, training load and adherence \\
Environmental monitors & Air Quality Monitor, CO\textsubscript{2} monitors, combustion product analyzer, acoustic dosimeters and sound level meters, radiation monitors and dosimeters, water and microbial monitoring & Context for interpreting physiology; cabin hazards; personal noise and radiation exposure \\
Space weather notifications & DONKI notifications of solar flares, coronal mass ejections and solar energetic particle events, through the NASA Open APIs & Early warning for the radiation shelter recommendation \\
Self report on station computers & Daily check-in, sleep report, cognition tests & Behavioral health, fatigue and cognition \\
Public NASA research data & OSDR, EDA, RadLab, NLSP, LSDA & Reference ranges, model development and retrospective evaluation \\
\hline
\end{tabular}
\vspace{1.2em}
\end{table}

\subsection{Operational Data from ISS Equipment}
\label{subsec:operational-data}

The medical structure of the ISS is the Crew Health Care System. It brings together health maintenance, environmental monitoring and exercise equipment, and the JAXA medical operations page describes its main components \cite{jaxa_medops_system}. AHI organizes the operational data into four categories that follow this structure.

\begin{itemize}
\item \textbf{Wearables} collect physiological data during daily activities without stopping the work of the crew. They are the main source of continuous data.
\item \textbf{Medical equipment} produces clinical measurements at scheduled times, such as blood pressure, an electrocardiogram (ECG) recording, an ultrasound examination or a urine analysis.
\item \textbf{Exercise equipment} describes the daily countermeasure sessions that crews perform to limit the loss of bone, muscle and cardiovascular fitness \cite{nasa_jsc_exercise_countermeasures}.
\item \textbf{Environmental monitors} describe the cabin. They do not measure health directly, but air quality, carbon dioxide, noise and radiation affect the crew and are needed to interpret physiological changes \cite{nasa_toxicology_envchem,nasa_ochmo_tb035_acoustics}. Acoustic dosimeters and personal radiation dosimeters are worn by the crew, so they also measure the exposure of each person.
\end{itemize}

Not every device is used by every crew member during the whole mission. Several wearables are research hardware used during specific investigations \cite{nasa_wearable_tech_iss,nasa_sleep_on_station}. AHI therefore treats each data stream as present only when the corresponding device is in use, and a stream that is expected but absent is reported as missing data and never as normal.

\medskip\noindent\textbf{Self report.} Mood, perceived stress, workload and sleep reports are not produced by hardware. NASA Standard Measures already collects sleep reports and cognition tests in flight \cite{nasa_standard_measures}. AHI assumes that such inputs are entered on station computers or on the tablet of the Bio-Monitor system and are labeled as self report.

\subsection{Collection Frequency}
\label{subsec:collection-frequency}

Not every indicator is collected at the same rate. \ref{tab:frequency} lists the cadence the platform is designed to handle for each kind of data. These cadences are integration targets for AHI. They are not a NASA operational specification.

\begin{table}[htbp]
\centering
\vspace{0.8em}
\caption{Cadence that AHI is designed to handle for each kind of data. Design targets, not a NASA specification.}
\label{tab:frequency}
\small
\begin{tabular}{p{4.6cm}p{4.2cm}p{5.2cm}}
\hline
\textbf{Kind of data} & \textbf{Cadence in AHI} & \textbf{Typical source} \\
\hline
Heart rate & Continuous or near continuous & Bio-Monitor \\
Breathing & Continuous or per session & Bio-Monitor \\
Blood oxygen saturation & Periodic or continuous, depending on the device & Bio-Monitor, pulse oximeter \\
ECG & Per session or 24 hour recording & BP/ECG monitor, Holter Monitor 2 \\
Sleep & Each rest period & Actiwatch Spectrum, Dry-EEG Headband \\
Physical activity & During the day & Actiwatch Spectrum, Bio-Monitor \\
Exercise performance & Per session & CEVIS, ARED, T2 \\
Blood pressure & Scheduled measurements & BP/ECG monitor, Bio-Monitor \\
Cabin CO\textsubscript{2} and air quality & Continuous & CO\textsubscript{2} monitors, Air Quality Monitor \\
Radiation & Continuous area dose; accumulated personal dose & Radiation monitors, personal dosimeters \\
Blood, saliva and urine & Scheduled collections & Collection kits, urine analyzer \\
Ultrasound & Scheduled examinations & Ultrasound 2 \\
Eye examination & Before flight, around flight days 30 and 90, 30 days before return, after flight & Vision screening, fundoscope, optical coherence tomography, Ultrasound 2 \\
Visual symptoms & Any day, in the daily check-in & Self report \\
Hearing thresholds & Periodic, for flights longer than 30 days & On-Orbit Hearing Assessment \\
Personal noise exposure & 24 hour crew-worn sessions; area surveys & Acoustic dosimeters, sound level meter \\
Space weather & On each new notification & DONKI, relayed by the ground node \\
Microbial and water quality & Scheduled sampling & Environmental Health System \\
\hline
\end{tabular}
\vspace{1.2em}
\end{table}

A consequence of these cadences is that each signal has its own definition of fresh and stale. A heart rate that has not arrived for several minutes is stale, while a blood sample result is expected only at the next scheduled collection. AHI stores the expected cadence of every stream and uses it to decide when data become missing.

\subsection{Space Weather Notifications}
\label{subsec:space-weather}

Personal dosimeters measure the dose already received. A solar particle event can raise the dose rate within hours, and NASA-STD-3001 requires real-time monitoring of space weather and a method to alert the crew, locally and remotely, when radiation levels are expected to exceed acceptable levels \cite{nasa_ochmo_tb020_radiation}. AHI therefore adds a second, external source of radiation information: the \textbf{Space Weather Database Of Notifications, Knowledge, Information (DONKI)} of the NASA Community Coordinated Modeling Center \cite{nasa_donki}. DONKI catalogs space weather events, such as solar flares, coronal mass ejections, solar energetic particle events and geomagnetic storms, together with the notifications issued by its forecasters.

\medskip\noindent\textbf{Access through the NASA Open APIs.} DONKI is available through the NASA Open APIs at \texttt{api.nasa.gov} \cite{nasa_open_apis}, with one endpoint per event type and one for notifications. Every request carries an API key. NASA issues personal keys free of charge, and a shared \texttt{DEMO\_KEY} exists for exploration with a much lower rate limit. When tested by the team on 2026-11-15, the notifications endpoint returned current messages, and the event endpoints rejected date ranges longer than 60 days, so the adapter requests short windows. The key is a secret of the ground node: it is stored in the server configuration, never in the interface code, in the repository or in this report, and it can be rotated without a new release.

\medskip\noindent\textbf{Useful, not essential.} DONKI is an internet service. AHI uses it when a connection exists, but its essential functions must keep working without one (\ref{sec:software-architecture}). Only the ground node calls the API. It filters the notifications that matter for the crew, stores them with their message identifier and issue time, and sends them on board through the same store and forward link as any other message. On board, the radiation rule does not wait for DONKI: a sudden rise in the dose rate measured by the station radiation monitors triggers the same recommendation by itself (\ref{sec:analytical-methodology}). DONKI gives earlier warning when the link exists, and the on board monitors give protection when it does not.

\medskip\noindent\textbf{An honest label.} DONKI messages state that the NOAA Space Weather Prediction Center is the official source of space weather forecasts for the United States Government. In a real mission, AHI would receive its alerts from the space weather support of the mission. DONKI is the public NASA source that the design and the demonstration use, and every notification keeps its source and message identifier.

\subsection{Public NASA Research Data}
\label{subsec:public-data}

No public source offers continuous physiological telemetry of astronauts. Public NASA repositories nevertheless hold real measurements that AHI uses in three ways: as historical reference shown next to crew data, as material to develop and test the analysis, and as a retrospective test of the engine outside simulated data.

\begin{itemize}
\item \textbf{NASA Open Science Data Repository (OSDR).} The main NASA repository for space life science, with studies, files and metadata \cite{nasa_about_osdr}. Among its human studies are the Inspiration4 complete blood count (OSD-569) and comprehensive metabolic panel (OSD-575), with four crew members at three time points before launch and four after return and no in-flight samples \cite{osdr_osd569,osdr_osd575}. They give AHI real laboratory values with units, the format that blood collection results take in the platform.
\item \textbf{Environmental Data Application (EDA).} ISS cabin carbon dioxide, temperature, humidity and radiation dose for the periods of research payload missions \cite{osdr_eda}. AHI uses it to replay real cabin periods. The data describe the cabin during a payload mission and not a given crew.
\item \textbf{RadLab.} An access point for radiation telemetry from several instruments \cite{osdr_radlab}. It was unavailable when tested on 2026-11-14, so its coverage and format are still to be checked.
\item \textbf{NASA Life Sciences Portal (NLSP)} and \textbf{Life Sciences Data Archive (LSDA).} Discovery portals for historical life science and human health data, including analog studies \cite{nasa_nlsp,nasa_lsda}. Many datasets require a formal request. They are the path to sleep, cognition and analog data. Analog participants are never presented as astronauts.
\item \textbf{NASA Standard Measures} and \textbf{Artemis~II Standard Measures.} Reference material and not datasets. They list the measures collected before, during and after missions and justify the catalog of indicators \cite{nasa_standard_measures,nasa_artemis_ii_standard_measures}.
\end{itemize}

A caution applies to all public data. It is not safe to assume that individual medical data of astronauts are publicly available, and in general they are not. Public datasets are small, often from short missions, and rarely include in-flight samples. They support grounding and testing, but they are not enough to train clinical models by themselves (\ref{sec:artificial-intelligence}).

\subsection{Data Provenance}
\label{subsec:data-provenance}

Every observation carries a \texttt{data\_origin} field, shown in the interface as a badge, so that a reader always knows where a value comes from:

\begin{itemize}
\item \texttt{DEVICE}: received from ISS equipment through a device adapter (\ref{sec:hardware-sensors}).
\item \texttt{MANUAL}: entered by a person, such as the daily check-in or a sleep report.
\item \texttt{LAB}: a result of a sample collected in flight and analyzed later, with the collection time and the result time kept apart.
\item \texttt{DERIVED}: computed by the platform, such as baselines, risk scores and forecasts.
\item \texttt{REAL\_PUBLIC}: public NASA data, such as research datasets shown as historical reference and space weather notifications from DONKI.
\item \texttt{SIMULATED}: produced by the simulator used for development, testing and demonstration. Simulated values are never mixed with crew data.
\end{itemize}

Each record also carries its source device, unit, measurement time, quality status and activity context. Imported public datasets follow a reproducible path: the original file is downloaded from an allow list of NASA hosts, its SHA-256 hash and download time are recorded, and a manifest keeps the source address, version, population, design and license. Values are kept exactly as reported, nothing is imputed, and fields the source does not state are recorded as not stated.


\section{Solution Overview}
\label{sec:solution-overview}

This section describes AHI as a whole in its final scope: the value it offers, how it is organized, who uses it and what it deliberately does not do. The following sections give the details: the equipment (\ref{sec:hardware-sensors}), the analysis (\ref{sec:analytical-methodology}), the role of artificial intelligence (\ref{sec:artificial-intelligence}) and the architecture (\ref{sec:software-architecture}).

\subsection{Value Proposition}
\label{subsec:value-proposition}

The guiding idea of AHI is to move \textbf{from health monitoring to predictive health protection}. AHI is an intelligent platform that integrates physiological, behavioral and environmental data to monitor the health of astronauts continuously, anticipate risks and recommend interventions, contributing to safer and more sustainable space missions. It is not one more vital signs display. It is designed to answer five questions about each crew member:

\begin{enumerate}
\item What is the current condition?
\item How is that condition evolving?
\item Is there a risk of developing a problem?
\item What could happen if no action is taken?
\item Which actions could reduce that risk?
\end{enumerate}

Three ideas distinguish the approach:

\begin{itemize}
\item \textbf{Predictive, not reactive.} The platform does not wait for severe symptoms. It identifies trends of deterioration and estimates future risk from the individual history and from scientific knowledge.
\item \textbf{Personalized.} Each astronaut has a personal history, physiology and baseline. AHI compares each person with their own normal instead of applying the same fixed limits to everyone.
\item \textbf{Autonomous decision support.} On missions far from Earth, the ability to identify risks, show the evidence and suggest actions on board matters more as the communication delay grows. AHI supports the decisions of the crew and of the flight surgeon and never replaces them.
\end{itemize}

\subsection{Three Levels of Analysis}
\label{subsec:three-levels}

The platform works on three levels, summarized in \ref{tab:levels}. Each level builds on the one before it.

\begin{table}[htbp]
\centering
\vspace{0.8em}
\caption{The three levels of analysis of AHI.}
\label{tab:levels}
\small
\begin{tabular}{p{3.6cm}p{1.6cm}p{4.6cm}p{4.2cm}}
\hline
\textbf{Level} & \textbf{Time} & \textbf{Examples} & \textbf{Result} \\
\hline
1. Real-time health monitoring & Present & Heart rate above the personal pattern, falling oxygen saturation, altered blood pressure, short sleep, high CO\textsubscript{2} & Anomaly detection and immediate alerts \\
2. Predictive health analytics & Future & Accumulated fatigue, cardiovascular deconditioning, loss of strength, worsening sleep & Risk estimates before a condition becomes critical \\
3. Prescriptive health intelligence & Action & Adjust training, rest and recovery, hydration plan, medical follow-up & Recommended interventions to reduce the predicted risk \\
\hline
\end{tabular}
\vspace{1.2em}
\end{table}

\subsection{Eight Health Dimensions}
\label{subsec:health-dimensions}

To avoid a narrow view limited to heart rate and blood pressure, AHI organizes monitoring into eight health dimensions: cardiovascular; respiratory and oxygenation; muscle and bone; immune; neurological and sensorimotor; sleep and circadian rhythm; psychological and behavioral; and metabolic, renal and nutritional. Each dimension has its own indicators, analyses and risks, described in \ref{sec:analytical-methodology}. Most dimensions are followed with frequent data. The immune and the metabolic, renal and nutritional dimensions depend largely on samples and laboratory analysis, so they are followed periodically.

\subsection{Hazards as the Axis of Risk}
\label{subsec:hazard-axis}

NASA groups the threats to human health in spaceflight into five hazards \cite{nasa_hrp_five_hazards}. AHI uses them as the central axis of its risk analysis, as shown in \ref{tab:hazard-monitoring}.

\begin{table}[htbp]
\centering
\vspace{0.8em}
\caption{The five hazards of human spaceflight and what AHI follows for each.}
\label{tab:hazard-monitoring}
\small
\begin{tabular}{p{3.0cm}p{5.4cm}p{5.6cm}}
\hline
\textbf{Hazard} & \textbf{What can happen} & \textbf{What AHI follows} \\
\hline
Space radiation & Cell damage, cancer risk, cardiovascular and neurological effects & Dose, exposure history and biological indicators \\
Isolation and confinement & Stress, psychological changes, sleep problems & Mood, sleep, stress and team interaction \\
Distance from Earth & Difficulty of obtaining immediate medical care & Medical autonomy, alert severity and forecast of deterioration \\
Gravity fields & Bone and muscle loss, cardiovascular and vestibular changes & Strength, exercise, blood pressure, heart rate, balance \\
Hostile or closed environments & Exposure to CO\textsubscript{2}, contaminants and microbes & CO\textsubscript{2}, air quality, temperature, humidity and immune indicators \\
\hline
\end{tabular}
\vspace{1.2em}
\end{table}

The hazards should not be analyzed in isolation, since they can feed off one another \cite{nasa_hrp_five_hazards}. A typical chain is the following: an astronaut starts to sleep less, sleep quality worsens, physiological recovery indicators change, exercise performance declines, and prolonged stress may influence the immune response. Each signal may look acceptable on its own. AHI is designed to recognize when several apparently independent indicators begin to point to the same deterioration, and this is where its analytical and artificial intelligence layers add the most value.

\subsection{Users}
\label{subsec:users}

\begin{table}[htbp]
\centering
\vspace{0.8em}
\caption{Users of AHI and what each needs.}
\label{tab:users}
\small
\begin{tabular}{p{3.6cm}p{10.4cm}}
\hline
\textbf{User} & \textbf{Need} \\
\hline
Crew member & A plain view of their own status, the actions they can take on board and a quick daily check-in \\
Crew medical officer & A view of the whole crew, alerts by priority and the evidence behind each one \\
Flight surgeon & Trends, forecasts and recommendations to review, approve or reject, with the full history \\
Mission control & The health status of the crew as an operational factor for planning \\
Researcher & Traceable data with origin and quality, for retrospective analysis \\
\hline
\end{tabular}
\vspace{1.2em}
\end{table}

Access is role based: each user sees what their role requires (\ref{sec:software-architecture}).

\subsection{Main Modules}
\label{subsec:main-modules}

AHI is organized into six modules, detailed in \ref{sec:software-architecture}:

\begin{enumerate}
\item \textbf{Device integration:} one adapter per device that converts native data into canonical observations.
\item \textbf{Health data platform:} storage of current and historical data with origin, quality and audit.
\item \textbf{Analytics and prediction engine:} baselines, anomaly detection, risk models and forecasts.
\item \textbf{Recommendation engine:} interventions linked to the identified risks, with evidence and follow-up.
\item \textbf{Explanation layer:} the Claude API, and a local model on board when contact with Earth is lost, turn structured results into plain explanations, under strict guardrails.
\item \textbf{Dashboards and review:} the crew, crew health, predictive and recommendation screens, and the flight surgeon review workflow.
\end{enumerate}

\subsection{What AHI Does Not Do}
\label{subsec:not-in-scope}

Some limits are part of the design and not missing features:

\begin{itemize}
\item \textbf{It does not diagnose.} AHI identifies risk indicators and requests appropriate evaluation. It does not name diseases. In particular it does not diagnose depression or anxiety from facial expressions, messages or heart rate.
\item \textbf{It does not prescribe medication or treatment.} Recommendations are operational actions and requests for review, never drugs or doses.
\item \textbf{It does not replace the flight surgeon or mission protocols.} High and critical items always require human review.
\item \textbf{It does not present a forecast as a certainty.} Every prediction carries its uncertainty and is never presented as the certainty that a disease will occur.
\item \textbf{It does not invent official limits.} Status comes from personal baselines. Any fixed threshold must come from a cited standard or from review by a flight surgeon.
\item \textbf{It does not treat missing data as normal.} An absent or stale signal is reported as no data.
\end{itemize}


\section{Hardware and Sensors}
\label{sec:hardware-sensors}

This section defines the equipment that AHI assumes as its data sources. The premise of the report is that all operational data, for the crew and for the cabin, come from these devices. The selection criterion is strict: \textbf{a device is included only if public NASA, CSA or JAXA sources describe it as used on the ISS.} Devices that were tested only on the ground or in analog habitats are left out, even when they are relevant.

The descriptions below repeat what the cited sources state. Where a source does not establish which signal a device exposes, the report says so, and AHI does not assume more. The existence of a device on board does not imply that a data interface is open or available to third parties (\ref{subsec:integration-assumptions}). \ref{fig:hw-biomonitor-info} to \ref{fig:hw-sv104a} show each device as published by NASA, CSA and JAXA.

\subsection{Wearables}
\label{subsec:wearables}

Wearables are the most important source for AHI, because they collect data without interrupting the activities of the crew.

\medskip\noindent\textbf{Bio-Monitor (CSA).} A wearable system made of a smart shirt, a headband and a tablet application. It records heart pulse and electrical activity, blood pressure, breathing rate and volume, skin temperature, blood oxygen saturation and physical activity, and it was launched to the ISS in December 2018 \cite{canadian_space_agency_biomonitor}. It was developed under CSA contracts from the Astroskin technology \cite{csa_wireless_body_sensors}. Bio-Monitor is the closest real example of the concept behind AHI: several sensors combined in one wearable system that consolidates physiological information. It is the primary continuous source of the platform.

\medskip\noindent\textbf{Actiwatch Spectrum.} A waterproof sleep and wake monitor worn on the wrist that measures motion and ambient light. It is used to analyze circadian rhythms, sleep and wake patterns and activity \cite{nasa_issmp_hardware,nasa_wearable_tech_iss}. The Actiwatch Spectrum Plus is in use on the station for the Standard Measures study \cite{nasa_sleep_on_station}. The same class of device, a wrist actigraph that records movement and light, was adopted by NASA for the Artemis II crew (\ref{fig:hw-actiwatch}).

\medskip\noindent\textbf{Dry-EEG Headband (Dreams).} A headband used by the ESA investigation Dreams that lets astronauts monitor their sleep wherever they sleep and over several nights \cite{nasa_sleep_on_station}. It is research hardware used during an investigation, not a device worn by every crew member. AHI uses it, when present, as a richer sleep source than wrist actigraphy.

\medskip\noindent\textbf{Excluded: Astroskin as a separate garment.} Astroskin is the commercial technology behind the Bio-Monitor. The Astroskin evaluation found in the sources was carried out with 16 crew members in the HERA ground analog \cite{toscano2017astroskin}, so it is represented in AHI only through the Bio-Monitor that flies on the ISS.

\hwfigure{0.92}{images/hardware/bio-monitor-infographic.jpg}{Bio-Monitor infographic: what the smart shirt and headband monitor, and how data flow from the crew member to an on-orbit controller and to scientists on the ground.}{fig:hw-biomonitor-info}{Canadian Space Agency, Bio-Monitor, \url{https://www.asc-csa.gc.ca/eng/sciences/bio-monitor.asp}}

\hwfigure{0.70}{images/hardware/bio-monitor-saint-jacques.jpg}{CSA astronaut David Saint-Jacques with the Bio-Monitor on board the ISS.}{fig:hw-biomonitor-photo}{Canadian Space Agency, Bio-Monitor, \url{https://www.asc-csa.gc.ca/eng/sciences/bio-monitor.asp}}

\hwfigure{0.85}{images/hardware/artemis-ii-actigraph.jpg}{Artemis II crew (left) and the wrist actigraph adopted by NASA to track their sleep, activity and light exposure (right). Developed at the University of S\~ao Paulo and commercialized by Condor Instruments, it is a wrist actigraph of the same class as the Actiwatch Spectrum used on the ISS.}{fig:hw-actiwatch}{Particle News, NASA Uses USP-Born Actigraph to Monitor Artemis II Crew, \url{https://particle.news/story/nasa-uses-usp-born-actigraph-to-monitor-artemis-ii-crew}}

\hwfigure{0.75}{images/hardware/dreams-headband.jpg}{Sleep Monitoring in Space with Dry-EEG Headband (Dreams) hardware in the Columbus module of the ISS, 20 September 2021 (NASA image iss065e423775).}{fig:hw-dreams}{NASA Image and Video Library, \url{https://images.nasa.gov/details/iss065e423775}}

\subsection{Medical Equipment}
\label{subsec:medical-equipment}

Medical equipment produces scheduled, clinical measurements. Its data arrive as sessions or as results, not as a continuous stream.

\begin{itemize}
\item \textbf{Blood pressure and ECG monitor.} JAXA lists a blood pressure and ECG monitor among the ISS medical equipment \cite{jaxa_medops_system}, and NASA lists health monitoring of blood pressure, ECG and heart rate among the systems it maintains on the ISS \cite{nasa_jsc_exercise_countermeasures}.
\item \textbf{Holter Monitor 2.} An ambulatory ECG that measures the electrical activity of the heart for up to 24 hours through electrode pads on the torso \cite{nasa_issmp_hardware}.
\item \textbf{Pulse oximeter.} A finger device that measures pulse and oxygen saturation, used with ISS medical project hardware \cite{nasa_issmp_hardware}.
\item \textbf{Ultrasound 2.} A General Electric Vivid q unit used for cardiac, muscle, vessel and blood flow analysis, with features to estimate muscle volume changes and analyze cardiac strain \cite{nasa_issmp_hardware}. It is used regularly on the station, and its scans of the heart and veins can be viewed by doctors on the ground in real time \cite{nasa_ultrasound_2025}. Ocular ultrasound is also part of the in-flight eye surveillance for SANS \cite{brunstetter2023brief}.
\item \textbf{Urine analyzer.} JAXA lists the Urolax urine analyzer among the ISS medical equipment \cite{jaxa_medops_system}. The source does not detail the parameters it reports, so AHI treats its output as a set of named results with units, defined when the interface is documented.
\item \textbf{Blood, saliva and urine collection.} The ISS medical project carries blood collection kits, saliva session packs, urine collection kits and a refrigerated centrifuge \cite{nasa_issmp_hardware}, and blood collection and saliva samples are part of NASA Standard Measures \cite{nasa_standard_measures}. These are not sensors. Samples are collected at set times and analyzed later, so their results reach AHI with a delay. NASA immune research uses such samples to study stress hormones and the reactivation of dormant viruses \cite{nasa_immune_biomarkers}.
\item \textbf{Optical coherence tomography (OCT).} A near infrared imaging device that produces cross sections of the retina and the optic nerve head. NASA uses it to follow the eye changes of spaceflight associated neuro-ocular syndrome: a Spectralis OCT of the same model used before flight is on the ISS, and in one study crew members were imaged around flight days 10, 30, 60, 90, 120 and 150 \cite{macias2020ocular}. NASA describes OCT as its most critical tool to detect and monitor SANS, on board the ISS since 2013 \cite{brunstetter2023brief}. AHI receives the structured measurements of each examination, such as retinal and nerve fiber layer thickness, per eye, and never the image itself.
\item \textbf{Fundoscope (retinal photography).} A handheld camera that photographs the retina and the optic nerve head. Retinal photography is part of the in-flight eye surveillance for SANS \cite{brunstetter2023brief}, and the examination can be made with remote support from doctors on the ground (\ref{fig:hw-fundoscope}). AHI receives the structured findings of the reviewer per eye, not the photograph.
\item \textbf{Vision screening.} Visual acuity, the Amsler grid, which shows distorted or missing areas of central vision, and a vision survey, performed at the same points of the flight as the imaging \cite{brunstetter2023brief}. These tests measure how well the person sees, while OCT, fundoscopy and ultrasound show the structure of the eye, and both are needed because structural changes can exist without symptoms. AHI does not use a webcam, a phone camera or an on-screen letter chart as a substitute for these examinations.
\item \textbf{On-Orbit Hearing Assessment (OOHA).} A hearing test performed on board with headphones, which measures the pure tone hearing thresholds of crew members who stay on orbit longer than 30 days. Flight surgeons on the ground review the results for significant changes. The ISS system, called EarQ, was upgraded in 2021 to the KUDUwave software, which allows a more comprehensive assessment \cite{coste2024hearing}. AHI receives each audiogram as thresholds per frequency and per ear, in decibels hearing level.
\end{itemize}

\hwfigure{0.45}{images/hardware/bp-ecg-monitor-jaxa.jpg}{Blood pressure and ECG monitor of the ISS medical operations equipment.}{fig:hw-bpecg}{JAXA, Medical Operations, \url{https://humans-in-space.jaxa.jp/en/biz-lab/med-in-space/healthcare/medops/system/}}

\hwfigure{0.55}{images/hardware/holter-monitor-2-components.jpg}{Flight certified Holter Monitor 2 components: recorder, electrode lead set and memory cards. The Holter Monitor 2 is an ambulatory ECG that records the electrical activity of the heart for up to 24 hours.}{fig:hw-holter}{NASA Ames Research Center, Holter Monitor (STS-126), \url{https://www.nasa.gov/ames/space-biosciences/holter-monitor-sts-126}}

\hwfigure{0.92}{images/hardware/pulse-oximeter-ccfp-slide.jpg}{Cerebral Cochlear Fluid Pressure Analyzer, whose finger pulse oximeter measures pulse and oxygen saturation.}{fig:hw-oximeter}{NASA ISS Medical Project, On-Orbit Hardware (2016), slide 4, \url{https://www.nasa.gov/sites/default/files/atoms/files/issmp_on-orbit_hardware.pdf}}

\hwfigure{0.85}{images/hardware/ultrasound-2-cassidy-parmitano.jpg}{NASA astronaut Chris Cassidy, Expedition 36, performs an ultrasound on ESA astronaut Luca Parmitano in the Columbus laboratory for the Spinal Ultrasound investigation (NASA image ISS036-E-035750).}{fig:hw-ultrasound}{NASA, image ISS036-E-035750, public domain, via Wikimedia Commons, \url{https://commons.wikimedia.org/wiki/File:ISS-36_Luca_Parmitano_and_Chris_Cassidy_in_the_Columbus_lab.jpg}}

\hwfigure{0.45}{images/hardware/urolax-jaxa.jpg}{Urolax urine analyzer of the ISS medical operations equipment.}{fig:hw-urolax}{JAXA, Medical Operations, \url{https://humans-in-space.jaxa.jp/en/biz-lab/med-in-space/healthcare/medops/system/}}

\hwfigure{0.92}{images/hardware/blood-collection-slide.jpg}{Blood collection consumables of the ISS Medical Project: open blood tube kit, blood sample kits and insulated blood pouch.}{fig:hw-blood}{NASA ISS Medical Project, On-Orbit Hardware (2016), slide 17, \url{https://www.nasa.gov/sites/default/files/atoms/files/issmp_on-orbit_hardware.pdf}}

\hwfigure{0.75}{images/hardware/oct-adenot.jpg}{ESA astronaut Sophie Adenot, Expedition 74, with the optical coherence tomography device used for eye examinations on the ISS, 3 April 2026. Doctors on the ground can follow the examination in real time (NASA image iss074e0433179).}{fig:hw-oct}{NASA Image and Video Library, credit ESA/Sophie Adenot, \url{https://images.nasa.gov/details/iss074e0433179}}

\hwfigure{0.75}{images/hardware/fundoscope-aunon-chancellor.jpg}{NASA astronaut Serena Au\~n\'on-Chancellor examines her eye with a fundoscope on board the ISS with remote support from doctors on the ground, 13 July 2018 (NASA image iss056e096896).}{fig:hw-fundoscope}{NASA Image and Video Library, \url{https://images.nasa.gov/details/iss056e096896}}

\hwfigure{0.75}{images/hardware/hearing-assessment.jpg}{Roscosmos cosmonaut Anna Kikina, Expedition 68, wears headphones for a hearing assessment inside the Quest airlock of the ISS, 21 October 2022 (NASA image iss068e017876).}{fig:hw-hearing}{NASA Image and Video Library, \url{https://images.nasa.gov/details/iss068e017876}}

\subsection{Exercise Equipment}
\label{subsec:exercise-equipment}

Crews exercise every day to limit the loss of bone, muscle and cardiovascular fitness caused by microgravity. NASA lists three regular devices on the ISS: the CEVIS cycle ergometer, the ARED resistive exercise device and the T2 treadmill with COLBERT, which allows walking and running in microgravity \cite{nasa_jsc_exercise_countermeasures}. The software that works with ARED and T2 collects, stores and displays data and shows the exercise prescription to the crew member \cite{nasa_jsc_exercise_countermeasures}.

\begin{itemize}
\item \textbf{CEVIS.} A cycle ergometer with electronically controlled workload. AHI uses workload, duration, cadence and the heart rate during the session, which is the basis for following cardiovascular fitness against the personal baseline.
\item \textbf{ARED.} Resistive exercise similar to weight training. AHI uses load, repetitions, sets and duration to follow strength and adherence to training.
\item \textbf{T2 with COLBERT.} A treadmill. AHI uses speed, duration and the load applied by the harness to follow aerobic capacity.
\end{itemize}

\medskip\noindent\textbf{Caution.} The public sources used here confirm the three devices and that their software collects data, but they do not list the exact fields each device exports. The fields above are the ones AHI is designed to consume, to be confirmed against the interface documentation of each device.

\hwfigure{0.42}{images/hardware/cevis-teal.jpg}{Astronaut Loral O'Hara (Increment 70) exercises on the new teal CEVIS ergometer on board the ISS.}{fig:hw-cevis}{NASA Johnson Space Center, Exercise Countermeasures, \url{https://www.nasa.gov/reference/jsc-exercise-countermeasures/}}

\hwfigure{0.42}{images/hardware/ared.jpg}{Koichi Wakata, Expedition 38 flight engineer, exercising on the Advanced Resistive Exercise Device (ARED) in Node 3.}{fig:hw-ared}{NASA Johnson Space Center, Exercise Countermeasures, \url{https://www.nasa.gov/reference/jsc-exercise-countermeasures/}}

\hwfigure{0.42}{images/hardware/t2-colbert.jpg}{JAXA astronaut Satoshi Furukawa, Expedition 70, jogs on the T2 treadmill with COLBERT.}{fig:hw-colbert}{NASA Johnson Space Center, Exercise Countermeasures, \url{https://www.nasa.gov/reference/jsc-exercise-countermeasures/}}

\subsection{Environmental Monitoring}
\label{subsec:environmental-monitoring}

The health of a crew does not depend only on the body. A closed cabin can cause or worsen physiological problems, and AHI uses its monitors as context for every crew signal.

\begin{itemize}
\item \textbf{Air Quality Monitor (AQM).} Two AQM units use gas chromatography and differential mobility spectrometry to detect and quantify 23 target volatile organic compounds \cite{nasa_toxicology_envchem}.
\item \textbf{CO\textsubscript{2} monitors.} A CO\textsubscript{2} monitor uses non-dispersive infrared to detect elevated carbon dioxide \cite{nasa_toxicology_envchem}. The Personal CO\textsubscript{2} Monitor investigation used sensors attached to clothing to measure the immediate surroundings of the wearer \cite{nasa_wearable_tech_iss}. Elevated CO\textsubscript{2} has been linked to headaches in ISS crews \cite{nasa_ochmo_tb004_co2}.
\item \textbf{Combustion product analyzer.} The compound specific analyzer for combustion products uses electrochemical sensors to detect compounds produced by fire \cite{nasa_toxicology_envchem}.
\item \textbf{Radiation monitors and personal dosimeters.} Radiation measurement devices installed in the ISS are monitored in real time, and each astronaut wears a personal dosimeter that records the accumulated dose \cite{jaxa_medops_system}. Examples on the ISS include the ESA Active Dosimeter and the CSA EVARM dosimeters \cite{nasa_wearable_tech_iss}. AHI uses the area dose rate for the radiation shelter rule and the personal dose for the career dose of each astronaut (\ref{sec:analytical-methodology}).
\item \textbf{Acoustic Monitor (crew-worn acoustic dosimeter).} NASA's name for the Svantek SV 102A+, a Class 1 dual channel noise dosimeter used on the ISS since November 2017, with three calibrated units on board at all times. Crew members wear it for a continuous 24 hours, covering work and sleep, to measure their typical noise exposure, and it is also placed at fixed locations. Its data are downloaded to a station computer through USB. The ISS flight rule on noise uses the 16 hour work period and 8 hour sleep period equivalent levels measured by crew-worn dosimeters, and sets actions such as the use of hearing protection \cite{limardo2021noise}. A newer personal dosimeter, the SV104A, was tested in a technology demonstration (\ref{fig:hw-sv104a}).
\item \textbf{Sound level meter (SLM).} Sound level meter measurements and acoustic noise dosimetry are routinely performed on the ISS \cite{nasa_ochmo_tb035_acoustics}. The meter surveys the sound level and its frequency content at points of the station, which shows where the noise comes from.
\item \textbf{Water and microbial monitoring.} The Environmental Health System includes monitoring of water, air and microbes \cite{jaxa_medops_system}. Water quality is checked with a colorimetric kit for biocide concentration and a total organic carbon analyzer \cite{nasa_toxicology_envchem}.
\item \textbf{Cabin temperature and humidity.} Station environmental data, of the kind published for research mission periods in the Environmental Data Application \cite{osdr_eda}.
\end{itemize}

\hwfigure{0.75}{images/hardware/aqm.jpg}{Air quality monitors used for volatile organic compound detection, positioned in the U.S. Lab on the ISS (highlighted in red).}{fig:hw-aqm}{NASA, Toxicology and Environmental Chemistry, \url{https://www.nasa.gov/directorates/esdmd/hhp/toxicology-and-environmental-chemistry/}}

\hwfigure{0.75}{images/hardware/personal-co2-monitor.jpg}{One of the wearable carbon dioxide monitors clipped to the wall near a crew sleeping compartment.}{fig:hw-co2}{NASA, Wearable Tech for Space Station Research, \url{https://www.nasa.gov/missions/station/iss-research/wearable-tech-for-space-station-research/}}

\hwfigure{0.80}{images/hardware/csa-cp.jpg}{The compound specific analyzer for combustion products (CSA-CP), used to monitor for evidence of fires or smoldering events on the ISS.}{fig:hw-csacp}{NASA, Toxicology and Environmental Chemistry, \url{https://www.nasa.gov/directorates/esdmd/hhp/toxicology-and-environmental-chemistry/}}

\hwfigure{0.75}{images/hardware/active-dosimeter.jpg}{ESA astronaut Thomas Pesquet holds one of the mobile units of the Active Dosimeter.}{fig:hw-dosimeter}{NASA, Wearable Tech for Space Station Research, \url{https://www.nasa.gov/missions/station/iss-research/wearable-tech-for-space-station-research/}}

\hwfigure{0.40}{images/hardware/toca.jpg}{The U.S. Total Organic Carbon Analyzer (TOCA), used to test water quality in real time on the ISS.}{fig:hw-toca}{NASA, Toxicology and Environmental Chemistry, \url{https://www.nasa.gov/directorates/esdmd/hhp/toxicology-and-environmental-chemistry/}}

\hwfigure{0.65}{images/hardware/environment-sampling-jaxa.jpg}{Image used by JAXA to illustrate station environment monitoring of air, water and microbes.}{fig:hw-env-jaxa}{JAXA, Medical Operations, \url{https://humans-in-space.jaxa.jp/en/biz-lab/med-in-space/healthcare/medops/system/}}

\hwfigure{0.75}{images/hardware/sound-level-meter.jpg}{NASA astronaut Steve Swanson, Expedition 40 commander, uses a sound level meter in the Destiny laboratory of the ISS, 1 August 2014 (NASA image ISS040-E-086654).}{fig:hw-slm}{NASA Image and Video Library, \url{https://images.nasa.gov/details/iss040e086654}}

\hwfigure{0.30}{images/hardware/sv104a-dosimeter.jpg}{The SV104A personal noise dosimeter, which measures noise dose and noise levels from 55 to 140 dB on board the ISS, tested in the Wireless Acoustics technology demonstration. Image courtesy of SVANTEK (NASA image jsc2025e044837).}{fig:hw-sv104a}{NASA Image and Video Library, \url{https://images.nasa.gov/details/jsc2025e044837}}

\subsection{Measured Variables}
\label{subsec:measured-variables}

\ref{tab:variables} relates each device to the variables AHI consumes and to the health dimensions of \ref{sec:analytical-methodology} that use them.

\begin{table}[htbp]
\centering
\vspace{0.8em}
\caption{ISS equipment assumed as data source, variables consumed by AHI and dimensions served.}
\label{tab:variables}
\footnotesize
\begin{tabular}{p{3.0cm}p{6.4cm}p{4.6cm}}
\hline
\textbf{Device} & \textbf{Variables consumed by AHI} & \textbf{Dimensions} \\
\hline
Bio-Monitor & Heart rate, ECG derived heart rate variability, blood pressure, breathing rate and volume, skin temperature, blood oxygen saturation, activity & Cardiovascular, respiratory, sleep \\
Actiwatch Spectrum & Motion, activity level, ambient light, sleep and wake timing & Sleep and circadian rhythm \\
Dry-EEG Headband & Sleep stages, sleep duration, awakenings & Sleep and circadian rhythm \\
BP/ECG monitor & Systolic and diastolic pressure, ECG, heart rate & Cardiovascular \\
Holter Monitor 2 & 24 hour ECG, rhythm & Cardiovascular \\
Pulse oximeter & Oxygen saturation, pulse & Respiratory \\
Ultrasound 2 & Cardiac, vascular, muscle and eye examination results & Cardiovascular, musculoskeletal, vision \\
Optical coherence tomography & Retinal and nerve fiber layer thickness, structured findings per eye & Vision \\
Fundoscope & Structured findings of the retinal photograph per eye & Vision \\
Vision screening & Visual acuity per eye, Amsler grid result, vision survey & Vision \\
On-Orbit Hearing Assessment & Hearing thresholds per frequency and ear & Hearing \\
Urine analyzer & Urinalysis results & Metabolic, renal and nutrition \\
Blood and saliva samples & Blood counts, chemistry, stress hormones, immune markers & Immune, metabolic \\
CEVIS & Workload, duration, heart rate during exercise & Cardiovascular, musculoskeletal \\
ARED & Load, repetitions, sets, duration & Musculoskeletal \\
T2 with COLBERT & Speed, duration, harness load & Musculoskeletal, cardiovascular \\
AQM, CO\textsubscript{2} monitors & Volatile organic compounds, cabin and personal CO\textsubscript{2} & Respiratory, environment \\
Radiation monitors, dosimeters & Area dose rate, personal mission dose, career dose & Environment \\
Acoustic Monitor, sound level meter & Equivalent noise levels for work and sleep, 24 hour noise dose, levels by frequency band & Hearing, sleep, environment \\
Water and microbial monitoring & Water quality, microbial sampling results & Immune, environment \\
\hline
\end{tabular}
\vspace{1.2em}
\end{table}

Some variables are measured directly and others are estimated. Heart rate variability is derived from the ECG, and sleep stages from the headband are estimates. Each value carries a quality status that records whether it is measured, estimated or questionable.

\subsection{Imaging and Periodic Examinations}
\label{subsec:imaging}

Several risks are followed by periodic examinations and not by continuous sensors. Bone mineral density is measured by densitometry before and after flight \cite{lang2004bone,leblanc2007skeletal}. The eye findings of spaceflight associated neuro-ocular syndrome are followed with optical coherence tomography, ultrasound and other examinations \cite{nasa_hrr_sans,macias2020ocular}, and hearing with the On-Orbit Hearing Assessment \cite{coste2024hearing}. AHI does not process images or audio. It receives the result of each examination as point measurements, with the date, the device, the eye or ear concerned and the examiner's structured findings, and places them on the timeline of the crew member as examination observations of the same canonical model. Each new result is compared with the person's preflight examination and with their previous results, never with a population average.

\subsection{Device Priority}
\label{subsec:device-priority}

Integrating every device at once is not realistic. \ref{tab:device-priority} gives the order in which AHI integrates the equipment. The first three already support monitoring and prediction for the dimensions with the most frequent data.

\begin{table}[htbp]
\centering
\vspace{0.8em}
\caption{Integration priority of the ISS equipment.}
\label{tab:device-priority}
\small
\begin{tabular}{p{1.4cm}p{4.0cm}p{8.6cm}}
\hline
\textbf{Priority} & \textbf{Equipment} & \textbf{Main data} \\
\hline
1 & Bio-Monitor & Heart rate, ECG derived variability, breathing, blood oxygen, skin temperature \\
2 & Actiwatch Spectrum & Sleep, activity and light \\
3 & CEVIS, ARED, T2 & Physical performance and training load \\
4 & BP/ECG monitor & Systolic and diastolic blood pressure, ECG \\
5 & CO\textsubscript{2} monitors and AQM & Environmental exposure \\
6 & Radiation dosimeters & Area dose rate, personal and career dose \\
7 & Acoustic Monitor, sound level meter & Personal noise exposure \\
8 & Vision screening, optical coherence tomography, fundoscope, On-Orbit Hearing Assessment & Eye and hearing examinations \\
9 & Sample results, ultrasound, urine analyzer, Holter, Dry-EEG Headband & Periodic clinical and research data \\
\hline
\end{tabular}
\vspace{1.2em}
\end{table}

\subsection{Conceptual Integration with ISS Equipment}
\label{subsec:integration-assumptions}

\textbf{AHI does not consume data from any ISS device.} The project has no access to the equipment, to its interfaces or to crew data, and every value in the AHI demonstrator comes from a simulator. Even so, the team studied how a real integration would work, from the moment a device records a value to the moment that value is stored in the database of the platform. The purpose of this study is that the data model, the database and the adapters are designed for the real data path from the start, so that a simulator can later be replaced by a real feed without changing the analysis. This subsection describes that study at a theoretical and conceptual level. Where a step rests on a public source it is cited, and every other step is a design assumption of the team.

\medskip\noindent\textbf{What the public sources establish about the data path.} The ISS research hardware already provides much of the infrastructure an integration needs. The Human Research Facility racks connect to the ISS Ethernet and video services, and their laptops are used to control equipment, to collect and store data and crew notes, and to provide uplink and downlink \cite{nasa_issmp_hardware}. The Ultrasound 2 connects to the Human Research Facility rack to allow real-time downlink of its video \cite{nasa_issmp_hardware}, and doctors on the ground can view the scans in real time \cite{nasa_ultrasound_2025}. The Actiwatch Spectrum kit includes a dock, a USB cable and the Actiware software \cite{nasa_issmp_hardware}, which means its data are recorded on the watch and transferred later. The Bio-Monitor includes a tablet application \cite{canadian_space_agency_biomonitor}, and the software of ARED and T2 collects, stores and displays exercise data \cite{nasa_jsc_exercise_countermeasures}. In ground testing of the Astroskin technology, 24 hour monitoring showed lessons about data loss and wireless connectivity \cite{toscano2017astroskin}. Radiation devices installed in the station are monitored in real time, while personal dosimeters record the accumulated dose \cite{jaxa_medops_system}. The Acoustic Monitor records a 24 hour session and is then connected to a station computer through USB to download its data \cite{limardo2021noise}.

\medskip\noindent\textbf{Integration patterns.} From these facts, the team identified four ways in which data leave a device. Each pattern needs a different kind of adapter, as summarized in \ref{tab:integration-patterns}.

\begin{itemize}
\item \textbf{Near real-time stream.} The device or its host application produces values continuously, and the adapter reads them as they arrive. This is the target pattern for the Bio-Monitor and for cabin environmental telemetry.
\item \textbf{Session file.} The device records a session and the data are transferred afterwards as a file, through a dock, a cable or the host computer. This is the pattern of the Actiwatch Spectrum, the Holter Monitor 2, the Dry-EEG Headband, the Acoustic Monitor and the exercise sessions of CEVIS, ARED and T2. The adapter reads the exported file and keeps the original measurement times.
\item \textbf{Examination record.} A clinical examination, such as an ultrasound scan, an optical coherence tomography, a retinal photograph, a vision screening, a hearing assessment, a blood pressure and ECG session or a urinalysis, produces a result reviewed by a person. The adapter receives the structured findings, not the raw image or waveform, with the date, the device and the examiner.
\item \textbf{Laboratory result.} Blood, saliva and urine samples are collected on board, stored and analyzed later, often on the ground. The result reaches AHI days or weeks after the collection, and the adapter keeps the collection time and the result time apart.
\end{itemize}

\begin{table}[htbp]
\centering
\vspace{0.8em}
\caption{Conceptual integration pattern for each ISS device. Sourced facts are cited in the text; the patterns are design assumptions of the team.}
\label{tab:integration-patterns}
\footnotesize
\begin{tabular}{p{3.0cm}p{2.6cm}p{4.4cm}p{4.0cm}}
\hline
\textbf{Device} & \textbf{Pattern} & \textbf{Point of capture (assumed)} & \textbf{Expected latency} \\
\hline
Bio-Monitor & Stream & Export from the tablet application & Seconds to minutes; gaps when out of range \\
Actiwatch Spectrum & Session file & Actiware export after docking & Hours, after each download \\
Dry-EEG Headband & Session file & Export after each night & Hours \\
Holter Monitor 2 & Session file & Recording download after up to 24 hours & One day \\
BP/ECG monitor & Examination record & Session result on the host computer & Minutes after the session \\
Pulse oximeter & Examination record & Session result & Minutes \\
Ultrasound 2 & Examination record & Structured findings after review & Hours to days \\
Urine analyzer & Examination record & Analyzer result & Minutes to hours \\
Optical coherence tomography & Examination record & Structured measurements per eye after review & Hours to days \\
Fundoscope & Examination record & Structured findings per eye after review & Hours to days \\
Vision screening & Examination record & Results entered on the station computer & Minutes \\
On-Orbit Hearing Assessment & Examination record & Audiogram on the station computer & Minutes; ground review later \\
Blood and saliva samples & Laboratory result & Ground laboratory report & Days to weeks \\
CEVIS, ARED, T2 & Session file & Exercise software export per session & Minutes to hours \\
CO\textsubscript{2} monitors, AQM & Stream & Station environmental telemetry & Minutes \\
Radiation monitors & Stream & Station radiation telemetry & Minutes \\
Personal dosimeters & Session file & Periodic dosimeter readout & Days or longer \\
Acoustic Monitor & Session file & USB download to a station computer after the 24 hour session & About one day \\
Sound level meter & Examination record & Survey results per location & Hours \\
DONKI (not a device) & External notification & Ground node polling of the NASA Open APIs, relayed on board & Minutes on the ground, plus link delay \\
Water and microbial monitoring & Laboratory result & Analysis report & Days \\
\hline
\end{tabular}
\vspace{1.2em}
\end{table}

\medskip\noindent\textbf{From the device to the database.} Whatever the pattern, every value follows the same ingestion pipeline on the on board node, described in \ref{sec:software-architecture}:

\begin{enumerate}
\item \textbf{Capture.} The adapter receives the native payload: a stream message, an exported file, an examination record or a laboratory report.
\item \textbf{Raw landing.} The original payload is stored unchanged, with its SHA-256 hash, the adapter version and the time of arrival. Nothing is lost if a later step fails, and any stored value can be traced back to the payload it came from.
\item \textbf{Parsing and validation.} The adapter reads the payload as data and never executes it. It checks the structure, the units and the physical plausibility of each value. A value that fails is kept with the status invalid, not discarded silently.
\item \textbf{Normalization.} Units are converted to the canonical unit of each metric, while the original value and unit are kept. Times are converted to UTC, and the clock offset of the device, when known, is recorded.
\item \textbf{Identification.} The value is linked to its subject, through the pseudonym assigned to the device for that period, and to its device, mission phase and activity context.
\item \textbf{Deduplication.} Each record carries an idempotent key built from the device, the metric and the measurement time, so that a file read twice or a batch resent after a link interruption creates no duplicate.
\item \textbf{Quality.} The adapter sets the quality status (valid, questionable, invalid or missing), the confidence and whether the value is estimated, and it updates the state of the device: online, degraded, stale or offline.
\item \textbf{Storage.} The canonical observation is written to the database of the on board node, and the analysis is triggered for the affected person and metric.
\item \textbf{Synchronization.} The record enters the store and forward queue to the ground node, which holds the full history (\ref{subsec:data-flow}).
\end{enumerate}

\medskip\noindent\textbf{Database design for real data.} The study shaped the database so that it can hold the four patterns without special cases. \ref{tab:integration-db} lists the main storage areas. Continuous observations are stored in a time series table, partitioned by time, because they make up most of the volume. Sessions, examinations and laboratory results are stored as parent records, and their individual values are canonical observations linked to them, so that a heart rate from a CEVIS session and a heart rate from the Bio-Monitor are queried in the same way but keep their different origin.

\begin{table}[htbp]
\centering
\vspace{0.8em}
\caption{Storage areas of the AHI database designed for real device data.}
\label{tab:integration-db}
\footnotesize
\begin{tabular}{p{3.4cm}p{6.6cm}p{4.0cm}}
\hline
\textbf{Storage area} & \textbf{Content} & \textbf{Fed by} \\
\hline
Raw payloads & Original payload, hash, adapter version, arrival time & Every adapter \\
Devices & Device identifier, model, version, assigned subject and period, state, battery and signal & Adapter status reports \\
Observations & Canonical observations (\ref{tab:observation}), partitioned by time & Streams and all child values \\
Sessions & Start, end, device, subject and context of a recording or exercise session & Session files \\
Examinations & Type, date, device, examiner and structured findings & Examination records \\
Laboratory results & Sample type, collection time, result time, analyte, value, unit and reported reference range & Laboratory results \\
Ingestion log & Accepted and rejected records with the reason for each rejection & The ingestion pipeline \\
\hline
\end{tabular}
\vspace{1.2em}
\end{table}

\medskip\noindent\textbf{Problems a real integration must solve.} The study also listed the difficulties that real data would bring and that the design already accounts for:

\begin{itemize}
\item \textbf{Late and out of order data.} Session files and laboratory results arrive long after the measurement. The database keeps the measurement time separate from the arrival time, and the analysis recomputes the affected baselines and trends when older data arrive.
\item \textbf{Clock differences.} Each device has its own clock. Times are converted to UTC and the known offset of each device is recorded, so that values from different devices can be aligned.
\item \textbf{Gaps and connectivity.} Wireless links can drop, as the Astroskin ground tests showed \cite{toscano2017astroskin}. A gap is stored as missing data, and the expected cadence of each stream decides when a signal becomes stale.
\item \textbf{Device changes.} Hardware is replaced over time, as with the new teal CEVIS ergometer \cite{nasa_jsc_exercise_countermeasures}. Each adapter is versioned, and a new model needs a new adapter version, not a change to the database or to the analysis.
\item \textbf{Shared devices.} Several devices are used by more than one crew member. The assignment of a device to a subject is stored with its period of validity, so that each value is linked to the right person.
\end{itemize}

\medskip\noindent\textbf{What a real integration would still require.} No open interface of ISS equipment is assumed. A real integration would require the interface documentation of each device, the permission of the agencies that own the equipment and the data, a review of the security of the station network, an agreement for the handling of crew medical data, and a validation of each adapter against recorded data before use. Until then, every adapter has a simulated counterpart that produces the same canonical observations, labeled \texttt{SIMULATED}, which is how the AHI demonstrator works.


\section{Analytical Methodology}
\label{sec:analytical-methodology}

This section describes what AHI measures and how it turns the measurements into status, predictions and recommendations. The analysis follows the three levels of \ref{sec:solution-overview}: real-time monitoring, predictive analytics and prescriptive intelligence. It starts with explainable statistics and explicit rules, and adds machine learning models only where enough real data exist to train and evaluate them (\ref{sec:artificial-intelligence}).

All numeric parameters in this section, such as band widths, level cutoffs, weights and time windows, are \textbf{configuration}. They are kept in one versioned configuration, they are not official NASA limits or clinical thresholds, and each must be grounded in a cited source or reviewed by a flight surgeon before operational use.

\subsection{Health Dimensions}
\label{subsec:dimensions}

AHI organizes monitoring into ten dimensions. For each one, \ref{tab:dimensions} lists what is measured, what is analyzed and which risks the analysis helps to identify. The sources are the ISS devices of \ref{sec:hardware-sensors}.

\begin{table}[htbp]
\centering
\vspace{0.8em}
\caption{The ten health dimensions of AHI.}
\label{tab:dimensions}
\scriptsize
\begin{tabular}{p{2.3cm}p{4.0cm}p{4.2cm}p{3.5cm}}
\hline
\textbf{Dimension} & \textbf{What is measured} & \textbf{What is analyzed} & \textbf{Risks it helps identify} \\
\hline
Cardiovascular & Heart rate, heart rate variability, blood pressure, ECG and rhythm, oxygen saturation, vascular findings from ultrasound & Persistent heart rate changes, variability below the personal pattern, pressure changes, deconditioning, rhythm changes, exercise response & Cardiovascular adaptation, orthostatic intolerance, circulatory changes \\
Respiratory and oxygenation & Breathing rate and volume, oxygen saturation, breathing during sleep and exercise, CO\textsubscript{2} exposure & Persistent fall in saturation, abnormal rise in breathing rate, link between air quality and physiology & Hypoxia, exposure to elevated CO\textsubscript{2} and contaminants \\
Muscle and bone & Strength, endurance, exercise load and volume, muscle findings from ultrasound, bone results from periodic exams and markers & Progressive loss of strength, lower tolerated load, reduced exercise capacity & Muscle atrophy, bone loss, reduced physical capacity \\
Immune & Blood cell counts, stress hormones, immune markers, signs of latent virus reactivation & Changes of immune markers against the personal baseline, links between stress, sleep and immunity & Altered immunity, higher susceptibility to infection \\
Neurological and sensorimotor & Reaction time, cognition test results, symptoms of space motion sickness, balance and coordination tests & Slower responses, worse attention and memory, coordination problems & Sensorimotor and vestibular changes \\
Vision & Visual acuity and Amsler grid per eye; visual symptoms reported by the crew; retinal photograph, optical coherence tomography and ocular ultrasound findings per eye & Change of each eye against the preflight examination, trend across the scheduled examinations, new symptoms between examinations, co-occurrence with blood pressure and time in flight & Eye changes of SANS, blurred vision \\
Hearing & Hearing thresholds per frequency and ear; personal noise exposure during work and sleep & Change of each ear against the preflight audiogram, trend across tests, noise dose next to hearing results & Hearing threshold shift, interference with communication, alarms and sleep \\
Sleep and circadian rhythm & Sleep duration, efficiency, stages when the headband is used, awakenings, light exposure, regularity & Accumulated sleep debt, circadian shift, fragmented sleep, links with performance and recovery & Sleep loss, fatigue, lower performance \\
Psychological and behavioral & Self reported mood, stress, workload, rest quality and isolation; structured questionnaires & Persistent change of mood or stress, mental fatigue, links between isolation, sleep and performance & Behavioral changes and team problems; indicators only, never a diagnosis \\
Metabolic, renal and nutritional & Urinalysis results, blood chemistry, fluid intake and diet reported by the crew & Dehydration, nutritional imbalance, kidney stone risk factors, links between diet, exercise and recovery & Renal problems, bone loss, inadequate nutrition \\
\hline
\end{tabular}
\vspace{1.2em}
\end{table}

The evidence behind the choice of dimensions is in \ref{sec:scientific-foundation}. Two constraints shape how each one is analyzed. First, the dimensions with frequent data (cardiovascular, respiratory, sleep, activity and exercise) support time series analysis, while the immune, metabolic, vision and hearing dimensions rely on periodic samples and examinations and allow only slow, sparse trends. Second, the psychological dimension uses indicators of risk and a request for appropriate evaluation. AHI never diagnoses depression, anxiety or any other condition from heart rate, messages or expressions.

\subsection{Level 1: Real-Time Monitoring}
\label{subsec:baseline}

\medskip\noindent\textbf{Personal baseline.} Each person is compared with their own history and not with a population average. For every person, metric and activity context (sleep, rest, work or exercise), the engine keeps a robust estimate of the normal value and a band around it. Contexts matter because a heart rate during exercise should be compared with past exercise and not with sleep. For the accepted samples $x_1,\dots,x_n$ of one context:

\begin{equation}
m=\operatorname{median}(x_i),\qquad \mathrm{MAD}=\operatorname{median}\left(\lvert x_i-m\rvert\right),\qquad s=\max\left(1.4826\,\mathrm{MAD},\;s_{\min}\right)
\end{equation}

\begin{equation}
\text{normal band}=\left[m-ks,\;m+ks\right]
\end{equation}

The median and the median absolute deviation (MAD) resist outliers, the factor 1.4826 makes the MAD comparable to a standard deviation for normal data, and the floor $s_{\min}$ keeps a very stable signal from making small changes look large. Preflight data provide the first baseline, which is then updated in flight. A sample enters the baseline only if it is not an outlier and not marked questionable, and new samples are held out while a slow drift is active, so that a gradual decline does not quietly become the new normal. A baseline is used only after enough samples, and until then the interface says that it is still learning the person's normal.

\medskip\noindent\textbf{Acute deviation.} For a value $x$ the signed robust z-score is $z=(x-m)/s$. The acute check uses the median of the last three readings, which filters a single spike:

\begin{equation}
z_{\text{acute}}=\frac{\operatorname{median}\left(x_t,x_{t-1},x_{t-2}\right)-m}{s}
\end{equation}

Increasing values of $z_{\text{acute}}$ in the direction that is bad for the metric give moderate, high and critical levels.

\medskip\noindent\textbf{Slow drift.} An acute check cannot see a gradual change. An exponentially weighted average $D_t$ of the z-scores, with a half-life $H$ that stays correct for irregular samples, follows slow changes:

\begin{equation}
\alpha_t=1-0.5^{\,\Delta t/H},\qquad D_t=D_{t-1}+\alpha_t\left(z_t-D_{t-1}\right)
\end{equation}

A drift that stays above a level for a minimum time raises an alert, so that a single bad hour does not.

\medskip\noindent\textbf{Periodic data.} Sleep, exercise sessions, cognition tests and check-ins arrive once a day or less, and they use runs of consecutive values outside the band. Laboratory, urinalysis and ultrasound results are compared with the person's previous results and with the reference ranges reported with them.

\medskip\noindent\textbf{Environment.} Cabin signals have no personal normal. They are compared with limits taken from cited standards, such as the NASA-STD-3001 limit for the one hour average CO\textsubscript{2} partial pressure \cite{nasa_ochmo_tb004_co2}, and with their own recent trend.

\medskip\noindent\textbf{Noise and hearing.} Personal noise exposure is compared with a cited limit and not with a personal normal: the 24 hour noise dose $D$ measured by the crew-worn dosimeter is compared with the NASA-STD-3001 limit $D\le100$ \cite{nasa_ochmo_tb035_acoustics}, and the work and sleep period levels are kept as separate observations, because noise during sleep also affects rest. Hearing is compared with the person. Each audiogram is compared, frequency by frequency and ear by ear, with the preflight audiogram of the same person, which is the baseline of this dimension:

\begin{equation}
\Delta T_{f,e}=T_{f,e}-T^{\text{pre}}_{f,e}
\end{equation}

where $T_{f,e}$ is the threshold at frequency $f$ for ear $e$, in decibels hearing level. A shift that persists in a later test is flagged for review by the flight surgeon; the size of shift that counts as significant is configuration that the flight surgeon sets, not a value invented by AHI. The platform shows the noise exposure and the hearing results on the same timeline, so that a reviewer can see whether a shift followed a period of high exposure. It reports this as co-occurrence and not as cause, all the more because low frequency shifts on the ISS are not typically caused by noise \cite{coste2024hearing}.

\medskip\noindent\textbf{Eyes.} The vision dimension combines tests of how well the person sees, such as visual acuity and the Amsler grid, with examinations of the structure of the eye, such as OCT, retinal photography and ocular ultrasound, because structural changes can exist without symptoms \cite{brunstetter2023brief}. Each result, for example the retinal and nerve fiber layer thickness measured by OCT, is compared per eye with the preflight examination and with the previous in-flight results. AHI shows the change and its trend across the scheduled examinations and requests review when an eye changes more than the configured amount.

Between examinations, the daily check-in lets the crew member report visual symptoms: blurred vision, difficulty focusing, distorted lines or images, or a change in near or distant vision. A new symptom creates a request to inform the crew surgeon, as NASA asks crew members to do for any vision change \cite{brunstetter2023brief}, and can bring the next screening forward. Since SANS severity is related to mission duration \cite{brunstetter2023brief} and its leading hypothesis is the headward fluid shift \cite{nasa_fluid_shifts_study}, the vision panel places time in flight, blood pressure and fluid related indicators, when available, on the same timeline as the eye results. These relations are shown as co-occurrence. AHI does not interpret images, does not decide whether a finding is SANS and does not diagnose SANS automatically.

\medskip\noindent\textbf{Career radiation dose.} NASA-STD-3001 limits the total career effective dose of an astronaut from spaceflight radiation to less than 600~mSv, for all ages and sexes \cite{nasa_ochmo_tb020_radiation}. This is a published, cited limit, so AHI follows it directly. For each astronaut, the career dose is the sum of the dose from previous missions, taken from the medical record, the dose of the current mission measured by the personal dosimeter, and an estimate for the time since the last readout, computed from the area monitors and labeled as derived:

\begin{equation}
D_{\text{career}}=D_{\text{prior}}+D_{\text{mission}}+\hat{D}_{\text{since readout}},\qquad
u=\frac{D_{\text{career}}}{600~\text{mSv}}
\end{equation}

The platform shows the fraction $u$ of the limit already used and projects it to the end of the mission at the recent dose rate $\dot{D}$, for the mission profile in use (ISS, Moon or Mars):

\begin{equation}
\hat{D}_{\text{end}}=D_{\text{career}}+\dot{D}\,T_{\text{remaining}}
\end{equation}

The projection is a planning aid for the flight surgeon and for mission planning. As stated in \ref{subsec:spaceflight-challenges}, a dosimeter value is compared with the limit only when it is reported as effective dose, and AHI does not estimate cancer or any other radiation health risk.

\medskip\noindent\textbf{Radiation shelter rule.} NASA-STD-3001 requires a method to alert the crew when radiation levels are expected to exceed acceptable levels, and the NASA technical brief notes that the time course of solar particle events usually allows the crew to reconfigure the vehicle and prepare a shelter, for example with water and stowage around the crew quarters \cite{nasa_ochmo_tb020_radiation}. AHI turns this into a rule with two independent triggers:

\begin{enumerate}
\item \textbf{Space weather notification.} A DONKI notification of a solar energetic particle event, or of a flare or coronal mass ejection with a predicted particle event, relayed on board by the ground node (\ref{subsec:space-weather}).
\item \textbf{On board measurement.} A rise of the area dose rate measured by the station radiation monitors well above its recent trend, which needs no link to Earth. The size of the rise is configuration.
\end{enumerate}

Either trigger creates a high priority recommendation for the whole crew: prepare the radiation shelter according to the mission protocol, and review any planned extravehicular activity. It is marked as an on board action, because the crew can act at once, and it is sent to the ground for confirmation. The recommendation shows its trigger, the notification identifier or the measured dose rate, and it is closed when the dose rate returns to its trend and the notification period ends. AHI recommends; the decision to shelter follows the flight rules of the mission.

\medskip\noindent\textbf{Data quality.} A signal with no valid value within its expected cadence becomes \texttt{NO\_DATA}, and a device reported as stale or offline has the same effect. Questionable values cannot escalate beyond a moderate level on their own. \textbf{Missing data is never treated as normal.}

\subsection{Level 2: Predictive Analytics}
\label{subsec:trend-analysis}

Prediction answers two questions: where is a signal going, and how likely is a deterioration over a given horizon. AHI uses several techniques, summarized in \ref{tab:techniques}.

\begin{table}[htbp]
\centering
\vspace{0.8em}
\caption{Analytical techniques of AHI and their use.}
\label{tab:techniques}
\small
\begin{tabular}{p{4.2cm}p{9.8cm}}
\hline
\textbf{Technique} & \textbf{Use} \\
\hline
Statistical baselines & Compare each indicator with the individual history \\
Anomaly detection & Identify unusual changes, univariate and multivariate \\
Time series forecasting & Project the trend of a signal with an uncertainty interval \\
Multivariate analysis & Relate different physiological, behavioral and environmental signals \\
Risk and survival models & Estimate the risk of an event within a time window \\
Explainable AI & Show which factors drove a prediction \\
\hline
\end{tabular}
\vspace{1.2em}
\end{table}

\medskip\noindent\textbf{Trend projection.} The first forecasting method is damped Holt linear smoothing, which is transparent and needs little data:

\begin{equation}
\ell_t=\alpha y_t+(1-\alpha)\left(\ell_{t-1}+\phi b_{t-1}\right),\qquad
b_t=\beta\left(\ell_t-\ell_{t-1}\right)+(1-\beta)\phi b_{t-1}
\end{equation}

\begin{equation}
\hat{y}_{t+h}=\ell_t+\left(\phi+\phi^2+\dots+\phi^h\right)b_t
\end{equation}

The damping $\phi<1$ makes the trend fade with the horizon instead of growing without limit. The interval is $\hat{y}\pm z_{q}\,\sigma_h$, where $\sigma_h$ is the standard deviation of recent residuals widened with the horizon. The \textit{time to band exit}, the first moment the projection leaves the personal band, turns a trend into an operational question. Projections are shown for 24, 48 and 72 hours where the data support them, and the interface says when there is not enough data.

\medskip\noindent\textbf{Combined patterns.} Explicit rules look at several signals together, for example a fatigue pattern made of short or inefficient sleep, low heart rate variability and a rising resting heart rate, or elevated cabin CO\textsubscript{2} while breathing rate rises and reaction time slows. A combined rule reports a co-occurrence and not a cause.

\medskip\noindent\textbf{Transparent risk score.} The rule hits of one astronaut are combined into an operational attention level (low, moderate, high or critical), a score from 0 to 100, the contributing factors in plain language and a confidence value. Each triggered rule contributes a weight $w_i$ that grows with its level, combined as a noisy-OR:

\begin{equation}
\text{score}=100\left(1-\prod_i\left(1-w_i\right)\right)
\end{equation}

Several independent moderate deviations escalate without exceeding 100. A critical level requires at least one critical finding, a lost signal weighs as unknown and not as worse, and self report alone cannot raise the level above moderate. Confidence falls with every missing or questionable key signal and describes how much data stands behind the level. The score is an operational level and not a probability.

\medskip\noindent\textbf{Machine learning models.} When enough real data exist, AHI adds learned models behind the same interface: Isolation Forest for multivariate anomaly detection, gradient boosted trees (XGBoost) and random forests for risk estimation and for ranking the factors that matter, and, once long histories exist, sequence models such as LSTM networks or the Temporal Fusion Transformer for multivariate forecasting. The initial set is the personal baseline, documented rules, Isolation Forest and XGBoost. Every learned model is versioned, documented with its features and training data, evaluated for false alerts and calibration, and explained by the factors that drove each output (\ref{sec:artificial-intelligence}).

\medskip\noindent\textbf{Predictions offered.} \ref{tab:predictions} lists the predictions AHI is designed to produce. They do not all have the same precision. Sleep and fatigue have frequent data and allow temporal models, while bone loss and immune changes depend on periodic clinical and laboratory evaluations, so their predictions are slower and less certain. No prediction is presented as a diagnosis or as the certainty that a disease will occur.

\begin{table}[htbp]
\centering
\vspace{0.8em}
\caption{Predictions AHI is designed to produce.}
\label{tab:predictions}
\footnotesize
\begin{tabular}{p{3.4cm}p{5.4cm}p{5.2cm}}
\hline
\textbf{Prediction} & \textbf{Data used} & \textbf{Expected result} \\
\hline
Accumulated fatigue & Sleep, heart rate variability, activity, heart rate & Risk of fatigue in the next 24 to 72 hours \\
Cardiovascular deterioration & ECG, blood pressure, heart rate variability, exercise & Trend of cardiovascular change \\
Muscle loss & Strength, load, repetitions, training & Trend of loss of performance \\
Sleep changes & Actigraphy, EEG headband, light exposure & Risk of sleep debt \\
Immune changes & Blood, saliva, stress hormones, history & Trend of immune change \\
Dehydration & Fluid intake, urinalysis, clinical indicators & Risk of fluid imbalance \\
Psychological deterioration & Questionnaires, sleep, stress, isolation & Risk of worse well-being \\
Environmental risk & CO\textsubscript{2}, radiation, contaminants, history & Accumulated exposure and possible effect on health \\
Vision change & Vision screening, OCT, retinal photograph and ultrasound findings, symptoms, time in flight & Trend of each eye against the preflight examination; never a SANS diagnosis \\
Hearing change & Audiograms, noise dose, sleep period noise & Trend of thresholds against the preflight audiogram, with exposure alongside \\
Career radiation dose & Personal dosimeter, area monitors, prior missions, mission profile & Fraction of the NASA-STD-3001 career limit used, projected to the end of the mission \\
\hline
\end{tabular}
\vspace{1.2em}
\end{table}

\subsection{ASTRA Deviation Analysis}
\label{subsec:astra}

ASTRA is a statistical analysis of the daily deviation of one person from their own baseline. It is offered by the ground node as a stateless operation (\texttt{POST /api/v1/analytics/deviation}) and its results are derived data. The caller sends the stay of one person so far, with one record per day (the daily value of each metric and the mission phase), and receives, for every in-flight and post-flight day, a level of \texttt{LOW}, \texttt{MODERATE}, \texttt{HIGH} or \texttt{NO\_DATA} together with the detectors that fired.

\medskip\noindent\textbf{How a day is judged.}
\begin{enumerate}
\item \textbf{Personal baseline.} The pre-flight days of the person form the baseline. Each value becomes a $z$ score against the baseline mean and an effective standard deviation that shrinks the personal deviation toward the deviation of a reference population, so a short baseline does not produce an unstable scale. At least seven valid pre-flight readings are needed for each metric. Weight is judged by its rate of change over the last 14 days.
\item \textbf{Expected adaptation.} The $z$ score is standardized by the adaptation expected on that day and by the spread of healthy people in that phase, both read from a reference pack calibrated on NASA bed rest data. Bed rest is an analog of spaceflight, not spaceflight, and this is stated wherever the result is shown.
\item \textbf{Statistical detectors.} Three detectors run on the standardized value: a $z$ score threshold held for two days, an EWMA of the deviation and a CUSUM for small steady drifts, each requiring several consecutive days, and a fourth detector on the rate of weight change. A day without a reading never resets a run and is never read as normal.
\item \textbf{Level.} \texttt{HIGH} requires the weight rate or at least two of the three statistical detectors on the same day. Any other confirmed alarm is \texttt{MODERATE}. No alarm is \texttt{LOW}. A day in which nothing could be evaluated is \texttt{NO\_DATA}.
\end{enumerate}

\medskip\noindent\textbf{Machine learning stays informational.} An Isolation Forest scores the combination of eleven features (the standardized vitals, their three day mean and the weight rate). Its score and flag are returned next to the level and never change it, in line with the principle that statistics come first and machine learning only adds information (\ref{sec:artificial-intelligence}).

\medskip\noindent\textbf{In the interface.} The attention panel shows the latest judged day of the selected astronaut: a level in plain words (within the personal range, moderate deviation, high deviation or not enough data), the main indicator and a note that the analysis is decision support and not a diagnosis. For the simulated crew, the history sent to the analysis is a virtual history calibrated on real NASA bed rest data and labeled simulated. If the ground node is unreachable, the panel simply does not appear: the application never depends on ASTRA.

\medskip\noindent\textbf{Limits.} The limits and weights are the ones calibrated in the analysis, not official limits. The reference pack comes from a bed rest analog with a small number of participants, so the thresholds must be recalibrated with real flight data before operational use (\ref{sec:future-work}).

\subsection{Level 3: Prescriptive Intelligence}
\label{subsec:recommendations}

The third level turns predictions into actions. For each relevant risk, the recommendation engine selects interventions from a versioned catalog, considering the current condition, the history, the evolution of the indicators, the schedule and workload, sleep, exercise, diet, the cabin environment, and medical restrictions and mission protocols. The catalog contains six kinds of action:

\begin{itemize}
\item \textbf{Exercise:} adjust training intensity, resistance or aerobic work and recovery periods. Example: a decline of muscle performance leads to a suggestion to review the training prescription with the specialist.
\item \textbf{Nutrition and hydration:} follow fluid and protein intake and flag possible deficiencies. Example: low fluid intake leads to a suggestion to follow a hydration plan approved by the medical team.
\item \textbf{Sleep and recovery:} regular rest periods, lighting, pauses between activities. Example: three consecutive nights of short sleep lead to a suggestion to review the rest schedule.
\item \textbf{Mental health and well-being:} breathing and relaxation exercises, leisure time, planned contact with family, and a psychological evaluation when needed. Example: persistent signs of stress lead to a suggestion of a relaxation activity and of contact with the psychological support team.
\item \textbf{Medical follow-up:} request a medical evaluation, repeat an examination or refer to a specialist. Example: persistent cardiovascular changes lead to a request for review by the flight surgeon, a report of blurred vision leads to a request to inform the crew surgeon, and a hearing threshold shift leads to a request to repeat the hearing assessment.
\item \textbf{Environmental protection:} radiation shelter, hearing protection and exposure reduction, following mission protocols. Example: a solar particle event leads to a suggestion to prepare the radiation shelter, and a high noise dose during sleep leads to a suggestion to use hearing protection and to check the noise sources near the crew quarters.
\end{itemize}

The catalog never recommends medication, doses or treatment. Each recommendation is a record with seven fields: the problem identified, the evidence that justifies it, the suggested action, the expected result, the time for reassessment, the status of the intervention and the evolution of the indicators after it. The last two fields close the loop: over time the platform learns which interventions are followed by better indicators, always within the protocols validated by the medical team.

\medskip\noindent\textbf{Autonomy and distance.} Each recommendation states who can act. \textit{On board} means that the crew can act now. \textit{Ground review} means that a flight surgeon or mission control should review it. Items that are high, critical or addressed to the flight surgeon always require human review. When the communication delay is long, a high or critical item that needs ground review also creates an on board first step, whose reason states the earliest possible reply from Earth:

\begin{equation}
t_{\text{reply}}=2\,t_{\text{one way}}+t_{\text{review}}
\end{equation}

\medskip\noindent\textbf{Scenarios with and without intervention.} For a selected action, AHI compares the projected course of the indicators with and without the intervention. The comparison is a planning aid. Its effect sizes must come from published evidence or from the observed results of past interventions, and it is labeled as a projection and not as evidence of efficacy.


\section{Artificial Intelligence}
\label{sec:artificial-intelligence}

AHI combines three kinds of artificial intelligence, each with a different job: statistical models that establish what is normal for each person, machine learning models that detect patterns and estimate risk, and a large language model (LLM) that writes the results in plain, direct language. The LLM is accessed through the \textbf{Claude API} from Anthropic \cite{anthropic_claude_api} whenever there is a link to Earth, and through a \textbf{local model specialized in AHI}, running on board, when there is not. The central rule is that \textbf{generative AI explains the results but is never responsible for the clinical calculations}. Status, risk, forecasts and the choice of recommendations come from the analytical engine of \ref{sec:analytical-methodology}, where every output can be traced to data and to a versioned model or rule.

\subsection{Roles of Each Layer}
\label{subsec:ai-roles}

\begin{table}[htbp]
\centering
\vspace{0.8em}
\caption{The layers of artificial intelligence in AHI.}
\label{tab:ai-layers}
\small
\begin{tabular}{p{3.4cm}p{5.0cm}p{5.6cm}}
\hline
\textbf{Layer} & \textbf{Methods} & \textbf{Responsibility} \\
\hline
Statistical & Robust baselines, drift detection, damped trend projection, explicit rules & Personal normal, acute and slow deviations, first forecasts; works from the first day \\
Machine learning & Isolation Forest, XGBoost, random forests; later LSTM and Temporal Fusion Transformer & Multivariate anomalies, risk estimates, factor ranking, multivariate forecasts \\
Language model, connected & Claude API with retrieval of curated sources & Plain language explanations of predictions, alerts and recommendations; summaries and handover text \\
Language model, on board & A local model specialized in the AHI context & The same explanations when the link to Earth is lost \\
\hline
\end{tabular}
\vspace{1.2em}
\end{table}

\subsection{Machine Learning Models}
\label{subsec:ml-models}

\medskip\noindent\textbf{Isolation Forest} isolates observations that are unusual in several signals at once, which complements the univariate checks. It is trained per person or per small group, on data that the baselines already classify as normal.

\medskip\noindent\textbf{XGBoost and random forests} estimate the risk of a defined event within a window, such as a fatigue episode in the next 72 hours, from features such as sleep debt, heart rate variability relative to baseline, training load and environmental exposure. Their feature importance and per prediction contributions identify which factors drove a result.

\medskip\noindent\textbf{Sequence models}, LSTM networks and the Temporal Fusion Transformer, forecast several signals together. They need long histories and are added only when such histories exist.

\medskip\noindent\textbf{The data problem.} Learned models need data, and data on astronauts are scarce: there is no public continuous telemetry, crews are small, and labeled outcomes are rare. AHI therefore follows a staged approach. It starts with baselines and documented rules, which need no training. It adds Isolation Forest and XGBoost when enough data from the mission and from authorized historical sources exist, and sequence models later. Public NASA data, analog studies and simulated data are used to develop and test the models, never as a substitute for validation on the target population (\ref{sec:data-sources}). A model trained only on simulated data learns the simulator and is not used for crew decisions.

\medskip\noindent\textbf{Governance.} Each model has a model card with its purpose, features, training data, version, evaluation results and known limits. Each output records the model version that produced it. A new model version runs in shadow mode, next to the current one and without affecting the crew view, until its evaluation is reviewed.

\subsection{Claude API Explanation Layer}
\label{subsec:llm-layer}

Statistical predictions are precise but hard to read under pressure. A line such as ``projected resting heart rate 71 beats per minute at 48 hours, 80\% interval 66 to 76, band exit in 31 hours, confidence 0.72'' is correct, but it does not tell a crew member what it means. AHI uses the Claude API to rewrite such outputs in a more direct way, for example: ``If the current trend continues, your resting heart rate is likely to rise above your usual range in about a day. The estimate is moderately certain because two nights of sleep data are missing.'' The numbers in the sentence are the numbers of the engine, and nothing is added.

The tasks of the layer are the following:

\begin{itemize}
\item \textbf{Write predictions directly:} turn each forecast, with its horizon, interval and confidence, into one or two plain sentences.
\item \textbf{Explain alerts and recommendations} already selected by the engine: what changed, compared with what, and why it matters.
\item \textbf{Summarize} the condition of a crew member or of the crew for a shift handover.
\item \textbf{Answer questions} about the data of the platform, citing the records and the curated scientific sources it used.
\item \textbf{Adapt to the reader:} plain language for the crew, technical detail for the flight surgeon.
\end{itemize}

The API receives only structured outputs of the engine, such as the metric, the value, the baseline, the forecast, the factors, the confidence and the recommendation, with pseudonymized identifiers and passages from a curated library of NASA and peer reviewed sources. It does not receive raw signals, names or identifying data, and it does not compute status or risk. The model in use is recorded with every generated text. The choice of model follows the current Claude models available through the API, and it is revised as new models are released.

\subsection{Operation Without Contact with Earth}
\label{subsec:offline-llm}

The Claude API is an internet service. It is reachable only while the spacecraft has a link to Earth. Space links are interrupted: during Artemis~II, the Orion spacecraft lost all contact with Earth for about 40 minutes while it passed behind the Moon \cite{ctv_artemis2_blackout}, and the communication delay grows with distance up to about 24 minutes one way for Mars \cite{nasa_space_communications_seven_things}. Interruptions are also possible on the ISS. Health monitoring cannot stop during these periods, and it is exactly when the crew is alone that a clear explanation matters most.

The analytical engine already runs on board and does not depend on the link (\ref{sec:software-architecture}). For the explanation layer, the proposed solution is to \textbf{train a local language model specialized in the context of AHI}, installed on the computers of the ISS or of the Artemis spacecraft and available to the astronauts without any connection. Its design follows four principles:

\begin{itemize}
\item \textbf{Specialized and small.} The local model does not need general knowledge. It needs to explain the outputs of AHI: the metric catalog, the eight health dimensions, the rules, the forecasts, the recommendation catalog and the curated scientific library. An open model small enough for the on board computers is adapted to this context with examples written and reviewed by the team and by flight surgeons.
\item \textbf{Same inputs, same guardrails.} It receives the same structured outputs as the Claude API and its texts go through the same checks (\ref{subsec:guardrails}).
\item \textbf{Automatic switch.} When the link is lost, the platform switches to the local model and labels its texts as generated on board. When the link returns, texts can again be written through the Claude API, and the history is synchronized with the ground.
\item \textbf{Validated before flight.} The local model is evaluated on reference cases before it is installed, for factual consistency, absence of diagnostic wording and correct use of citations, and a new version is installed only after review.
\end{itemize}

If neither model is available, the platform shows the structured result with template text. Monitoring, alerts and recommendations never depend on a language model.

\subsection{Guardrails}
\label{subsec:guardrails}

The same rules apply to the Claude API and to the local model:

\begin{itemize}
\item \textbf{No diagnosis and no prescription.} The model does not name diseases, does not suggest medication or doses, and does not change the level or the recommendation selected by the engine.
\item \textbf{Grounded statements only.} Every number in a generated text must exist in the structured input, and every scientific statement must cite a passage from the curated library. A text that fails these checks is not shown, and the structured result is shown instead.
\item \textbf{No invented limits.} The model never states a threshold or an official limit that is not in its input.
\item \textbf{Uncertainty kept visible.} Confidence, missing data and limitations from the input must appear in the text, and a forecast is never written as a certainty.
\item \textbf{Origin visible.} Generated text is labeled as generated, with the model and the place where it ran, and the structured result is always one click away.
\end{itemize}

\subsection{Open Decisions}
\label{subsec:ai-open-decisions}

Four decisions remain open and need input from mission medical and security staff. The first is the data agreement for the Claude API: which structured, pseudonymized data may leave the spacecraft, and under which retention terms. The second is the choice of the open model for the local layer and of the on board hardware that runs it. The third is how the local model is updated during a long mission, given limited bandwidth. The fourth is the shared evaluation of both language layers, with a set of reference cases that checks factual consistency, the absence of diagnostic wording and the correct use of citations.


\section{Software Architecture}
\label{sec:software-architecture}

This section describes the architecture of AHI in its final scope. Two constraints shape it. First, the crew must keep monitoring, alerts and on board recommendations when the link to Earth is slow or interrupted, so the essential analysis runs on board. Second, the flight surgeon, mission control and researchers work on the ground with the full history, so data are synchronized to a ground system. The architecture is therefore a pair of nodes, one on board and one on the ground, built from the same six modules.

\subsection{Components}
\label{subsec:components}

\begin{table}[htbp]
\centering
\vspace{0.8em}
\caption{The six modules of AHI and where they run.}
\label{tab:components}
\small
\begin{tabular}{p{3.2cm}p{7.4cm}p{3.4cm}}
\hline
\textbf{Module} & \textbf{Responsibility} & \textbf{Runs} \\
\hline
Device integration & One adapter per ISS device. Converts native output to canonical observations, reports device state, accepts late batches & On board \\
Health data platform & Current and historical observations, baselines, model outputs, recommendations and the audit trail, with origin and quality on every record & On board (recent window) and ground (full history) \\
Analytics and prediction engine & Baselines, anomaly detection, combined rules, risk score, forecasts and learned models & On board (essential) and ground (full models and retraining) \\
Recommendation engine & Selection of interventions from the catalog, autonomy, follow-up of their effect & On board and ground \\
Explanation layer & Claude API with retrieval of curated sources and output checks; a local model specialized in AHI when contact is lost (\ref{sec:artificial-intelligence}) & Claude API through the link to Earth; local model on board \\
Dashboards and review & Crew view, crew health overview, predictive and recommendation screens, flight surgeon review workflow & Crew view on board; all screens on the ground \\
\hline
\end{tabular}
\vspace{1.2em}
\end{table}

The analytics engine is a library with no dependency on the interface or on storage, so the same code runs on board, on the ground and in an offline evaluation harness. The on board node keeps a recent window of data sufficient for baselines and short forecasts. The ground node keeps the full history and trains and evaluates the learned models, and approved model versions are sent to the on board node.

\subsection{APIs}
\label{subsec:apis}

The ground node exposes a REST interface under a versioned path, and the on board node exposes the same subset it needs. The main groups of operations are the following:

\begin{itemize}
\item \textbf{Ingestion:} batches of canonical observations, answered with the counts of accepted and rejected records and the reason for each rejection. Ingestion is idempotent, so a batch resent after a link interruption creates no duplicates.
\item \textbf{Health queries:} current status, history, baseline, anomalies, risk, predictions and recommendations for each astronaut, and the cabin environment.
\item \textbf{Actions:} acknowledgment of an alert, logging of an intervention and its result, review decisions of the flight surgeon, and messages between crew and ground.
\item \textbf{Devices and data:} device status, import of public datasets and readiness checks.
\item \textbf{Export:} a FHIR R4 bundle of patient, device and observation resources, with pseudonymized identifiers and the origin of every value.
\end{itemize}

\subsection{Entity-Relationship Model}
\label{subsec:entity-relationship-model}

The central entity is the canonical health observation, one record per measurement, with the fields in \ref{tab:observation}. Vital signs, exercise sessions, cabin readings, laboratory results, examination findings and self report all use the same shape.

\begin{table}[htbp]
\centering
\vspace{0.8em}
\caption{Fields of the canonical health observation.}
\label{tab:observation}
\small
\begin{tabular}{p{3.6cm}p{10.4cm}}
\hline
\textbf{Field} & \textbf{Rule} \\
\hline
Observation identifier & Unique; an idempotent key from the source allows deduplication \\
Subject & A pseudonym such as AST-01, or CABIN for the environment \\
Metric code and unit & One code from the metric catalog; canonical unit after normalization, original value kept \\
Value & A number, or null when missing. Never zero for a missing value \\
Measured and received time & ISO 8601 in UTC; for samples, also the collection and result times \\
Data origin & \texttt{DEVICE}, \texttt{MANUAL}, \texttt{LAB}, \texttt{DERIVED}, \texttt{REAL\_PUBLIC} or \texttt{SIMULATED} \\
Source & Device and adapter version, or dataset accession and file for imported data \\
Quality & Status (valid, questionable, invalid or missing), confidence from 0 to 1, and whether the value is estimated \\
Context & Mission phase and activity context \\
\hline
\end{tabular}
\vspace{1.2em}
\end{table}

Around this record the model has the following entities. An astronaut is assigned devices and is the subject of many observations, and each device produces observations. For each astronaut, metric and context the engine keeps a baseline. An astronaut has anomalies, predictions, alerts, check-ins, messages and one current risk assessment for each dimension. A risk assessment suggests recommendations, each recommendation can lead to an intervention, and each intervention records its follow-up indicators. Every prediction and recommendation records the version of the model or rule that produced it. Audit events record acknowledgments, interventions, review decisions, escalations and messages, with time and actor.

\subsection{Data Flow}
\label{subsec:data-flow}

\begin{enumerate}
\item The Bio-Monitor collects heart rate, breathing and other physiological signals.
\item The Actiwatch Spectrum sends sleep and activity data.
\item The exercise equipment provides performance data for each session, and the medical equipment and environmental monitors provide their measurements.
\item Each adapter identifies the astronaut, the device, the time and the context of every measurement and produces canonical observations.
\item The health data platform consolidates current and historical data and checks quality, units and duplicates.
\item The analytics and prediction engine identifies anomalies, computes risks and estimates future trends.
\item The dashboards present the current situation and the possible scenarios, and the explanation layer describes them in plain language.
\item The recommendation engine suggests interventions and routes the relevant cases to medical review.
\end{enumerate}

The on board node synchronizes with the ground through a store and forward queue. Data and messages wait in the queue during a loss of signal and are sent in order when the link returns, and every message shows the communication delay of the mission profile (ISS, Moon or Mars).

\subsection{Technology Stack}
\label{subsec:technology-stack}

\begin{table}[htbp]
\centering
\vspace{0.8em}
\caption{Technology stack of AHI.}
\label{tab:stack}
\small
\begin{tabular}{p{3.4cm}p{10.6cm}}
\hline
\textbf{Layer} & \textbf{Technology} \\
\hline
Interface & React, TypeScript, Tailwind CSS; three.js for the 3D body; a charting library for time series \\
Server & Python, FastAPI, Pydantic, as a modular monolith with background jobs \\
Storage & PostgreSQL with TimescaleDB for time series; Redis for queues and cache \\
Analytics & pandas, NumPy, SciPy, scikit-learn, XGBoost; a deep learning library for sequence models \\
Language model & Claude API with retrieval over a curated library of NASA and peer reviewed sources; a local model specialized in AHI on board for periods without contact \\
Interoperability & FHIR R4 export; LOINC codes where they exist \\
Delivery & Containers; continuous integration with lint, type checks, tests and build \\
\hline
\end{tabular}
\vspace{1.2em}
\end{table}

\medskip\noindent\textbf{Architectural decision.} A modular monolith with asynchronous processing was chosen over microservices. It is simpler to deploy on limited on board computers, easier to validate, and keeps clear module boundaries so that a module can be separated later if needed.

\subsection{Security and Access}
\label{subsec:security}

Crew health data are among the most sensitive data of a mission. Access is role based, with roles for crew member, crew medical officer, flight surgeon, mission control, researcher and data engineer, and each role sees only what it needs. Data are encrypted in transit and at rest, identifiers are pseudonymized outside the medical roles, and the audit trail is persistent and tamper evident. Researchers receive de-identified extracts. Data imported from outside are treated as untrusted: they are hashed, parsed as data and never executed.


\section{Visual Identity, 3D Model and UX}
\label{sec:visual-identity-ux}

The visual language of AHI is a dark mission control look: deep navy backgrounds, translucent glass panels, a teal accent and five status colors that always come with words and an icon. This section defines the design system of the platform, the frontend technologies that implement it and the 3D anatomical model at the center of the interface. The interface is in English, because the work is presented to an international audience, and it states plain language first and the technical term second.

Three principles guided the design. \textbf{Plain language first}, so that a person under stress reads "Act soon" before "Warning". \textbf{Honesty about data}, so that the origin of the data and the loss of a signal are always visible. And \textbf{one source of truth with two views}, so that the astronaut and the flight surgeon see the same state in layouts made for their needs.

\subsection{Frontend Technologies}
\label{subsec:frontend-technologies}

The interface is a single page web application written in TypeScript. \ref{tab:frontend-stack} lists the main tools and the role of each one. The versions are those declared in the project at the time of writing.

\begin{table}[htbp]
\centering
\vspace{0.8em}
\caption{Main frontend technologies of AHI.}
\label{tab:frontend-stack}
\footnotesize
\begin{tabular}{p{3.6cm}p{1.6cm}p{8.8cm}}
\hline
\textbf{Tool} & \textbf{Version} & \textbf{Role} \\
\hline
React & 19 & Component model of the whole interface \\
TypeScript & 7 & Typed code and the shared data contract between the engine and the interface \\
Vite & 8 & Development server and production build, with code splitting so the 3D viewer loads on demand \\
Tailwind CSS & 4 & Design tokens and styling of every screen \\
zustand & 5 & Application state store that owns the analytical engine and publishes only the parts of the state that changed \\
Recharts & 3 & Trend charts with the personal band, gaps for missing data and forecast cones \\
lucide-react & 1 & Icon set used with every status word \\
three.js & 0.187 & WebGL rendering of the 3D body \\
React Three Fiber & 9 & React renderer for three.js, so the 3D scene follows the same component model \\
drei & 10 & three.js helpers: orbit controls, procedural lighting environment, contact shadows, bounding volume hierarchies for picking \\
postprocessing & 6 & Bloom on status highlights and filmic tone mapping \\
glTF Transform and meshoptimizer & 4.5 and 1.3 & Build time optimization and compression of the 3D model files \\
Blender & 4.5 LTS & Headless preparation of the anatomical source model at build time \\
Vitest and oxlint & 5 and 1 & Unit tests and static checks \\
\hline
\end{tabular}
\vspace{1.2em}
\end{table}

The analytical engine is written in the same language and runs in the same application, so the interface reads its results directly from the store. Charts and the 3D scene read live values without forcing the rest of the interface to render again, which keeps the screen responsive while the data update every simulated minute.

% AHI design tokens (frontend/web/src/index.css and src/lib/status.ts)
\definecolor{ahiSpace950}{HTML}{04070D}
\definecolor{ahiSpace900}{HTML}{070C16}
\definecolor{ahiSpace850}{HTML}{0A1120}
\definecolor{ahiSpace800}{HTML}{0E1729}
\definecolor{ahiSpace700}{HTML}{16223A}
\definecolor{ahiSpace600}{HTML}{22324F}
\definecolor{ahiInk100}{HTML}{E8EEFC}
\definecolor{ahiInk300}{HTML}{A9B6CF}
\definecolor{ahiInk500}{HTML}{7D8BA6}
\definecolor{ahiAccent}{HTML}{5EEAD4}
\definecolor{ahiAccentStrong}{HTML}{2DD4BF}
\definecolor{ahiSky}{HTML}{7DD3FC}
\definecolor{ahiNominal}{HTML}{34D399}
\definecolor{ahiWatch}{HTML}{FBBF24}
\definecolor{ahiWarning}{HTML}{FB923C}
\definecolor{ahiCritical}{HTML}{F43F5E}
\definecolor{ahiNoData}{HTML}{94A3B8}
\definecolor{ahiViolet}{HTML}{A78BFA}
\definecolor{ahiVioletText}{HTML}{DDD6FE}
\newcommand{\swatch}[1]{{\setlength{\fboxsep}{0pt}\setlength{\fboxrule}{0.3pt}\fcolorbox{black!35}{#1}{\rule{0pt}{0.85em}\rule{2.2em}{0pt}}}}
\newcommand{\statuspill}[2]{\tikz[baseline=-0.6ex]\node[fill=#1,rounded corners=5pt,inner xsep=7pt,inner ysep=3.5pt,text=ahiSpace900,font=\footnotesize\bfseries]{#2};}

\subsection{Logo}
\label{subsec:logo}

The logo of AHI, shown in \ref{fig:ahi-logo}, is a mark drawn in vector form and a wordmark. The mark combines a circle, a tilted orbit, a small point on that orbit and a heartbeat trace across the center: a body in orbit and its vital signs, which is what the platform watches. Its strokes use a gradient from the teal accent to a light sky blue, and it sits on the deep navy background of the interface. In the application the mark appears in the top bar next to the wordmark ``AHI'', and the documentation portal uses the same mark.

\hwfigure{0.80}{images/branding/ahi-logo.png}{AHI logo: the mark (circle, orbit, orbiting point and heartbeat trace) with the wordmark on the interface background. The full name below the wordmark was added for this report.}{fig:ahi-logo}{Rendered by the authors from the AHI logo source (\texttt{frontend/web/src/features/topbar/TopBar.tsx} and \texttt{docs-site/public/logo.svg}).}

\subsection{Color Palette}
\label{subsec:color-palette}

Status is the only use of the five signal colors, never decoration. \ref{tab:status-language} defines the status language, which is declared once and used by every screen and by the 3D body. Color is never the only cue: every status pairs its color with an icon and a word, and the missing data status changes shape as well.

\begin{table}[htbp]
\centering
\vspace{0.8em}
\caption{Status language: wording, technical label, color and shape.}
\label{tab:status-language}
\small
\begin{tabular}{p{2.2cm}p{2.4cm}p{2.4cm}p{2.8cm}p{3.8cm}}
\hline
\textbf{Status} & \textbf{Crew wording} & \textbf{Technical label} & \textbf{Color} & \textbf{Dot and animation} \\
\hline
Low & All good & Nominal & \swatch{ahiNominal}\ \texttt{\#34D399} & Filled \\
Moderate & Keep an eye & Watch & \swatch{ahiWatch}\ \texttt{\#FBBF24} & Filled \\
High & Act soon & Warning & \swatch{ahiWarning}\ \texttt{\#FB923C} & Filled, slow pulse \\
Critical & Act now & Critical & \swatch{ahiCritical}\ \texttt{\#F43F5E} & Filled, fast pulse \\
No data & Signal lost & No data & \swatch{ahiNoData}\ \texttt{\#94A3B8} & Hollow and dashed, never filled \\
\hline
\end{tabular}
\vspace{1.2em}
\end{table}

The no data status is deliberately different from the low status, so that \textbf{missing data never reads as normal}. In the 3D body the nominal regions use calm, cool tints, so that any warm color reads as attention, and a region with no data is grey, dimmed and still.

Backgrounds use a scale of deep navy tones from near black to a lighter slate. Text uses three levels of light grey for primary, secondary and muted content. A teal accent marks focus rings, links and the spotlight of the guided demonstration. Data origin badges follow a consistent scheme, and whenever simulated data are shown, for development, training or demonstration, a violet badge marked simulated is always visible in the top bar.

\ref{fig:status-preview} shows the five status colors as they appear on the interface background, and \ref{tab:palette} lists the full palette, with the token name used in the code, the color and its use.

\par\bigskip\noindent
\begin{minipage}{\linewidth}
\centering
\begin{tikzpicture}
\node[fill=ahiSpace900,rounded corners=6pt,minimum width=0.96\linewidth,minimum height=3.1cm] (bg) {};
\node[anchor=north west,text=ahiInk300,font=\footnotesize] at ([xshift=10pt,yshift=-8pt]bg.north west) {Status language on the interface background (space-900)};
\node at ([yshift=4pt]bg.center) {\statuspill{ahiNominal}{All good}\hspace{6pt}\statuspill{ahiWatch}{Keep an eye}\hspace{6pt}\statuspill{ahiWarning}{Act soon}\hspace{6pt}\statuspill{ahiCritical}{Act now}\hspace{6pt}\tikz[baseline=-0.6ex]\node[draw=ahiNoData,dashed,line width=1pt,rounded corners=5pt,inner xsep=7pt,inner ysep=3.5pt,text=ahiNoData,font=\footnotesize\bfseries]{Signal lost};};
\node[anchor=south west,font=\footnotesize] at ([xshift=10pt,yshift=8pt]bg.south west) {\textcolor{ahiInk100}{Primary text}\hspace{10pt}\textcolor{ahiInk300}{Secondary text}\hspace{10pt}\textcolor{ahiInk500}{Muted text}\hspace{10pt}\textcolor{ahiAccent}{Accent and focus}\hspace{10pt}\tikz[baseline=-0.6ex]\node[fill=ahiViolet!22!ahiSpace900,draw=ahiViolet!55!ahiSpace900,rounded corners=3pt,inner xsep=4pt,inner ysep=2pt,text=ahiVioletText,font=\scriptsize\bfseries]{SIMULATED};};
\end{tikzpicture}
\captionof{figure}{Status pills, text levels, accent and the simulated data badge drawn with the real colors of the interface.}
\label{fig:status-preview}
\figsource{Drawn by the authors from the AHI design tokens (\texttt{frontend/web/src/index.css} and \texttt{frontend/web/src/lib/status.ts}).}
\end{minipage}
\par\bigskip

\begin{table}[htbp]
\centering
\vspace{0.8em}
\caption{AHI color palette, from the design tokens of the interface.}
\label{tab:palette}
\footnotesize
\begin{tabular}{p{2.6cm}p{2.9cm}p{8.5cm}}
\hline
\textbf{Token} & \textbf{Color} & \textbf{Use} \\
\hline
\multicolumn{3}{l}{\textit{Backgrounds}} \\
\texttt{space-950} & \swatch{ahiSpace950}\ \texttt{\#04070D} & Deepest background, behind the 3D scene \\
\texttt{space-900} & \swatch{ahiSpace900}\ \texttt{\#070C16} & Main page background \\
\texttt{space-850} & \swatch{ahiSpace850}\ \texttt{\#0A1120} & Panels \\
\texttt{space-800} & \swatch{ahiSpace800}\ \texttt{\#0E1729} & Raised panels and cards \\
\texttt{space-700} & \swatch{ahiSpace700}\ \texttt{\#16223A} & Borders and hover surfaces \\
\texttt{space-600} & \swatch{ahiSpace600}\ \texttt{\#22324F} & Dividers and inactive controls \\
\multicolumn{3}{l}{\textit{Text}} \\
\texttt{ink-100} & \swatch{ahiInk100}\ \texttt{\#E8EEFC} & Primary text and values \\
\texttt{ink-300} & \swatch{ahiInk300}\ \texttt{\#A9B6CF} & Secondary text \\
\texttt{ink-500} & \swatch{ahiInk500}\ \texttt{\#7D8BA6} & Muted text, axis labels \\
\multicolumn{3}{l}{\textit{Brand and accent}} \\
\texttt{accent} & \swatch{ahiAccent}\ \texttt{\#5EEAD4} & Focus rings, links, demonstration spotlight, start of the logo gradient \\
\texttt{accent-strong} & \swatch{ahiAccentStrong}\ \texttt{\#2DD4BF} & Pressed and active states \\
Logo gradient end & \swatch{ahiSky}\ \texttt{\#7DD3FC} & End of the logo gradient \\
\multicolumn{3}{l}{\textit{Status (only for status, never decoration)}} \\
\texttt{nominal} & \swatch{ahiNominal}\ \texttt{\#34D399} & Low: all good \\
\texttt{watch} & \swatch{ahiWatch}\ \texttt{\#FBBF24} & Moderate: keep an eye \\
\texttt{warning} & \swatch{ahiWarning}\ \texttt{\#FB923C} & High: act soon \\
\texttt{critical} & \swatch{ahiCritical}\ \texttt{\#F43F5E} & Critical: act now \\
\texttt{nodata} & \swatch{ahiNoData}\ \texttt{\#94A3B8} & No data: signal lost \\
\multicolumn{3}{l}{\textit{Data origin}} \\
Simulated badge & \swatch{ahiViolet}\ \texttt{\#A78BFA} & Violet badge for simulated data (Tailwind violet; text in \texttt{\#DDD6FE}) \\
\hline
\end{tabular}
\vspace{1.2em}
\end{table}


\subsection{Typography}
\label{subsec:typography}

The interface uses the Inter family for text and JetBrains Mono for numbers and codes. Numbers use tabular figures, so that live values do not jitter as digits change. In presentation mode no text is smaller than 13 pixels, chart ticks included. The fonts are bundled with the application, so that the on board interface makes no external request.

\subsection{Components}
\label{subsec:ui-components}

The interface is assembled from a small set of shared components:

\begin{itemize}
\item \textbf{Status pill, icon and dot:} the status language of \ref{tab:status-language}.
\item \textbf{Metric value:} numbers that change smoothly, and that snap when reduced motion is requested or the page is hidden.
\item \textbf{Trend chart:} the personal band labeled "Your normal", gaps wherever data is missing, the forecast cone and markers for alerts, actions and check-ins. Gaps stay gaps: missing data is never filled or smoothed over.
\item \textbf{Sparkline:} six hours of history in crew cards and in the vitals strip.
\item \textbf{Risk gauge:} the score from 0 to 100 with its level.
\item \textbf{Baseline explanation:} a "Why is this my normal?" control that shows the window, the readings, the median and the MAD.
\item \textbf{Autonomy badge:} "Can act on board now", "Needs flight surgeon" or the minimum wait for an Earth reply.
\item \textbf{Origin badge:} the data origin of every value.
\end{itemize}

Several decisions follow from these components. The Crew view shows at most two actions, to reduce the load, and the rest is available in the Ground view. Every recommendation says who can act, because communication delay makes autonomy part of the answer. The critical banner floats over the header of the central area, so that the 3D body keeps its size during an alert.

Accessibility is part of the design: global keyboard shortcuts, keyboard navigation of timeline markers, a visible focus ring, a link that skips to what needs attention as the first focusable element, and accessible labels on every control. The target is conformance with WCAG 2.1 level AA, verified by a formal audit and by tests with a screen reader. The ground screens are designed for large displays, and the Crew view also for the tablets and laptops used on board.

\subsection{3D Anatomical Model}
\label{subsec:3d-model}

The body in the center of the interface is a \textbf{real anatomical atlas}: the open Z-Anatomy model, ``the libre 3D atlas of anatomy'', itself based on BodyParts3D \cite{z_anatomy}. It was chosen because it covers the whole body with named structures organized by anatomical system, which is what layers, selection and the link to health data need, and because its Blender source can be processed automatically and reproducibly. Z-Anatomy is distributed under the Creative Commons Attribution-ShareAlike 4.0 license, so the derived files are shared under the same license with its attribution, shown in a caption under the body that is always visible. The atlas is one adult reference body shown for every crew member. It is not a scan of an astronaut, and its colors come from the health data, never from the model.

\medskip\noindent\textbf{Asset pipeline.} The published files are generated by scripts and never edited by hand. One command runs the whole chain in about five minutes:

\begin{enumerate}
\item \textbf{Source.} The Z-Anatomy archive is downloaded from one allowed host and checked against a pinned SHA-256 hash.
\item \textbf{Blender export.} Blender 4.5 LTS runs headless, with embedded scripts disabled. It keeps only physical bodies, converts vessel and nerve curves into tubes, joins about 2,230 source objects into 250 structures, removes about four fifths of the triangles, moves the body to the viewer frame and writes one file per anatomical system.
\item \textbf{Kidneys.} The upstream kidney model has a non commercial license, so AHI builds its own kidneys in code and fits them to the size and position of the atlas kidneys.
\item \textbf{Optimization.} glTF Transform and meshoptimizer remove duplicates, quantize positions and normals and compress every file.
\item \textbf{Validation.} The published files are checked against a manifest contract that lists every structure with its system, tissue, side, bounding box and origin. Automated tests fail if a required structure is missing or an excluded object appears.
\end{enumerate}

Some upstream content was left out on purpose: the inner ear and the kidney, because of non commercial licenses, reproductive structures, so that the model is a neutral reference body, and helper objects such as labels. \ref{tab:3d-systems} summarizes the result.

\begin{table}[htbp]
\centering
\vspace{0.8em}
\caption{Files of the 3D anatomical model, one per anatomical system.}
\label{tab:3d-systems}
\footnotesize
\begin{tabular}{p{2.6cm}p{6.6cm}p{1.6cm}p{1.6cm}p{1.2cm}}
\hline
\textbf{System} & \textbf{Content} & \textbf{Structures} & \textbf{Triangles} & \textbf{MB} \\
\hline
Integumentary & Skin, the whole body surface & 1 & 48,503 & 0.24 \\
Skeletal & Skull, vertebrae, ribs, sternum and limb bones & 41 & 184,971 & 0.95 \\
Muscular & Large muscles one by one, small ones by group & 133 & 299,818 & 1.66 \\
Cardiovascular & Heart, coronary vessels, great vessels, arteries and veins & 29 & 89,893 & 0.50 \\
Respiratory & Lungs, trachea, bronchi, larynx, diaphragm & 6 & 44,991 & 0.22 \\
Nervous & Brain, cerebellum, brainstem, spinal cord, nerves & 17 & 133,356 & 0.66 \\
Digestive & Tongue to intestines, liver, gallbladder, pancreas & 10 & 41,990 & 0.22 \\
Urinary & Kidneys (AHI geometry), ureters, bladder & 4 & 10,826 & 0.06 \\
Other & Eyes, glands, thymus, spleen & 9 & 26,008 & 0.13 \\
\hline
\textbf{Total} & 9 files & \textbf{250} & \textbf{880,356} & \textbf{4.64} \\
\hline
\end{tabular}
\vspace{1.2em}
\end{table}

\medskip\noindent\textbf{Rendering.} The viewer, built with three.js and React Three Fiber, loads the systems on demand: the muscles, the largest file, are downloaded only when they are needed. Each tissue class has its own physically based material, such as translucent skin, matte bone, soft muscle and red, blue and yellow tubes for arteries, veins and nerves. A lighting environment built in code, filmic tone mapping and bloom only on status highlights complete the look. Every file comes from the application's own server, with no external decoder or image, so the viewer works without internet access, which matters on board.

\ref{fig:atlas-layers} shows the atlas with four layer presets, and \ref{fig:atlas-states} shows how health status is drawn on it.

\hwfigure{1.0}{images/anatomy/atlas-layers.jpg}{The Z-Anatomy atlas in the AHI viewer, front view: (a) default layers, with translucent skin, skeleton, vessels and the monitored organs; (b) skeleton; (c) muscles; (d) organs and vessels without the skeleton.}{fig:atlas-layers}{AHI demonstrator, anatomy lab page (\texttt{frontend/web/anatomy-lab.html}), captured by the authors. 3D anatomy: Z-Anatomy atlas, CC BY-SA 4.0.}

\medskip\noindent\textbf{Interaction.} The user can orbit, zoom and switch between front, back, left and right views. Pointing at a structure shows its name, system and the status of its region; selecting it flies the camera to it, turns the rest of the body into a see-through ghost and opens its detail panel. Layers switch anatomical systems on and off, and the skin opacity can be adjusted.

\medskip\noindent\textbf{Link to health data.} One module maps each structure to the body region and the metrics it reports under. A region in alert gets a pulsing rim in its status color, a region with no data turns grey and still, and structures that are only context, such as the digestive organs or the eyes, never show a status. The heart beats at the live heart rate and the lungs breathe at the live breathing rate of the selected astronaut, and both stop when their data are missing, so that missing data is never drawn as normal.

\hwfigure{0.80}{images/anatomy/atlas-states.jpg}{Health status drawn on the atlas: (a) acute cardiac event, with the heart critical, the lungs in warning and the brain to watch, each labeled with its status word and value; (b) sensor failure, where every region turns grey and still and is labeled no data, never normal. Region statuses set in the lab page.}{fig:atlas-states}{AHI demonstrator, anatomy lab page (\texttt{frontend/web/anatomy-lab.html}), captured by the authors. 3D anatomy: Z-Anatomy atlas, CC BY-SA 4.0.}

\medskip\noindent\textbf{Fallback.} If the atlas cannot load, the viewer switches to a simpler model generated in code, with the same status language and live motion, so the dashboard always has a body to show. If that also fails, the rest of the interface keeps working.

\medskip\noindent\textbf{Limits.} The model is one adult reference body at one scale, derived from a male source with a neutral groin. It does not adapt to sex, height or body composition, and it has no textures.

\subsection{Screens}
\label{subsec:screens}

The platform has four main screens for the ground team and one view for the astronaut, all over the same state.

\begin{enumerate}
\item \textbf{Astronaut health overview:} general health status, vital signs, active alerts, recent history, the level of risk in each of the eight dimensions and the time of the last update of each sensor.
\item \textbf{Predictive health analytics:} historical charts, physiological trends, forecasts for 24, 48 and 72 hours, the scenario without intervention and the scenario with intervention, and the confidence and uncertainty of every forecast.
\item \textbf{Health recommendations:} identified risks, recommendations with their justification, pending actions and the evolution of the indicators after each intervention, with the review workflow of the flight surgeon.
\item \textbf{Crew health monitoring:} the health of the whole crew, comparisons with each person's own normal, priority alerts, the history of each astronaut and the overall health situation of the mission.
\item \textbf{Crew view:} for the astronaut, a large status word, one sentence, at most two actions that can be taken on board, the daily check-in and the person's own values against their own normal.
\end{enumerate}

The ground screens follow a deliberate reading order: who (the crew), what (the body and the headline), why and what to do (the attention panel), when (the status timeline) and the live values. \ref{fig:ground-view} and \ref{fig:crew-view} show how this design was rendered in the AHI demonstrator, with simulated data.

Further screens are an alerts drawer with the history of events, an Earth link panel that shows the delay and the messages, a daily check-in dialog, a mission overview with dose, exercise and sleep summaries, a shift handover report with a FHIR export, and a data and model panel with the devices, the model cards and the data origins. A fixed notice on every screen states that the tool is decision support.

\hwfigure{0.96}{images/ahi-ground-view.png}{Crew health monitoring in the AHI demonstrator, nominal day: crew, 3D body, attention panel, status timeline, live vitals, cabin and devices. Simulated data.}{fig:ground-view}{AHI demonstrator, screenshot by the authors with simulated data, published in the project documentation portal (\texttt{docs-site/public/screenshots/ground-view.png}).}

\hwfigure{0.96}{images/ahi-crew-view.png}{Crew view in the AHI demonstrator: a status word, the top actions with autonomy badges, the daily check-in and values against the personal normal. Simulated data.}{fig:crew-view}{AHI demonstrator, screenshot by the authors with simulated data, published in the project documentation portal (\texttt{docs-site/public/screenshots/crew-view.png}).}


\section{Global Access and Inclusion}
\label{sec:global-access}

Space exploration is international, and so are its crews. The ISS is operated by agencies from several countries, crew members and flight surgeons do not share one first language, and future missions will include people with different backgrounds and abilities. A health platform is only useful if every one of them can understand it quickly, including under stress, without medical or statistical training and without depending on sight. AHI treats this as a requirement of the product, not as an extra, and this section specifies the three features that carry it:

\begin{enumerate}
\item \textbf{Plain language explanations and a health assistant} that turn statistical results into short, direct sentences for lay readers, by text and by voice (\ref{subsec:lay-explanations} and \ref{subsec:assistant}).
\item \textbf{Internationalization} of the whole platform into English, Portuguese, Spanish, Japanese and German (\ref{subsec:languages}).
\item \textbf{Audio description and voice operation}, so that blind and low vision users, people with dyslexia and astronauts busy with a task can follow and operate the platform without looking at a screen (\ref{subsec:accessibility}).
\end{enumerate}

\subsection{Principles}
\label{subsec:global-principles}

Five principles apply to all three features, and \ref{tab:inclusion-users} relates the groups of users to their needs and to the features that serve them.

\begin{itemize}
\item \textbf{One truth, many forms.} The analytical engine produces one structured result. Text, translation, speech and sound are different ways of presenting the same result, and none of them may change a value, a level or a recommendation.
\item \textbf{Plain language first, technical term second.} Every user sees the plain version first and can open the technical version, with the numbers, at any time.
\item \textbf{More than one channel.} Every piece of information has at least two forms: a word and an icon, a sound and a text, a spoken sentence and a written one. No information depends on color, sight or hearing alone.
\item \textbf{Works without internet.} The internet is used when available, as in the demonstrator, which simulates the AI with the Claude API. Without it, the contingency plan uses translations, speech models and a language model installed on board (\ref{subsec:offline-llm}).
\item \textbf{Grounded in published guidance.} The writing rules, the words for uncertainty and the accessibility features are based on published standards and studies: WCAG 2.1 for accessibility \cite{w3c_wcag21}, plain language guidance for the wording \cite{plain_language_guide}, and research on how verbal expressions of uncertainty are understood \cite{budescu2009uncertainty,mastrandrea2010ipcc}. Generated texts are checked automatically against their inputs, which is part of the design. This project does not run user studies or reviews by native speakers: tests with blind, low vision and dyslexic users and the review of the translations are requirements before operational use (\ref{sec:future-work}).
\end{itemize}

\begin{table}[htbp]
\centering
\vspace{0.8em}
\caption{Groups of users, their needs and the AHI features that serve them.}
\label{tab:inclusion-users}
\footnotesize
\begin{tabular}{p{3.4cm}p{5.0cm}p{5.6cm}}
\hline
\textbf{Group} & \textbf{Need} & \textbf{Features} \\
\hline
Crew member without medical training & Understand their own data and what to do & Plain language explanations, health assistant, ``Explain this'' on every value \\
Speaker of another language & Use the platform in their own language & Five languages, reviewed medical glossary, assistant and speech in the same language \\
Blind user & Operate the platform without sight & Audio description of every screen, keyboard and voice navigation, spoken alerts \\
User with low vision & Read the screen comfortably & Large text, high contrast, zoom, spoken description on demand \\
User with dyslexia & Read without effort & Read aloud with highlighting, wider text spacing, short sentences \\
Astronaut busy with a task & Stay informed without looking & Hands-free mode, short spoken summaries, sounds per status level, spoken questions \\
User with reduced hearing & Not miss spoken information & Every sound and spoken message also written; adjustable volume and speed \\
Flight surgeon & Precise, technical detail & Technical version of every explanation, full numbers and sources \\
\hline
\end{tabular}
\vspace{1.2em}
\end{table}

\subsection{Explaining Statistics to Lay Readers}
\label{subsec:lay-explanations}

The engine of AHI speaks the language of statistics: robust z-scores, bands of median plus or minus a multiple of the MAD, drift indices, forecast intervals, confidence values and a noisy-OR risk score (\ref{sec:analytical-methodology}). These outputs are precise, but most crew members are not statisticians, and a value such as $z=-2.4$ does not tell anyone what to do. The language model of AHI (\ref{subsec:llm-layer}) translates these outputs into plain language. It does not compute anything: it receives the structured result of the engine and writes it in words.

\medskip\noindent\textbf{What needs to be translated.} \ref{tab:lay-concepts} lists the statistical concepts that appear in AHI and the plain wording the language model is instructed to use for each one.

\begin{table}[htbp]
\centering
\vspace{0.8em}
\caption{Statistical concepts of AHI and how they are explained to lay readers. Example sentences use simulated values.}
\label{tab:lay-concepts}
\footnotesize
\begin{tabular}{p{3.0cm}p{4.4cm}p{6.6cm}}
\hline
\textbf{Concept} & \textbf{What the engine outputs} & \textbf{Plain explanation} \\
\hline
Personal baseline & Median and MAD band per context & ``Your normal range for resting heart rate is 58 to 66, based on your last two weeks at rest.'' \\
Deviation & Robust z-score & ``Today's value is clearly lower than usual for you.'' The number is shown only in the technical version. \\
Slow drift & Exponentially weighted drift index & ``Over the last five days your heart rate variability has slowly gone down.'' \\
Forecast & Damped trend projection with interval & ``If nothing changes, it will probably leave your normal range in about a day.'' \\
Uncertainty & Interval width and confidence & ``This estimate is fairly uncertain, because two nights of sleep data are missing.'' \\
Risk level & Noisy-OR score from 0 to 100 & ``Act soon. Three signs point the same way: short sleep, lower heart rate variability and a higher resting heart rate.'' The score is never described as a probability. \\
Co-occurrence & Combined rule & ``Cabin CO\textsubscript{2} rose at the same time as your breathing rate. This does not show that one caused the other.'' \\
Missing data & No data status & ``There is no heart rate data from the shirt for 40 minutes, so this value is unknown, not normal.'' \\
Cited limit & Value against a standard & ``You have used 31\% of the career radiation limit set by the NASA standard.'' \\
\hline
\end{tabular}
\vspace{1.2em}
\end{table}

\medskip\noindent\textbf{Writing rules.} The instructions given to the language model, and the checks applied to its output, follow a fixed set of rules:

\begin{itemize}
\item \textbf{Start with what it means.} The first sentence says the status and what to do; details come after.
\item \textbf{Compare with the person, not with a population.} ``Lower than usual for you'', never ``abnormal''.
\item \textbf{Few, rounded numbers.} At most two numbers per sentence, rounded to what the reader needs, with units written in words when read aloud.
\item \textbf{Consistent words for uncertainty.} Confidence is mapped to a small fixed set of words, such as ``likely'', ``fairly uncertain'' and ``very uncertain''. The mapping is configuration, identical in every language, so the same confidence always produces the same word, as in the calibrated uncertainty language of scientific assessments \cite{mastrandrea2010ipcc}. Because readers interpret such words in different ways, a study of that language recommends pairing the words with numeric ranges \cite{budescu2009uncertainty}, so the numeric version is always one tap away. The mapping must be reviewed by flight surgeons before operational use.
\item \textbf{Say what is missing.} Missing or questionable data that lowered the confidence are always named.
\item \textbf{No diagnosis, no cause.} No disease names, no medication, and co-occurrence is never written as cause.
\item \textbf{Short and calm.} Short sentences, everyday words and a calm tone, with no alarming adjectives. Urgency is carried by the status word, not by the wording.
\end{itemize}

\medskip\noindent\textbf{How an explanation is produced.} For every explanation the platform follows the same pipeline:

\begin{enumerate}
\item The engine produces the structured result: metric, value, baseline band, deviation, forecast, interval, confidence, factors, missing data and recommendation.
\item The platform adds the reader profile (crew or flight surgeon), the language of the user and, if needed, passages from the curated library.
\item The language model writes the explanation following the writing rules.
\item Automatic checks verify that every number in the text exists in the input, that the status word matches the level, that the uncertainty word matches the confidence mapping, that no forbidden term appears and that sentences stay short.
\item If every check passes, the text is shown with a ``Show the numbers'' control that opens the technical version. If a check fails, a template sentence built directly from the structured result is shown instead.
\end{enumerate}

\medskip\noindent\textbf{Two readers, one result.} The same result produces two versions. The crew version is the plain text above. The flight surgeon version keeps the plain first sentence and adds the band, the z-score, the interval, the confidence value and the factors, in technical terms. Both versions come from the same input, so they can never disagree.

\medskip\noindent\textbf{``Explain this'' everywhere.} Every value, chart, alert and recommendation has an ``Explain this'' control. It opens the plain explanation of that item, in writing and, if the user wishes, read aloud. It is the shortest path from a number to its meaning, and it is the entry point of the health assistant.

\subsection{Health Assistant}
\label{subsec:assistant}

The health assistant is a chat that answers questions about the health of the astronaut in easy, direct language for lay readers. It is meant for the crew member who wants to understand their own data without medical training, and for any user who prefers to ask instead of searching the screens. It can be used by typing or by speaking, and it answers in writing and aloud. Typical questions are:

\begin{itemize}
\item ``Why is my status \textit{Keep an eye} today?''
\item ``Is my sleep getting better or worse this week?''
\item ``What does the radiation alert mean for me?''
\item ``What can I do before my next exercise session?''
\item ``What does \textit{fairly uncertain} mean in this forecast?''
\end{itemize}

\medskip\noindent\textbf{How an answer is built.} The assistant does not answer from general knowledge. For each question it identifies what is asked (a status, a trend, a recommendation, a term or a scientific background question), collects the relevant structured results from the analytical engine, retrieves the relevant passages of the curated library, and asks the language model to compose a short answer from those inputs only, in the language of the user and with the writing rules of \ref{subsec:lay-explanations}. The answer goes through the same checks before it is shown, and it ends with the records and sources it used, one tap away. If the checks fail, the assistant shows the structured data instead and says so. Follow-up questions keep the context of the conversation, so that ``And yesterday?'' refers to the metric just discussed.

\medskip\noindent\textbf{Voice.} The user holds a button, or says a wake phrase in hands-free mode, and speaks the question. Speech to text turns it into a question, and text to speech reads the answer at an adjustable speed. The written answer stays on the screen. The ISS is a noisy place (\ref{sec:hardware-sensors}), so voice use is designed for a headset, which also keeps spoken health information private in a shared module.

\medskip\noindent\textbf{Limits of the assistant.} The assistant answers only with data that the role of the user may see (\ref{subsec:security}), so a crew member cannot ask about another crew member. It never diagnoses, never suggests medication and never contradicts the level or the recommendation of the engine; a question of that kind is answered with an offer to contact the flight surgeon. It is never the path for an emergency: critical alerts reach the crew through the alert system, whether or not anyone asks. Conversations are recorded as text in the audit trail, and recorded audio is not kept after transcription.

\medskip\noindent\textbf{With and without internet.} In the demonstrator the assistant uses the Claude API, which simulates the AI; without a connection, the contingency plan uses the local model on board, with the same inputs and checks (\ref{subsec:offline-llm}).

\subsection{Internationalization}
\label{subsec:languages}

The interface of AHI is available in \textbf{English, Portuguese, Spanish, Japanese and German}. Each user chooses their language, which follows them on every screen, in the explanations, in the health assistant, in the speech and in the exported reports. English is the reference language, so that a crew member and a flight surgeon who use different languages can always switch to a common one during a discussion.

\medskip\noindent\textbf{What is translated.} Everything the user reads or hears: the interface texts, the status words, the metric and dimension names, the recommendation catalog, the template sentences, the generated explanations, the assistant answers, the audio description and the spoken alerts. Stored values, codes and identifiers are never translated, so the data are the same whatever the language.

\medskip\noindent\textbf{Status words in five languages.} The status language is the most critical text of the platform, because it carries the urgency. \ref{tab:status-i18n} shows the proposed translations. They are a working proposal of the team, not yet reviewed, and must be reviewed by native speakers with health background before operational use.

\begin{table}[htbp]
\centering
\vspace{0.8em}
\caption{Status words in the five languages of AHI. Proposal of the team, subject to review by native speakers.}
\label{tab:status-i18n}
\footnotesize
\begin{tabular}{p{2.3cm}p{2.4cm}p{2.3cm}p{2.3cm}p{3.0cm}}
\hline
\textbf{English} & \textbf{Portuguese} & \textbf{Spanish} & \textbf{Japanese} & \textbf{German} \\
\hline
All good & Tudo bem & Todo bien & \ja{良好} & Alles gut \\
Keep an eye & Fique atento & Atenci\'on & \ja{要注意} & Im Blick behalten \\
Act soon & Aja em breve & Actuar pronto & \ja{早めに対応} & Bald handeln \\
Act now & Aja agora & Actuar ahora & \ja{今すぐ対応} & Jetzt handeln \\
Signal lost & Sinal perdido & Se\~nal perdida & \ja{信号なし} & Signal verloren \\
\hline
\end{tabular}
\vspace{1.2em}
\end{table}

\medskip\noindent\textbf{Local conventions, same data.} \ref{tab:locale-formats} summarizes how presentation changes with the language. Numbers follow the decimal mark of each language, units stay in the International System, and times stay in UTC, as in mission operations, so that a time read aloud in any language refers to the same moment.

\begin{table}[htbp]
\centering
\vspace{0.8em}
\caption{Presentation conventions per language. The stored values never change.}
\label{tab:locale-formats}
\footnotesize
\begin{tabular}{p{2.4cm}p{2.6cm}p{3.0cm}p{5.6cm}}
\hline
\textbf{Language} & \textbf{Decimal mark} & \textbf{Example value} & \textbf{Notes} \\
\hline
English & Point & 36.8 \textdegree C & Reference language \\
Portuguese & Comma & 36,8 \textdegree C & Plain language in Brazilian and European usage reviewed \\
Spanish & Comma & 36,8 \textdegree C & Neutral Spanish for international crews \\
Japanese & Point & 36.8 \textdegree C & Bundled font with Japanese characters; denser text \\
German & Comma & 36,8 \textdegree C & Long compound words; layouts tested for length \\
\hline
\end{tabular}
\vspace{1.2em}
\end{table}

\medskip\noindent\textbf{How it is built.}

\begin{itemize}
\item \textbf{No text in the code.} Every visible text lives in a message catalog per language, with rules for plurals and variables, so that adding a language is a translation task and not a code change.
\item \textbf{One medical glossary per language.} Status words, metric names and recommendation texts come from one glossary per language, so that \textit{Act soon} is meant to be as urgent in Japanese as in English. The language model and the speech models use the same glossary (\ref{subsec:guardrails}). The glossary of this report is a proposal of the team; its review by native speakers with health background is a requirement before operational use.
\item \textbf{The language model writes, it does not translate.} Explanations and assistant answers are written directly in the language of the user from the structured result, instead of being written in English and machine translated, and they pass the same checks in every language.
\item \textbf{Layouts that stretch.} German words are long and Japanese text is dense, so components must hold the longest translation without cutting a status word, which can be checked automatically with artificially lengthened text.
\item \textbf{Fonts on board.} Inter covers English, Portuguese, Spanish and German, and the bundled Noto Sans JP covers Japanese (\ref{subsec:typography}).
\item \textbf{Mixed language crews.} Each user sees their own language. Shared artifacts, such as the shift handover report, are produced in English and, on request, in the language of the reader.
\item \textbf{Fallback.} A text missing in one language falls back to English, and the missing key is logged for the translators.
\item \textbf{Offline.} All catalogs are bundled with the application, and no translation service is called while the platform runs.
\end{itemize}

\subsection{Audio Description and Voice Operation}
\label{subsec:accessibility}

Accessibility is part of the design, with conformance to WCAG 2.1 level AA as the design target \cite{w3c_wcag21}. The features below follow the success criteria of that standard; conformance has not been verified in this project and would require a formal audit and tests with users before operational use. The foundation is the same for every user: global keyboard shortcuts, keyboard navigation of timeline markers, a visible focus ring, a link that skips to what needs attention as the first focusable element, accessible labels on every control, live regions that announce new alerts to screen readers, and a status language that never relies on color alone, since every status has its own word, icon and shape (\ref{tab:status-language}).

On top of this foundation, AHI offers \textbf{audio description}: the platform describes aloud what is on the screen and lets the user operate it by voice. It is generated from the same state as the screens, as a structured description of each screen (its headline, its regions and its items), so that the spoken version and the visual version always show the same thing.

\medskip\noindent\textbf{Describing what is visual.}

\begin{itemize}
\item \textbf{The 3D body.} The body is described as a list of regions with their status, the most urgent first, for example: ``Heart: act now, heart rate 142, above your normal of 58 to 72. Lungs: keep an eye. All other regions: all good.'' The user can move from region to region with the keyboard or by voice, and selecting a region reads its detail.
\item \textbf{Charts.} A chart is described by its trend, its position against the personal normal, the gaps of missing data and the forecast, for example: ``Heart rate variability, last seven days: going down slowly, below your normal since Tuesday, two hours of missing data last night; forecast: likely to stay low tomorrow.''
\item \textbf{Alerts.} New alerts are announced as soon as they appear, in order of urgency. A critical alert interrupts any other speech; lower levels wait for a pause.
\item \textbf{Recommendations.} Each recommendation is read with its action, who can act and when it will be reassessed.
\end{itemize}

\medskip\noindent\textbf{Voice operation.} The main actions of the crew view can be done by voice, as listed in \ref{tab:voice-commands}. Commands exist in the five languages, and every command has a keyboard equivalent.

\begin{table}[htbp]
\centering
\vspace{0.8em}
\caption{Main voice commands of AHI (English form).}
\label{tab:voice-commands}
\footnotesize
\begin{tabular}{p{3.8cm}p{10.2cm}}
\hline
\textbf{Command} & \textbf{Effect} \\
\hline
``Status'' & Reads the status word of the user and the one sentence summary \\
``Read alerts'' & Reads the open alerts, the most urgent first \\
``Explain'' plus a metric & Reads the plain explanation of that metric (\ref{subsec:lay-explanations}) \\
``What should I do?'' & Reads the recommended actions that can be taken on board \\
``Done'' & Logs the current recommended action as done, after confirmation \\
``Next'', ``Back'', ``Repeat'' & Moves through regions, items or the last message \\
``Slower'', ``Faster'' & Changes the speed of speech \\
``Ask'' plus a question & Sends the question to the health assistant (\ref{subsec:assistant}) \\
``Stop'' & Stops speaking immediately \\
\hline
\end{tabular}
\vspace{1.2em}
\end{table}

\medskip\noindent\textbf{For each group of users.}

\begin{itemize}
\item \textbf{Blind users} operate the whole crew view by keyboard, screen reader or voice, with every screen available in its spoken form.
\item \textbf{Users with low vision} use large text, high contrast and zoom without loss of content, with the spoken description available at any time.
\item \textbf{Users with dyslexia} can have any text read aloud with the current sentence highlighted, increase letter, word and line spacing, and rely on the short sentence rule.
\item \textbf{Astronauts busy with a task,} during exercise, maintenance or an experiment, use hands-free mode: AHI reads short status summaries at chosen moments, announces new alerts and answers spoken questions.
\item \textbf{Users with reduced hearing}, including crew members whose hearing assessment shows a threshold shift (\ref{sec:analytical-methodology}), can raise the volume of speech or rely on the written version, since every spoken message is also shown.
\end{itemize}

\medskip\noindent\textbf{Sounds per status level.} Each status level has its own short sound, distinct in pitch and rhythm, so that the level can be recognized without looking. Sounds never replace words: every sound is followed by the status word and has a written equivalent.

\medskip\noindent\textbf{Privacy and environment.} Audio is designed for a headset, because the cabin is noisy and because spoken health information should not be heard by the whole module. Reading private data aloud on a speaker requires an explicit choice of the user. Volume, speed, voice and level of detail are personal settings.

\subsection{Guardrails and Validation Required}
\label{subsec:inclusion-evaluation}

The guardrails of \ref{subsec:guardrails} apply to every text of this section, in every language and in every form, written or spoken. The design choices of this section rest on the published guidance cited above; they have not been validated with users in this project. Before operational use, the following evaluations would be required, with the metrics of \ref{sec:results}:

\begin{itemize}
\item \textbf{Explanation fidelity and readability:} the share of explanations whose numbers, status word and uncertainty word match their input, and their reading level, per language.
\item \textbf{Comprehension:} whether lay readers understand the status, the trend and the recommended action from the plain explanation alone, measured in a usability study.
\item \textbf{Translation quality:} review of the glossary and of a sample of explanations by native speakers with health background.
\item \textbf{Speech accuracy:} word error rate of spoken questions and commands in a recorded noisy cabin environment, per language.
\item \textbf{Accessibility:} a formal WCAG audit and task tests with blind, low vision and dyslexic users and with users in hands-free mode.
\end{itemize}


\section{Implementation and Demonstration}
\label{sec:implementation-demo}

This section describes how AHI is built in stages and how it is demonstrated. Real crew data from ISS equipment are not available to the project, so the demonstration replays simulated data through the same device adapters, analysis and screens that real data would use. Every simulated value is labeled \texttt{SIMULATED}, and the platform is decision support and not a medical device.

\subsection{Implementation Plan}
\label{subsec:features}

The platform is built in stages that follow the device priority of \ref{sec:hardware-sensors}. A complete, consistent integration of a few devices with real prediction is worth more than a shallow integration of all of them. \ref{tab:features} lists the stages.

\begin{table}[htbp]
\centering
\vspace{0.8em}
\caption{Implementation stages of AHI.}
\label{tab:features}
\footnotesize
\begin{tabular}{p{2.6cm}p{5.6cm}p{5.8cm}}
\hline
\textbf{Stage} & \textbf{Devices and data} & \textbf{Capabilities} \\
\hline
1. Core & Bio-Monitor, Actiwatch Spectrum, CEVIS, ARED, T2 & Canonical model, adapters, personal baselines, anomaly detection, data quality, crew and ground views \\
2. Prediction & Same, plus blood pressure and ECG & Trend forecasts, combined rules, transparent risk score, recommendations with autonomy, scenarios with and without intervention \\
3. Environment & CO\textsubscript{2} monitors, AQM, radiation monitors and dosimeters; replay of real cabin data & Environmental context and cabin rules, accumulated exposure \\
4. Periodic data & Sample results, urine analyzer, ultrasound findings, Holter, Dry-EEG Headband & Immune and metabolic dimensions, examination results on the timeline \\
5. Learned models & Mission and authorized historical data & Isolation Forest and XGBoost behind the same interface, with model cards and shadow mode \\
6. Explanation and review & Curated source library & Language model explanation layer with guardrails, flight surgeon review workflow, persistent audit \\
7. Radiation and space weather & Personal dosimeters, area monitors, DONKI through the NASA Open APIs & Career dose against the NASA-STD-3001 limit, projection per mission profile, radiation shelter rule \\
8. Vision and hearing & Vision screening, optical coherence tomography, fundoscope, Ultrasound 2, visual symptoms in the check-in, On-Orbit Hearing Assessment, Acoustic Monitor, sound level meter & Vision and hearing as two new dimensions: eye and hearing results against the preflight examinations, noise dose, co-occurrence view \\
9. Global access & Accounts and roles; message catalogs; speech models & Role-based access and authentication, five languages, health assistant by text and voice, audio description \\
\hline
\end{tabular}
\vspace{1.2em}
\end{table}

\medskip\noindent\textbf{Language features in the demonstrator.} The demonstrator is a web application, planned to be deployed on Vercel, so it runs with internet access. Its explanations and health assistant use the Claude API to simulate the AI language features (\ref{subsec:offline-llm}). The internet is not essential to the design: a real system on board would follow the contingency plan of \ref{subsec:offline-llm} without a connection. The interface labels these texts as generated through the API, and when the API is not reachable the demonstrator falls back to template text, which is the same behavior that the contingency plan has without any model.

\subsection{Scenarios}
\label{subsec:scenarios}

A seeded, reproducible simulator produces a crew of four with fictional callsigns and the cabin, through simulated counterparts of the device adapters. Its noise is correlated in time, not white, so that signals behave like physiological series. Scripted scenarios exercise specific behaviors of the platform, as listed in \ref{tab:scenarios}. Their magnitudes and durations are demonstration parameters and do not describe any real physiology.

\begin{table}[htbp]
\centering
\vspace{0.8em}
\caption{Scripted scenarios used for demonstration and testing. All simulated.}
\label{tab:scenarios}
\small
\begin{tabular}{p{3.4cm}p{6.4cm}p{4.2cm}}
\hline
\textbf{Scenario} & \textbf{What happens} & \textbf{Behavior exercised} \\
\hline
Nominal day & Everyone stays within their personal ranges & No false alerts \\
Acute cardio event & Heart rate rises at rest while oxygen saturation falls & Fast critical alert and escalation \\
Sensor failure & The Bio-Monitor shirt becomes degraded, stale and offline & Missing data never shown as normal \\
Cabin CO\textsubscript{2} rise & Carbon dioxide rises for hours; breathing and reaction time shift for everyone & Shared environmental cause, co-occurrence \\
Progressive fatigue & Sleep worsens, heart rate variability falls and resting heart rate rises over days & Slow multi signal pattern \\
Exercise performance decline & Exercise output falls day after day & Slow single metric decline \\
\hline
\end{tabular}
\vspace{1.2em}
\end{table}

Two scenarios are planned for the new features of the final scope. In a \textit{solar particle event}, a DONKI notification arrives through the ground node and, a while later, the on board dose rate rises; the platform recommends preparing the radiation shelter from the first trigger, and the scenario is repeated with the link to Earth down to show that the on board trigger works alone. In a \textit{hearing threshold shift}, a period of high noise during sleep is followed by an audiogram with a low frequency shift; the platform requests a repeat assessment and shows the exposure next to the result without claiming a cause.

\subsection{Presentation Script}
\label{subsec:presentation-script}

The demonstration follows one crew member through a mission scenario in five steps. The figures of this subsection are screenshots of the AHI demonstrator running the scripted scenarios with simulated data.

\begin{enumerate}
\item \textbf{Data collection.} Simulated Bio-Monitor, Actiwatch Spectrum and exercise equipment feed the platform through their adapters. Each value carries its device, time, context and origin.
\item \textbf{Anomaly identification.} Sleep, heart rate variability and exercise performance change at the same time. No single value crosses a line, but the combined pattern against the person's own normal is flagged (\ref{fig:timeline-fatigue}).
\item \textbf{Health prediction.} The predictive screen shows the trend for 24, 48 and 72 hours, with its uncertainty, and the scenario without intervention (\ref{fig:body-detail}).
\item \textbf{Recommendation.} The platform suggests recovery, rest, hydration and medical follow-up actions, each with its evidence, its autonomy and its time for reassessment, and the language model explains the pattern in plain words.
\item \textbf{Outcome evaluation.} After the intervention is logged, the indicators are followed, and the screen compares their evolution with the projection without intervention (\ref{fig:whatif}).
\end{enumerate}

\hwfigure{0.96}{images/ahi-timeline-fatigue.png}{Progressive fatigue: LYRA moves to act soon because shorter sleep, lower heart rate variability and a higher resting heart rate show up together; the status timeline shows the pattern building over time.}{fig:timeline-fatigue}{AHI demonstrator, screenshot by the authors with simulated data, published in the project documentation portal (\texttt{docs-site/public/screenshots/timeline-fatigue.png}).}

\hwfigure{0.96}{images/ahi-body-detail.png}{Body system detail: the heart of LYRA is selected on the atlas, the rest of the body is ghosted, and the panel shows each metric against the personal normal with its history and forecast.}{fig:body-detail}{AHI demonstrator, screenshot by the authors with simulated data, published in the project documentation portal (\texttt{docs-site/public/screenshots/body-detail.png}).}

\hwfigure{0.96}{images/ahi-whatif.png}{What-if comparison for sleep extension: heart rate variability projected with and without the action, against the personal normal, labeled as a simulation for discussion and not evidence of clinical efficacy.}{fig:whatif}{AHI demonstrator, screenshot by the authors with simulated data, published in the project documentation portal (\texttt{docs-site/public/screenshots/whatif.png}).}

The demonstration also shows the acute cardiac event, to show fast alerting and escalation (\ref{fig:critical-alert}), the sensor failure, to show that missing data is never shown as normal, and the cabin CO\textsubscript{2} rise, to show co-occurrence without a claim of cause. Two mission tools complete it: the Earth link, which shows the communication delay and why on board steps come first (\ref{fig:mars-comms}), and the shift handover report with its FHIR export (\ref{fig:handover}).

\hwfigure{0.96}{images/ahi-critical-alert.png}{Acute cardiac event: critical banner, heart and lungs highlighted on the atlas, risk score 100 and recommended actions with their autonomy.}{fig:critical-alert}{AHI demonstrator, screenshot by the authors with simulated data, published in the project documentation portal (\texttt{docs-site/public/screenshots/critical-alert.png}).}

\hwfigure{0.96}{images/ahi-mars-comms.png}{Earth link under the farthest Mars profile: about 22 minutes one way and 44 minutes for a round trip, so a question to Earth waits at least 44 minutes and on board steps come first.}{fig:mars-comms}{AHI demonstrator, screenshot by the authors with simulated data, published in the project documentation portal (\texttt{docs-site/public/screenshots/mars-comms.png}).}

\hwfigure{0.96}{images/ahi-handover-report.png}{Shift handover report: crew at a glance, status, open alerts and pending actions, with print and FHIR R4 export.}{fig:handover}{AHI demonstrator, screenshot by the authors with simulated data, published in the project documentation portal (\texttt{docs-site/public/screenshots/handover-report.png}).}

A note for presenters: the same sentence must accompany every claim. The data in the demonstration are simulated, and the tool supports decisions and does not make them.


\section{Results}
\label{sec:results}

This section defines how AHI is evaluated and what results are expected. \textbf{It reports no clinical performance.} No real crew data were available to the project, so no sensitivity, specificity or accuracy on astronauts can be claimed. The evaluation is designed in three stages of increasing realism, and each stage answers a different question.

\subsection{Metrics}
\label{subsec:metrics}

\begin{table}[htbp]
\centering
\vspace{0.8em}
\caption{Evaluation metrics of AHI.}
\label{tab:checks}
\small
\begin{tabular}{p{4.0cm}p{6.0cm}p{4.0cm}}
\hline
\textbf{Metric} & \textbf{Definition} & \textbf{Applies to} \\
\hline
False alert rate & Alerts at moderate or above per crew day when nothing is wrong & Level 1, Level 2 \\
Detection time & Time from the start of a deterioration to the first alert & Level 1, Level 2 \\
Forecast error & Mean absolute error of projections at 24, 48 and 72 hours & Level 2 \\
Interval coverage & Share of observed values inside the stated forecast interval & Level 2 \\
Discrimination & Area under the ROC curve of the risk models, where labeled events exist & Level 2 \\
Calibration & Agreement between estimated and observed event frequency & Level 2 \\
Missing data handling & Share of missing or stale signals reported as no data & All levels \\
Recommendation review & Share of recommendations accepted, modified or rejected by the flight surgeon & Level 3 \\
Explanation fidelity & Share of generated texts whose numbers and citations match their inputs & Explanation layer \\
Usability & Task success and time for crew and flight surgeons & Interface \\
Assistant answer quality & Share of answers judged correct, grounded and in plain language by reviewers, per language & Health assistant \\
Speech recognition accuracy & Word error rate of spoken questions, in a recorded noisy cabin environment, per language & Voice and audio description \\
Access control correctness & Share of test requests where a role sees or does only what it is allowed to & Security \\
Shelter warning time & Time from the first trigger to the shelter recommendation, with and without the link to Earth & Radiation rule \\
\hline
\end{tabular}
\vspace{1.2em}
\end{table}

\subsection{Evaluation}
\label{subsec:evaluation}

\medskip\noindent\textbf{Stage 1: simulated scenarios.} A headless harness replays the engine on the scripted scenarios of \ref{sec:implementation-demo} over many seeds. The nominal scenario measures the false alert rate, since any alert there is false. The other scenarios measure detection time and check the intended behavior: the acute event reaches the critical level within minutes, slow patterns reach the high level and stop there, and a sensor failure produces no data and never a normal status. This stage shows that the rules and models behave as designed. It is optimistic by construction, because the same simulator generates and evaluates the data.

\medskip\noindent\textbf{Stage 2: retrospective replay of real data.} Real public data are replayed through the engine: ISS cabin periods from the Environmental Data Application \cite{osdr_eda}, laboratory panels from OSDR \cite{osdr_osd569,osdr_osd575} and, where they can be obtained, sleep and cognition data from analog studies \cite{nasa_nlsp,nasa_lsda}. This stage tests the engine on data it was not tuned on, for false alert rate, forecast error and interval coverage. Analog results are reported as analog results.

\medskip\noindent\textbf{Stage 3: operational shadow mode.} With authorization, AHI runs next to the existing medical operations on real crew data, without affecting decisions. Flight surgeons review its alerts and recommendations, which gives labeled cases for discrimination, calibration and recommendation review. This stage depends on access that the project does not have.

\begin{table}[htbp]
\centering
\vspace{0.8em}
\caption{Expected results, stated as hypotheses to be tested.}
\label{tab:detection}
\footnotesize
\begin{tabular}{p{4.4cm}p{9.6cm}}
\hline
\textbf{Hypothesis} & \textbf{How it is tested} \\
\hline
A personal baseline gives fewer false alerts than fixed population limits for the same detection time & Stage 1 and 2: compare both methods on the same data \\
Combined rules detect slow multi signal deterioration earlier than any single metric & Stage 1: progressive fatigue scenario; Stage 3: reviewed cases \\
Missing data is never reported as normal & Stage 1: sensor failure scenario and automated tests \\
Forecast intervals cover close to their stated share of values & Stage 2: replay of cabin and daily metrics \\
Learned models improve on rules only when enough real data exist & Stage 2 and 3: compare with the rules on held out data \\
Generated explanations contain no number or limit absent from their input & Explanation fidelity on reference cases \\
The local model of the contingency plan reaches the explanation fidelity of the Claude API on the reference cases & The same reference cases, in the five languages, run on both \\
Blind users and users busy with a task can complete the main tasks by audio and voice alone & Usability study with audio description and hands-free mode \\
The on board dose rate trigger recommends the shelter without any link to Earth & Stage 1: solar particle event scenario with the link down \\
\hline
\end{tabular}
\vspace{1.2em}
\end{table}

\medskip\noindent\textbf{Threats to validity.} \ref{tab:threats} lists the main reasons why results of the earlier stages should not be generalized.

\begin{table}[htbp]
\centering
\vspace{0.8em}
\caption{Threats to the validity of the evaluation.}
\label{tab:threats}
\small
\begin{tabular}{p{5.0cm}p{9.0cm}}
\hline
\textbf{Threat} & \textbf{Effect} \\
\hline
The same simulator generates and evaluates & Detection in Stage 1 is optimistic and shows internal consistency only \\
Few scenarios & Risk of overfitting rules to a few stories; more scenarios and seeds are needed \\
Public data are small and short & Inspiration4 panels cover four people on a three day flight with no in-flight samples \\
Analog populations differ & Analog sleep and cognition data are not astronaut data \\
Few labeled outcomes & Clinical sensitivity and specificity cannot be estimated before Stage 3 \\
\hline
\end{tabular}
\vspace{1.2em}
\end{table}

\subsection{Evidence}
\label{subsec:evidence}

The evidence produced by each stage is stored with the version of the engine and of every model, so that any change of a rule or model appears as a difference in the results. Each model card records its evaluation, and the evaluation can be reproduced with one command. \ref{tab:missing-evidence} lists the evidence that only later stages can provide.

\begin{table}[htbp]
\centering
\vspace{0.8em}
\caption{Evidence that requires real data or expert review.}
\label{tab:missing-evidence}
\small
\begin{tabular}{p{7.0cm}p{7.0cm}}
\hline
\textbf{Evidence} & \textbf{Requires} \\
\hline
Accuracy on real astronaut data & Authorized access to crew data (Stage 3) \\
Sensitivity and specificity against a clinical reference & Labeled outcomes reviewed by flight surgeons \\
Calibration of risk models & Enough labeled events \\
Review of thresholds and recommendations & Flight surgeons \\
Usability with crew and flight surgeons & A study with representative users \\
Accessibility & A formal audit, screen reader tests and tests with blind, low vision and dyslexic users \\
Quality of translations & Review by native speakers with health background in each language \\
\hline
\end{tabular}
\vspace{1.2em}
\end{table}


\section{Limitations and Ethics}
\label{sec:limitations-ethics}

AHI is a monitoring, prediction and decision support platform. It is \textbf{not a medical device}, is not certified and has not been validated on real crews. It does not diagnose, does not recommend medication or treatment, and does not replace the flight surgeon, mission control or established medical protocols. These limits are part of the product and are stated on every screen. This section collects the limitations of the approach and the ethical rules that govern it.

\subsection{Scientific Limitations}
\label{subsec:scientific-limitations}

\begin{itemize}
\item \textbf{The hardware premise is an assumption.} The report assumes that data come from ISS equipment. The project has no access to these devices or to their data, and their exact output formats were not verified. Some devices, such as the Dry-EEG Headband, are research hardware used during specific investigations and not by every crew member.
\item \textbf{Astronaut data are scarce.} Crews are small, public continuous telemetry does not exist, and labeled outcomes are rare. Learned models can be trained and validated only with authorized data, and models trained on simulated data are not used for crew decisions.
\item \textbf{Thresholds and weights are configuration.} No fixed number in AHI is an official NASA limit. Every parameter must be grounded in a cited standard or reviewed by a flight surgeon before operational use.
\item \textbf{Predictions differ in precision.} Sleep and fatigue have frequent data; bone, immune and metabolic changes depend on periodic samples and examinations, so their predictions are slow and uncertain.
\item \textbf{The risk score is not a probability} until a learned model is calibrated on real events.
\item \textbf{Activity context may be wrong.} If the context is taken from the schedule, a deviation from the schedule compares a value with the wrong baseline.
\item \textbf{Intervention effects are uncertain.} The comparison with and without intervention is only as good as the evidence for its effect sizes.
\item \textbf{Self report is not a clinical instrument.} AHI makes no psychiatric or psychological assessment and uses self report only as an indicator of risk.
\item \textbf{Correlation is not causation.} Signals that move together are reported as co-occurrence and not as cause. This applies in particular to noise exposure and hearing thresholds, since low frequency shifts on the ISS are not typically caused by noise \cite{coste2024hearing}.
\item \textbf{The career dose is only as good as its inputs.} It depends on the dose of previous missions in the medical record, on dosimeters that report effective dose and on an estimate between readouts. AHI follows the published limit and does not estimate radiation health risk.
\item \textbf{Public data have limits.} The Inspiration4 panels describe four people on a three day flight with no in-flight samples, and the cabin data of the Environmental Data Application describe the cabin during payload missions and not a crew.
\end{itemize}

\subsection{Technical Limitations}
\label{subsec:technical-limitations}

\begin{itemize}
\item \textbf{Device integration depends on access.} No open interface of ISS equipment is assumed. Each adapter depends on documentation and permission that the project does not have.
\item \textbf{On board computing is limited.} Learned models may need to be trained on the ground, and the local language and speech models must be small enough for the on board computers, so their explanations and answers may be simpler than those that the Claude API produces in the demonstrator.
\item \textbf{The contingency for language features is not demonstrated.} The demonstrator simulates the AI with the Claude API, which needs the internet. Without a connection, the contingency plan relies on a local model, whose quality has still to be shown (\ref{sec:results}).
\item \textbf{Space weather data come from an external service.} DONKI is reached through the NASA Open APIs by the ground node only, depends on an API key and on the availability of the service, and is not the official forecast source for operations. The on board dose rate trigger covers its absence.
\item \textbf{Speech and translation can fail.} Speech recognition is harder in a noisy cabin and with accents, and a translation can shift the meaning of a medical term. Every spoken answer is also written, and status words come from one glossary per language, but errors remain possible.
\item \textbf{Inclusion features are not validated.} The writing rules, the translations and the accessibility features are based on published standards and studies \cite{w3c_wcag21,plain_language_guide,budescu2009uncertainty}. The project has not reviewed the translations with native speakers, has not audited WCAG conformance and has not tested the features with blind, low vision or dyslexic users; these are requirements before operational use.
\item \textbf{The communication link is imperfect.} Delay, loss of signal and limited bandwidth affect synchronization, and the store and forward queue reduces but does not remove their effect.
\item \textbf{The language model can be wrong.} Guardrails check numbers, citations and wording, but they cannot guarantee a correct explanation, which is why the structured result is always available.
\item \textbf{Hardware changes.} Devices are replaced over time, and each change needs a new adapter and a new validation.
\end{itemize}

\subsection{Privacy}
\label{subsec:privacy}

Crew health data are personal and sensitive, and AHI is designed around that. Access is role based, data are encrypted in transit and at rest, identifiers are pseudonymized outside the medical roles, and the audit trail is persistent and tamper evident (\ref{sec:software-architecture}). Researchers receive de-identified extracts. The FHIR export carries pseudonymized identifiers only. The demonstrator uses simulated data only, so sending its texts to the Claude API exposes no crew data. With real crew data, a language model reached through the internet would receive data outside the mission systems: even with structured, pseudonymized inputs, such use would require an agreement with the crew and the medical authority, and the local model of the contingency plan avoids it (\ref{sec:artificial-intelligence}). Access is granted by role and every action is recorded with its author. Voice features add a risk of their own: spoken health information can be heard by others in a shared module, so audio is designed for a headset, private data are not read aloud on a speaker without an explicit choice, and recorded audio is not kept after transcription.

Populations are never mixed. Data from long duration astronauts, commercial astronauts on a short flight, analog participants and non human studies are different populations. Analog results are never presented as astronaut results, and fields that a source does not state are recorded as not stated and never imputed.

\subsection{Risks}
\label{subsec:risks}

The main risks of the use or misuse of a tool like AHI, and their mitigations, are listed in \ref{tab:risks}. The mitigations reduce the risks but do not remove them. The table was prepared by the authors and not through a formal hazard analysis.

\begin{table}[htbp]
\centering
\vspace{0.8em}
\caption{Principal risks and mitigations.}
\label{tab:risks}
\footnotesize
\begin{tabular}{p{4.0cm}p{5.6cm}p{4.4cm}}
\hline
\textbf{Risk} & \textbf{Mitigation} & \textbf{Residual limit} \\
\hline
Missing data read as normal & A no data status with a different shape and color; confidence falls; tests assert it & Applies only to the monitored signals \\
The score read as a probability & The interface describes it as an operational level until a model is calibrated & A number invites over-reading \\
False reassurance when nothing is flagged & No status means within the personal normal only & Slow adaptation can move the baseline \\
Over-reliance on automation & Human review required for high and critical items; no autonomous decision & Alert fatigue must be measured \\
Configuration values taken as official limits & Every fixed number is labeled with its source or as configuration & Numbers can be quoted out of context \\
Correlation read as cause & Combined rules report co-occurrence & Users may still infer a cause \\
Diagnostic or prescriptive wording & Plain status words; a catalog without medication; guardrails on generated text & Wording needs clinical review \\
Generated text hiding uncertainty & Grounding checks, labels on generated text, structured result always available & No check is complete \\
Simulated data mistaken for real & A visible simulated badge and the origin on every record & Depends on the reader \\
Assistant used as a doctor & Plain refusal of diagnosis and medication, offer to contact the flight surgeon, role scope & Users may still over-trust a fluent answer \\
Translation changes the urgency of a status & One glossary per language; English always available & Glossary not yet reviewed by native speakers \\
Shelter recommendation missed or late & Two independent triggers, one fully on board; tests with the link down & Depends on the radiation monitors \\
Access beyond the role & Role checks in every screen, API and the assistant; recorded emergency access & Misconfigured roles \\
\hline
\end{tabular}
\vspace{1.2em}
\end{table}

The ethical rules of the project are the following. Decision support is for people, and every high or critical item requires human review. Plain language comes first and no diagnostic wording is used. The platform must be usable by every crew member, whatever their language or disability. Missing data is never shown as normal. Origin, quality and source are recorded for every record, and models, parameters and rules are versioned, so that every alert can be traced to its data and its rule or model. Limitations and uncertainty are shown to the user. No operational claim is made before validation on real, appropriately consented data.


\section{Future Work}
\label{sec:future-work}

The items below extend AHI beyond the scope described in this report. They are plans in rough order of priority and not commitments.

\subsection{Validation}
\label{subsec:future-validation}

\begin{enumerate}
\item Retrospective replay of real ISS cabin periods and of public laboratory and analog data, as the first test of the engine outside simulated data (\ref{sec:results}).
\item Operational shadow mode with authorized crew data, with alerts and recommendations reviewed by flight surgeons, to obtain labeled cases.
\item A citable source for every threshold, with document, section and version, and review of thresholds, wording and the recommendation catalog by flight surgeons.
\item A usability study with crew members and flight surgeons, with tasks taken from the user journeys, in the five languages, and a formal accessibility audit with blind, low vision and dyslexic users.
\item Training and evaluation of the local language model and of the speech models on the on board hardware, compared with the Claude API used in the demonstrator on the same reference cases.
\item Prospective clinical validation, which is outside the scope of a project of this kind and would need a medical sponsor.
\end{enumerate}

\subsection{Data and Devices}
\label{subsec:future-real-data}

\begin{enumerate}
\item A table that relates each metric to the real device that measures it, to the NASA Standard Measure and to the HRP hazard, confirmed against the interface documentation of each device.
\item Further ISS hardware described in public sources and not part of the current premise, such as the T-Mini core temperature headband and the Wearable Monitoring vest \cite{nasa_wearable_tech_iss}, and the new teal CEVIS ergometer \cite{nasa_jsc_exercise_countermeasures}.
\item Access to historical data through the Life Sciences Data Archive and the NASA Life Sciences Portal \cite{nasa_lsda,nasa_nlsp}, with formal requests where required.
\item Real radiation data from RadLab once its coverage, units and format can be checked \cite{osdr_radlab}.
\end{enumerate}

\subsection{Missions Beyond Low Earth Orbit}
\label{subsec:future-platform}

The mission scenarios of \ref{subsec:mission-scenarios} show that the risks grow with the duration of the mission and its distance from Earth. The steps below extend AHI from the ISS to those scenarios.

\begin{enumerate}
\item Integration of the dynamic phases: the acceleration records required by NASA-STD-3001 for ascent, entry and landing \cite{nasa_ochmo_tb024_occupant}, imported after flight and placed on the timeline next to the check-ins after launch and after landing.
\item Adaptation to Artemis and lunar missions, whose Standard Measures differ from those of the ISS \cite{nasa_artemis_ii_standard_measures}, with smaller devices and less on board computing.
\item Longer autonomy for Mars missions, with communication delays of up to about 24 minutes one way \cite{nasa_space_communications_seven_things}: a larger on board model set, a flight approved version of the local language and speech models, and on board steps for every high item.
\item Space weather beyond low Earth orbit, where the magnetic field of the Earth no longer protects the crew: integration of the space weather support of the mission alongside DONKI, and shelter planning that accounts for extravehicular activity on the lunar surface.
\item Bandwidth aware synchronization, with priorities so that alerts and summaries reach Earth before raw data.
\end{enumerate}

\subsection{Analytics}
\label{subsec:future-analytics}

\begin{enumerate}
\item Activity recognition from sensors instead of the schedule, which would remove the assumption of a correct activity context.
\item Sequence models for multivariate forecasting once long histories exist.
\item Learning of intervention effects from the follow-up of logged interventions, within protocols approved by the medical team.
\item Integration of bone examination results, so that bone loss can be followed as point measurements on the timeline in the same way as the eye and hearing examinations \cite{lang2004bone}.
\item More languages for the interface and the assistant, chosen with the partner agencies of each mission, and sign language video for deaf users on the ground.
\end{enumerate}


\section{Conclusion}
\label{sec:conclusion}

\subsection{Summary}
\label{subsec:summary}

This report presented Astronaut Health Intelligence (AHI), a platform for the prevention of health problems and the predictive protection of crew health on space missions. AHI integrates physiological, behavioral and environmental data from equipment in use on the International Space Station: the Bio-Monitor, the Actiwatch Spectrum and the Dry-EEG Headband among the wearables; the blood pressure and ECG monitor, the Holter monitor, the pulse oximeter, the Ultrasound 2, the urine analyzer and the sample collection kits among the medical equipment; the optical coherence tomography and the On-Orbit Hearing Assessment for the eyes and ears; CEVIS, ARED and T2 for exercise; and the air quality, carbon dioxide, combustion, noise, radiation, water and microbial monitors for the cabin, complemented by the space weather notifications of NASA DONKI.

The platform organizes health into ten dimensions, with vision and hearing as two separate dimensions, and uses the five hazards of human spaceflight as the axis of its risk analysis. It works on three levels. Real-time monitoring compares each person with their own normal and detects acute changes, slow drifts and missing data. Predictive analytics projects trends for 24 to 72 hours and estimates risks with transparent rules first and learned models as real data allow. Prescriptive intelligence recommends exercise, nutrition, sleep, well-being and medical follow-up actions, each with its evidence, its autonomy and its follow-up. It follows the career radiation dose of each astronaut against the NASA-STD-3001 limit and recommends preparing the radiation shelter when a solar particle event is reported or measured. A language model explains the results in plain language and powers a health assistant by text and voice, under strict guardrails, and never computes them. Internet access is useful but not essential: the demonstrator simulates the AI with the Claude API, and an on board backend with a contingency plan, including local language and speech models, keeps the platform working without a connection. Role-based access, five languages and audio description make the platform usable by an international and diverse crew.

\subsection{Contributions}
\label{subsec:contributions}

\begin{enumerate}
\item \textbf{A hardware grounded design,} in which every data source is a device that public sources describe as used on the ISS, with what each source does and does not establish.
\item \textbf{A canonical health observation model} with origin, quality, source and context on every record, and one adapter per device, so that the analysis does not depend on where a value comes from.
\item \textbf{An explainable analysis chain} from personal baselines to risk, forecasts and recommendations, in which every output can be traced to its data and to a versioned rule or model.
\item \textbf{A treatment of missing data as a first class state,} reported as such and never as normal.
\item \textbf{An operational view of distance,} in which recommendations carry an autonomy tag and a long Earth link places on board steps first.
\item \textbf{A contingency design for language features,} in which the explanations and the assistant are built behind one interface, served by the Claude API in the demonstrator and by a local on board model when there is no internet access.
\item \textbf{Inclusion as a requirement,} with five languages, a plain language assistant by text and voice, and audio description for blind and low vision users, people with dyslexia and astronauts busy with a task.
\item \textbf{A staged evaluation plan,} from simulated scenarios to the replay of real public data and an operational shadow mode, with the evidence each stage can and cannot provide.
\end{enumerate}

The limits are stated as clearly as the design. AHI does not diagnose, does not prescribe, and does not replace the flight surgeon. Its hardware premise is an assumption about access that the project does not have, and its clinical performance can only be established with authorized crew data and expert review. Within those limits, AHI shows how heterogeneous health signals from real space station equipment can be turned into early, explainable and honest support for the decisions of crews and their medical teams.


% Start the bibliography on a new page before the template places its hyperlink anchor.
% Without this the table of contents link for References points to the last page of the Conclusion.
\clearpage
\printbibliography

\begin{appendices}
\section{Articles and NASA Sources}
\label{app:articles-nasa-sources}

\section{Algorithms}
\label{app:algorithms}

\section{Diagrams}
\label{app:diagrams}

\section{Tests}
\label{app:tests}

\end{appendices}

\end{document}
% End of Document %%%%%%%%%%%%%%%%%%%%%%%%%%%%%%%%%%%%%%%%%%%%%%%%%%%%%%
